\part{Hodge Theory II}\label{hodge-theory-ii}\thispagestyle{fancy}\par{}P. Deligne.
  "Théorie de Hodge II".
  \emph{Pub. Math. de l'IHÉS} \textbf{40} (1971) pp. 5–58.
  publications.ias.edu/node/361\chapter{Introduction}\thispagestyle{fancy}\label{hodge-theory-ii-introduction}\par{}\emph{Work presented as a doctoral thesis at l'Université d'Orsay.}\par{}By Hodge, the cohomology space \(\operatorname {H}^n(X,\mathbb {C})\) of a compact Kähler variety \(X\) is endowed with a "Hodge structure" of weight \(n\), i.e. a natural bigrading
  \[     \operatorname {H}^n(X,\mathbb {C})     = \bigoplus _{p+q=n} \operatorname {H}^{p,q}   \]
  which satisfies \(\overline {\operatorname {H}^{p,q}}=\operatorname {H}^{q,p}\).
  We will show here that the complex cohomology of a non-singular, not necessarily compact, algebraic variety is endowed with a structure of a slightly more general type, which presents \(\operatorname {H}^n(X,\mathbb {C})\) as a "successive extension" of Hodge structures of decreasing weights, contained between \(2n\) and \(n\), whose Hodge numbers \(h^{p,q}=\dim \operatorname {H}^{p,q}\), are zero for both \(p>n\) and \(q>n\).\par{}The reader will find an explanation in [\hyperref[hodge-theory-i]{Hodge Theory I}] of the yoga that underlies this construction.\par{}The proof, which is essentially algebraic, relies on one hand on Hodge theory, and on the other on Hironaka's resolution of singularities, which allows us, via a spectral sequence, to "express" the cohomology of a non-singular quasi-projective algebraic variety in terms of the cohomology of non-singular projective varieties.\par{}\cref{hodge-theory-ii-1} contains, apart from reminders on filtrations gathered together for the ease of the reader, two key results:

  
 \begin{enumerate}
    
 \item[(a)]
      \hyperref[hodge-theory-ii-1.2.10]{Theorem~\ref*{hodge-theory-ii-1.2.10}}, which will only be used via its corollary, \hyperref[hodge-theory-ii-2.3.5]{Theorem~\ref*{hodge-theory-ii-2.3.5}}, which gives the fundamental properties of "mixed Hodge structures".
    

    
 \item[(b)]
      The "two filtrations lemma", \hyperref[hodge-theory-ii-1.3.16]{Theorem~\ref*{hodge-theory-ii-1.3.16}}.
    

  \end{enumerate}\par{}\cref{hodge-theory-ii-2} recalls Hodge theory and introduces mixed Hodge structures.\par{}The heart of this work is \cref{hodge-theory-ii-3.2}, which defines the mixed Hodge structure on \(\operatorname {H}^n(X,\mathbb {C})\), and establishes some degenerations of spectral sequences.\par{}\cref{hodge-theory-ii-4} gives diverse applications, all following from \hyperref[hodge-theory-ii-4.1.1]{Theorem~\ref*{hodge-theory-ii-4.1.1}} and the theory of the \((K/k)\)-trace, for the resulting Hodge structures (\hyperref[hodge-theory-ii-4.1.2]{Corollary~\ref*{hodge-theory-ii-4.1.2}}).
  The principal ones are \hyperref[hodge-theory-ii-4.2.6]{Theorem~\ref*{hodge-theory-ii-4.2.6}} and \hyperref[hodge-theory-ii-4.4.15]{Corollary~\ref*{hodge-theory-ii-4.4.15}}.\chapter{Filtrations}\thispagestyle{fancy}\label{hodge-theory-ii-1}\section{Filtered objects}\label{hodge-theory-ii-1.1}\subsection{}\label{hodge-theory-ii-1.1.1}\par{}Let \(\mathscr {A}\) be an abelian category.\par{}We will be considering \(\mathbb {Z}\)-filtrations, finite, in general, on the objects of \(\mathscr {A}\):\begin{definition}{1.1.2}\setnumber{1.1.2}\label{hodge-theory-ii-1.1.2}\skipsinglepar\par{}A \emph{decreasing} (resp. \emph{increasing}) \emph{filtration} \(F\) of an object \(A\) of \(\mathscr {A}\) is a family \((F^n(A))_{n\in \mathbb {Z}}\) (resp. \((F_n(A))_{n\in \mathbb {Z}}\)) of sub-objects of \(A\) satisfying
    \[       \forall  n,m       \quad  n\leqslant  m \implies  F^m(A)\subset  F^n(A)     \]
    (resp. \(n\leqslant  m\implies  F_n(A)\subset  F_m(A)\)).\par{}A \emph{filtered object} is an object endowed with a filtration.
    When there is no chance of confusion, we often denote by the same letter filtrations on different objects of \(\mathscr {A}\).\par{}If \(F\) is a decreasing (resp. increasing) filtration on \(A\), then we set \(F^\infty (A)=0\) and \(F^{-\infty }(A)=A\) (resp. \(F_{-\infty }(A)=0\) and \(F_\infty (A)=A\)).\par{}The \emph{shifted filtrations} of a decreasing filtration \(W\) are defined by
    \[       W[n]^p(A)       = W^{n+p}(A)     \]
    for \(n\in \mathbb {Z}\).\end{definition}\subsection{}\label{hodge-theory-ii-1.1.3}\par{}If \(R\) is a decreasing (resp. increasing) filtration of \(A\), then the \(F_n(A)=F^{-n}(A)\) (resp the \(F^n(A)=F_{-n}(A)\)) form an increasing (resp. decreasing) filtration of \(A\).
    
    This allows us in principal to consider only decreasing filtrations;
    \emph{unless otherwise explicitly mentioned, when we say "filtration" we always mean "decreasing filtration"}.\label{hodge-theory-ii-1.1.4}\par{}A filtration \(F\) of \(A\) is said to be \emph{finite} if there exist \(n\) and \(m\) such that \(F^n(A)=A\) and \(F^m(A)=0\).\label{hodge-theory-ii-1.1.5}\par{}A \emph{morphism} from a filtered object \((A,F)\) to a filtered object \((B,F)\) is a morphism \(f\) from \(A\) to \(B\) that satisfies \(f(F^n(A))\subset  F^n(B)\) for all \(n\in \mathbb {Z}\).\par{}Filtered objects (resp. finite filtered objects) of \(\mathscr {A}\) form an additive category in which inductive limits and finite projective limits exist (and thus kernels, cokernels, images, and coimages of a morphism).\par{}A morphism \(f\colon (A,F)\to (B,F)\) is said to be \emph{strict}, or \emph{strictly compatible with the filtrations}, if the canonical arrows from \(\operatorname {Coim}(f)\) to \(\operatorname {Im}(f)\) is an isomorphism of filtered objects (cf. \hyperref[hodge-theory-ii-1.1.11]{Proposition~\ref*{hodge-theory-ii-1.1.11}}).\label{hodge-theory-ii-1.1.6}\par{}Let \((-)^\circ \) be the contravariant identity functor from \(\mathscr {A}\) to the dual category \(\mathscr {A}^\circ \).
    If \((A,F)\) is a filtered object of \(\mathscr {A}\), then the \((A/F^n(A))^\circ \) can be identified with sub-objects of \(A^\circ \).
    The \emph{dual} filtration on \(A^\circ \) is defined by
    \[       F^n(A^\circ )       = (A/F^{1-n})^\circ .     \]\par{}The double dual of \((A,F)\) can be identified with \((A,F)\).
    This construction identifies the dual of the category of filtered objects of \(\mathscr {A}\) with the category of filtered objects of \(\mathscr {A}^\circ \).\label{hodge-theory-ii-1.1.7}\par{}If \((A,F)\) is a filtered object of \(\mathscr {A}\), then its \emph{associated graded} is the object of \(\mathscr {A}^\mathbb {Z}\) defined by
    \[       \operatorname {Gr}^n(A)       = F^n(A)/F^{n+1}(A).     \]
    The convention \cref{hodge-theory-ii-1.1.6} is justified by the simple formula
    \[       \operatorname {Gr}^n(A^\circ )       = \operatorname {Gr}^{-n}(A)^\circ      \]
    which follows from the self-dual diagram

    \begin{eqn}{1.1.7.1}\setnumber{1.1.7.1}\label{hodge-theory-ii-1.1.7.1}\skipsinglepar\end{eqn}\label{hodge-theory-ii-1.1.8}\par{}Let \((A,F)\) be a filtered object, and \(j\colon  X\hookrightarrow  A\) a sub-object of \(A\).
    The \emph{filtration induced by \(F\)} (or simply the \emph{induced filtration}) on \(X\) is the unique filtration on \(X\) such that \(j\) is strictly compatible with the filtrations;
    we have
    \[       F^n(X)       = j^{-1}(F^n(A))       = X\cap  F^n(A).     \]\par{}Dually, the \emph{quotient filtration} on \(A/X\) (the unique filtration such that \(p\colon  A\to  A/X\) is strictly compatible with the filtrations) is given by
    \[       F^n(A/X)       = p(F^n(A))       \cong  (X+F^n(A))/X       \cong  F^n(A)/(X\cap  F^n(A)).     \]\begin{lemma}{1.1.9}\setnumber{1.1.9}\label{hodge-theory-ii-1.1.9}\skipsinglepar\par{}If \(X\) and \(Y\) are sub-objects of \(A\), with \(X\subset  Y\), then on \(Y/X\xrightarrow {\sim }\operatorname {Ker}(A/X\to  A/Y)\) the quotient filtration of \(Y\) agrees with that induced by that of \(A/X\).
    In the diagram
    
  \begin {tikzcd}
      A/Y
      && A/X
        \ar [ll]
    \\& A
        \ar [ul] \ar [ur]
      && Y/X
        \ar [ul]
    \\X
        \ar [rr] \ar [ur]
      && Y
        \ar [ul] \ar [uu] \ar [ur]
    \end {tikzcd}

    the arrows are strict.\end{lemma}\label{hodge-theory-ii-1.1.10}\par{}We call the filtration \hyperref[hodge-theory-ii-1.1.9]{Lemma~\ref*{hodge-theory-ii-1.1.9}} on \(Y/X\) the \emph{filtration induced by that of \(A\)} (or simply the \emph{induced filtration}).
    By \hyperref[hodge-theory-ii-1.1.9]{Lemma~\ref*{hodge-theory-ii-1.1.9}}, its definition is self-dual.

    In particular, if \(\Sigma \colon  A\xrightarrow {f}B\xrightarrow {G}C\) is a **!!\textbf{\emph{TO-DO: o-suite}}?!!** sequence, and if \(B\) is filtered, then \(\operatorname {H}(\Sigma )=\operatorname {Ker}(g)/\operatorname {Im}(f)=\operatorname {Ker}(\operatorname {Coker}(f)\to \operatorname {Coim}(g))\) is endowed with a canonical induced filtration.\par{}The reader can show that:\begin{proposition}{1.1.11}\setnumber{1.1.11}\label{hodge-theory-ii-1.1.11}\skipsinglepar\begin{enumerate}
    
 \item[(i)]
      Let \(f\colon (A,F)\to (B,F)\) be a morphism of filtered objects with finite filtrations.
      For \(f\) to be strict, it is necessary and sufficient that the sequence
      \[         0         \to  \operatorname {Gr}(\operatorname {Ker}(f))         \to  \operatorname {Gr}(A)         \to  \operatorname {Gr}(B)         \to  \operatorname {Gr}(\operatorname {Coker}(f))         \to  0       \]
      be exact.
    


    
 \item[(ii)]
      Let \(\sigma \colon (A,F)\to (B,F)\to (C,F)\) be a **!!\textbf{\emph{TO-DO: o-suite}}?!!** sequence of strict morphisms.
      We then have
      \[         \operatorname {H}(\operatorname {Gr}(\Sigma ))         \cong  \operatorname {Gr}(\operatorname {H}(\Sigma ))       \]
      canonically.
      In particular, if \(\Sigma \) is exact in \(\mathscr {A}\), then \(\operatorname {Gr}(\Sigma )\) is exact in \(\mathscr {A}^\mathbb {Z}\).
    

  \end{enumerate}\end{proposition}\par{}In a category of modules, to say that a morphism \(f\colon (A,F)\to (B,F)\) is strict implies that every \(b\in  B\) of filtration \(\geqslant  n\) (i.e. \(b\in  F^n(B)\)) that is in the image of \(A\) is already in the image of \(F^n(A)\):
  \[     f(F^n(A))     = f(A)\cap  F^n(B).   \]\label{hodge-theory-ii-1.1.12}\par{}If \(\otimes \colon \mathscr {A}_1\times \ldots \times \mathscr {A}_n\to \mathscr {B}\) is a multiadditive right-exact functor, and if \(A_i\) is an object of finite filtration of \(\mathscr {A}_i\) (for \(1\leqslant  i\leqslant  n\)), then we define a filtration on \(\bigotimes _{i=1}^n A_i\) by
    \[       F^k(\bigotimes _{i=1}^n A_i)       = \sum _{\sum  k_i=k} \operatorname {Im}(\bigotimes _{i=1}^n F^{k_i}(A_i)\to \bigotimes _{i=1}^n A_i)     \]
    (a sum of sub-objects).\par{}Dually, if \(H\) is left-exact, then we set
    \[       F^k(H(A_i))       = \bigcap _{\sum  k_i=k} \operatorname {Ker}(H(A_i)\to  H(A_i/F^{k_i}(A_i))).     \]\par{}If \(H\) is exact, then the two definitions are equivalent.\par{}We extend these definitions to contravariant functors in certain variables by \cref{hodge-theory-ii-1.1.6}.
    In particular, for the left-exact functor \(\operatorname {Hom}\), we set
    \[       F^k(\operatorname {Hom}(A,B))       = \{f\colon  A\to  B \mid  f(F^n(A))\subset  F^{n+k}(B)\,\,\forall  n\}.     \]
    We thus have
    \[       \operatorname {Hom}((A,F),(B,F))       = F^0(\operatorname {Hom}(A,B)).     \]

    Under the above hypotheses, we have obvious morphisms
    \[       \begin {aligned}         \bigotimes _{i=1}^n \operatorname {Gr}(A_i)         &\to  \operatorname {Gr}(\bigotimes _{i=1}^n A_i)       \\\operatorname {Gr} H(A_i)         &\to  H(\operatorname {Gr}(A_i)).       \end {aligned}     \]
    If \(H\) is exact, then these are isomorphisms and inverse to one another.\par{}These constructions are compatible with composition of functors, in a sense whose details we leave to the reader.\section{Opposite filtrations}\label{hodge-theory-ii-1.2}\label{hodge-theory-ii-1.2.1}\par{}Let \(A\) be an object of \(\mathscr {A}\) endowed with filtrations \(F\) and \(G\).
    By definition, \(\operatorname {Gr}_F^n(A)\) is a quotient of a sub-object of \(A\) and, as such, is endowed with a filtration induced by \(G\) \cref{hodge-theory-ii-1.1.10}.
    Passing to the associated graded defines a bigraded object \((\operatorname {Gr}_G^n\operatorname {Gr}_F^m(A))_{n,m\in \mathbb {Z}}\).
    By a lemma of Zassenhaus, \(\operatorname {Gr}_G^n\operatorname {Gr}_F^m(A)\) and \(\operatorname {Gr}_F^m\operatorname {Gr}_G^n(A)\) are canonically isomorphic: if we define the induced filtrations \cref{hodge-theory-ii-1.1.10} as quotient filtrations of the induced filtrations on a sub-object, then we have
    \[       \begin {aligned}         \operatorname {Gr}_G^n\operatorname {Gr}_F^m(A)         \cong  (F^m(A)\cap  G^n(A)) &/ ((F^{m+1}(A)\cap  G^n(A)) + (F^m(A)\cap  G^{n+1}(A)))       \\&=       \\\operatorname {Gr}_F^m\operatorname {Gr}_G^n(A)         \cong  (G^n(A)\cap  F^m(A)) &/ ((G^{n+1}(A)\cap  F^m(A)) + (G^n(A)\cap  F^{m+1}(A)))       \end {aligned}     \]\label{hodge-theory-ii-1.2.2}\par{}Let \(H\) be a third filtration of \(A\).
    It induces a filtration on \(\operatorname {Gr}_F(A)\), and thus on \(\operatorname {Gr}_G\operatorname {Gr}_F(A)\).
    It also induces a filtration on \(\operatorname {Gr}_F\operatorname {Gr}_G(A)\).
    We note that these filtrations do not in general correspond to one another under the isomorphism \cref{hodge-theory-ii-1.2.1}.
    In the expression \(\operatorname {Gr}_H\operatorname {Gr}_G\operatorname {Gr}_F(A)\), \(G\) and \(H\) thus play a symmetric role, but not \(F\) and \(G\).\begin{definition}{1.2.3}\setnumber{1.2.3}\label{hodge-theory-ii-1.2.3}\skipsinglepar\par{}Two \emph{finite} filtrations \(F\) and \(\overline {F}\) on \(A\) are said to be \emph{\(n\)-opposite} if \(\operatorname {Gr}_F^p\operatorname {Gr}_{\overline {F}}^q(A)=0\) for \(p+q\neq  n\).\end{definition}\label{hodge-theory-ii-1.2.4}\par{}If \(A^{p,q}\) is a bigraded object of \(\mathscr {A}\) such that

    
 \begin{enumerate}
      
 \item[(a)]
        \(A^{p,q}=0\) except for a finite number of pairs \((p,q)\), and
      

      
 \item[(b)]
        \(A^{p,q}=0\) for \(p+q\neq  n\)
      

    \end{enumerate}


    
    then we define two \(n\)-opposite filtrations of \(A=\sum _{p,q}A^{p,q}\) by setting

    \begin{eqn}{1.2.4.1}\setnumber{1.2.4.1}\label{hodge-theory-ii-1.2.4.1}\skipsinglepar\[         F^p(A)         = \sum _{p'\geqslant  p} A^{p',q'}       \tag{1.2.4.1}       \]\end{eqn}

    \begin{eqn}{1.2.4.2}\setnumber{1.2.4.2}\label{hodge-theory-ii-1.2.4.2}\skipsinglepar\[         \overline {F}^q(A)         = \sum _{q'\geqslant  q} A^{p',q'}.       \tag{1.2.4.2}       \]\end{eqn}

    We have

    \begin{eqn}{1.2.4.3}\setnumber{1.2.4.3}\label{hodge-theory-ii-1.2.4.3}\skipsinglepar\[         \operatorname {Gr}_F^p\operatorname {Gr}_{\overline {F}}^q(A)         = A^{p,q}.       \tag{1.2.4.3}       \]\end{eqn}\par{}Conversely:\begin{proposition}{1.2.5}\setnumber{1.2.5}\label{hodge-theory-ii-1.2.5}\skipsinglepar\begin{enumerate}
    
 \item[(i)]
      Let \(F\) and \(\overline {F}\) be finite filtrations on \(A\).
      For \(F\) and \(\overline {F}\) to be \(n\)-opposite, it is necessary and sufficient that, for all \(p,q\),
      \[         [p+q=n+1]         \implies  [F^p(A)\oplus \overline {F}^q(A) \xrightarrow {\sim } A].       \]
    


    
 \item[(ii)]
      If \(F\) and \(\overline {F}\) are \(n\)-opposite, and if we set
      \[         \begin {cases}           A^{p,q} = 0           &\text {for }p+q\neq  n         \\A^{p,q} = F^p(A)\cap \overline {F}^q(A)           &\text {for }p+q=n         \end {cases}       \]
      then \(A\) is the direct sum of the \(A^{p,q}\), and \(F\) and \(\overline {F}\) come from the bigrading \(A^{p,q}\) of \(A\) by the procedure of \cref{hodge-theory-ii-1.2.4}.
    


    
  
    
    \begin{proof}\begin{enumerate}
        
 \item[(i)]
          The condition \(\operatorname {Gr}_F^p\operatorname {Gr}_{\overline {F}}^q(A)=0\) for \(p+q>n\) implies that \(F^p\cap \overline {F}^q=(F^{p+1}\cap \overline {F}^q)+(F^p\cap \overline {F}^{q+1})\) for \(p+q>n\).
          By hypothesis, \(F^p\cap \overline {F}^q\) is zero for large enough \(p+q\);
          by decreasing induction, we thus deduce that the condition \(\operatorname {Gr}_F^p\operatorname {Gr}_{\overline {F}}^q(A)=0\) for \(p+q>n\) is equivalent to the condition \(F^p(A)\cap \overline {F}^q(A)=0\) for \(p+q>n\).
          Dually (\cref{hodge-theory-ii-1.1.6}, \cref{hodge-theory-ii-1.1.7}, \cref{hodge-theory-ii-1.1.10}), the condition \(\operatorname {Gr}_F^p\operatorname {Gr}_{\overline {F}}^q(A)=0\) for \(p+q<n\) is equivalent to the condition \(A=F^p(A)+\overline {F}^q(A)\) for \((1-p)+(1-q)>-n\), i.e. for \(p+q\leqslant  n+1\), and the claim then follows.
        


        
 \item[(ii)]
          If \(F\) and \(\overline {F}\) are \(n\)-opposite, then we can prove by decreasing induction on \(p\) that

          \begin{eqn}{1.2.5.1}\setnumber{1.2.5.1}\label{hodge-theory-ii-1.2.5.1}\skipsinglepar\[               \bigoplus _{p'\geqslant  p} A^{p',q'}               \xrightarrow {\sim } F^p(A).             \tag{1.2.5.1}             \]\end{eqn}

          
 \begin{enumerate}
            
 \item[(a)]
              For \(F^p(A)=0\), the claim is evident.
            

            
 \item[(b)]
              The decomposition \(A=F^{p+1}(A)\oplus \overline {F}^{n-p}(A)\) induces on \(F^p(A)\supset  F^{p+1}(A)\) a decomposition
              \[                 F^p(A)                 = F^{p+1}(A)\oplus (F^p(A)\cap \overline {F}^{n-p}(A))               \]
              and we conclude by induction.
            

          \end{enumerate}


          For \(p\) small enough, we have \(F^p(A)=A\).
          By \hyperref[hodge-theory-ii-1.2.5.1]{Equation~\ref*{hodge-theory-ii-1.2.5.1}}, the \(A^{p,q}\) thus form a bigrading of \(A\), and \(F\) satisfies \hyperref[hodge-theory-ii-1.2.4.1]{Equation~\ref*{hodge-theory-ii-1.2.4.1}}.
          The fact that \(\overline {F}\) satisfies \hyperref[hodge-theory-ii-1.2.4.2]{Equation~\ref*{hodge-theory-ii-1.2.4.2}} then follows by symmetry.
        

      \end{enumerate}\end{proof}
  

  \end{enumerate}\end{proposition}\label{hodge-theory-ii-1.2.6}\par{}The constructions \cref{hodge-theory-ii-1.2.4} and \hyperref[hodge-theory-ii-1.2.5]{Proposition~\ref*{hodge-theory-ii-1.2.5}} establish equivalences of categories that are quasi-inverse to one another between objects of \(\mathscr {A}\) endowed with two finite \(n\)-opposite filtrations and bigraded objects of \(\mathscr {A}\) of the type considered in \cref{hodge-theory-ii-1.2.4}.\begin{definition}{1.2.7}\setnumber{1.2.7}\label{hodge-theory-ii-1.2.7}\skipsinglepar\par{}Three \emph{finite} filtrations \(W\), \(F\), and \(\overline {F}\) on \(A\) are said to be \emph{opposite} if
    \[       \operatorname {Gr}_F^p\operatorname {Gr}_{\overline {F}}^q\operatorname {Gr}_W^n(A)       = 0     \]
    for \(p+q+n\neq 0\).\end{definition}\par{}This condition is symmetric in \(F\) and \(\overline {F}\).
  It implies that \(F\) and \(\overline {F}\) induce on \(W^n(A)/W^{n+1}(A)\) two \((-n)\)-opposite filtrations.
  We set
  \[     A^{p,q}     = \operatorname {Gr}_F^p\operatorname {Gr}_{\overline {F}}^q\operatorname {Gr}_F^{-p-q}(A)   \]
  whence decompositions \cref{hodge-theory-ii-1.2.4}, \hyperref[hodge-theory-ii-1.2.5]{Proposition~\ref*{hodge-theory-ii-1.2.5}}
  \[     W^n(A)/W^{n+1}(A)     = \bigoplus _{p+q=-n} A^{p,q}   \tag{1.2.7.1}   \]
  which makes \(\operatorname {Gr}_W(A)\) into a bigraded object.\begin{lemma}{1.2.8}\setnumber{1.2.8}\label{hodge-theory-ii-1.2.8}\skipsinglepar\par{}Let \(W\), \(F\), and \(\overline {F}\) be three finite opposite filtrations, and \(\sigma \) a sequence \((p_i,q_i)_{i\geqslant 0}\) pairs of integers satisfying

    
 \begin{enumerate}
      
 \item[(a)]
        \(p_i\leqslant  p_j\) and \(q_i\leqslant  q_j\) for \(i\geqslant  j\), and
      

      
 \item[(b)]
        \(p_i+q_i=p_0+q_0-i+1\) for \(i>0\).
      

    \end{enumerate}


    Set \(p=p_0\), \(q=q_0\), \(n=-p-q\), and
    \[       A_\sigma        = \left ( \sum _{0\leqslant  i}(W^{n+i}(A)\cap  F^{p_i}(A)) \right )       \cap  \left ( \sum _{0\leqslant  i}(W^{n+i}(A)\cap \overline {F}^{q_i}(A)) \right ).     \]
    Then the projection from \(W^n(A)\) to \(\operatorname {Gr}_W^n(A)\) induces an isomorphism
    \[       A_\sigma        \xrightarrow {\sim } A^{p,q} \subset  \operatorname {Gr}_W^n(A).     \]\begin{proof}\par{}We will prove by induction on \(k\) the following claim:\par{}\emph{The projection from \(W^n(A)/W^{n+k}\) to \(\operatorname {Gr}_W^n(A)\) induces an isomorphism from
        \[           \begin {aligned}             \Bigg [             &\left ( \sum _{i<k} (W^{n+i}(A)\cap  F^{p_i}(A)) + W^{n+k}(A) \right )           \\\cap              &\left ( \sum _{i<k} (W^{n+i}(A)\cap \overline {F}^{q_i}(A)) + W^{n+k}(A) \right )             \Bigg ] / W^{n+k}(A)           \end {aligned}           \tag{$*_k$}         \]
        to \(A^{p,q}\subset \operatorname {Gr}_W^n(A)\).}\par{}For \(k=1\), this is exactly the definition of \(A^{p,q}\).\par{}By (i) of \hyperref[hodge-theory-ii-1.2.5]{Proposition~\ref*{hodge-theory-ii-1.2.5}} we have

      \begin{eqn}{1.2.8.1}\setnumber{1.2.8.1}\label{hodge-theory-ii-1.2.8.1}\skipsinglepar\[           F^{p_k}(\operatorname {Gr}_W^{n+k}(A)) \oplus  \overline {F}^{q_k}(\operatorname {Gr}_W^{n+k}(A))           \xrightarrow {\sim } \operatorname {Gr}_W^{n+k}(A).         \tag{1.2.8.1}         \]\end{eqn}
      Set
      \[         \begin {aligned}           B &= \sum _{i<k} (W^{n+i}(A)\cap  F^{p_i}(A))         \\C &= \sum _{i<k} (W^{n+i}(A)\cap \overline {F}^{q_i}(A))         \\B' &= (W^{n+k}(A)\cap  F^{p_k}(A)) + W^{n+k+1}(A)         \\C' &= (W^{n+k}(A)\cap \overline {F}^{q_k}(A)) + W^{n+k+1}(A)         \\D &= W^{n+k}(A)         \\E &= W^{n+k+1}(A).         \end {aligned}       \]\par{}\hyperref[hodge-theory-ii-1.2.8.1]{Equation~\ref*{hodge-theory-ii-1.2.8.1}} can then be written as
      \[         \begin {aligned}           B'+C' &= D         \\B'\cap  C' &= E.         \end {aligned}       \]
      We also have, since \(p_k\leqslant  p_i\) (for \(i\leqslant  k\)),
      \[         B\cap  D         \subset  F^{p_k}(A) \cap  W^{n+k}(A)         \subset  B'       \]
      and, since \(q_k\leqslant  q_i\) (for \(i\leqslant  k\)),
      \[         C\cap  D         \subset  \overline {F}^{q_k}(A)\cap  W^{n+k}(A)         \subset  C'.       \]\par{}The claim (\(*_{k+1}\)) and then follows from (\(*_k\)) and the following lemma.\end{proof}\end{lemma}\begin{lemma}{1.2.9}\setnumber{1.2.9}\label{hodge-theory-ii-1.2.9}\skipsinglepar\par{}Let \(B\), \(C\), \(B'\), \(C'\), \(D\), and \(E\) be sub-objects of \(A\).
    Suppose that
    \[       \begin {gathered}         B'+C' = D         \qquad  B'\cap  C' = E       \\B\cap  D \subset  B'         \qquad  C\cap  D \subset  C'.       \end {gathered}     \]
    Then
    \[       ((B+B')\cap (C+C'))/E       \xrightarrow {\sim } ((B+D)\cap (C+D))/D.     \]\begin{proof}\par{}To prove surjectivity, we write
      \[         \begin {aligned}           ((B+B')\cap (C+C'))+D           &= (((B+B')\cap (C+C'))+B')+C'         \\&= ((B+B')\cap (C+C'+B'))+C'         \\&= (B+B'+C')\cap (C+C'+B')         \\&= (B+D)\cap (C+D).         \end {aligned}       \]\par{}To prove injectivity, we write
      \[         (B+B')\cap (C+C')\cap  D         = ((B+B')\cap  D)\cap ((C+C')\cap  D).       \]
      Since \(B'\subset  D\), we have
      \[         \begin {aligned}           (B+B')\cap  D           &= (B\cap  D)+B'         \\&= B'         \end {aligned}       \]
      and similarly
      \[         (C+C')\cap  D         = C'       \]
      and
      \[         \begin {aligned}           (B+B')\cap (C+C')\cap  D           &= B'\cap  C'         \\&= E.         \end {aligned}       \]\end{proof}\end{lemma}\par{}This finishes the proof of \hyperref[hodge-theory-ii-1.2.8]{Lemma~\ref*{hodge-theory-ii-1.2.8}}, noting that \hyperref[hodge-theory-ii-1.2.8]{Lemma~\ref*{hodge-theory-ii-1.2.8}} is equivalent to (\(*_k\)) for large \(k\).\begin{theorem}{1.2.10}\setnumber{1.2.10}\label{hodge-theory-ii-1.2.10}\skipsinglepar\par{}Let \(\mathscr {A}\) be an abelian category, and provisionally denote by \(\mathscr {A}'\) the category of objects of \(\mathscr {A}\) endowed with three opposite filtrations \(W\), \(F\), and \(\overline {F}\).
    The morphisms in \(\mathscr {A}'\) are the morphisms of \(\mathscr {A}\) that are compatible with the three filtrations.

    
 \begin{enumerate}
      
 \item[(i)]
        \(\mathscr {A}'\) is an abelian category.
      

      
 \item[(ii)]
        The kernel (resp. cokernel) of an arrow \(f\colon  A\to  B\) in \(\mathscr {A}'\) is the kernel (resp. cokernel) of \(f\) in \(\mathscr {A}\), endowed with the filtrations induced by those of \(A\) (resp. the quotients of those of \(B\)).
      

      
 \item[(iii)]
        Every morphism \(f\colon  A\to  B\) in \(\mathscr {A}'\) is strictly compatible with the filtrations \(W\), \(F\), and \(\overline {F}\);
        the morphism \(\operatorname {Gr}_W(f)\) is compatible with the bigradings of \(\operatorname {Gr}_W(A)\) and \(\operatorname {Gr}_W(B)\);
        the morphisms \(\operatorname {Gr}_F(f)\) and \(\operatorname {Gr}_{\overline {F}}(f)\) are strictly compatible with the filtration induced by \(W\).
      

      
 \item[(iv)]
        The "forget the filtrations" functors, \(\operatorname {Gr}_W\), \(\operatorname {Gr}_F\), and \(\operatorname {Gr}_{\overline {F}}\), and
        \[           \begin {gathered}             \operatorname {Gr}_W\operatorname {Gr}_F             \simeq  \operatorname {Gr}_F\operatorname {Gr}_W           \\\simeq  \operatorname {Gr}_{\overline {F}}\operatorname {Gr}_F\operatorname {Gr}_W           \\\simeq  \operatorname {Gr}_{\overline {F}}\operatorname {Gr}_W             \simeq  \operatorname {Gr}_W\operatorname {Gr}_{\overline {F}}           \end {gathered}         \]
        from \(\mathscr {A}'\) to \(\mathscr {A}\) are exact.
      

    \end{enumerate}\end{theorem}\par{}Denote by \(\sigma _0(p,q)\) and \(\sigma _1(p,q)\) the sequences
  \[     \begin {aligned}       \sigma _0(p,q)       &= (p,q), (p,q), (p,q-1), (p,q-2), (p,q-3), \ldots      \\\sigma _0(p,q)       &= (p,q), (p,q), (p-1,q), (p-2,q), (p-3,q), \ldots      \end {aligned}   \]
  and, with the notation of \hyperref[hodge-theory-ii-1.2.8]{Lemma~\ref*{hodge-theory-ii-1.2.8}}, set
  \[     A_i^{p,q}     = A_{\sigma _i(p,q)}     \qquad \text {for } i=0,1.   \]\par{}If \(f\colon  A\to  B\) is compatible with \(W\), \(F\), and \(\overline {F}\), then we have

  \begin{eqn}{1.2.10.1}\setnumber{1.2.10.1}\label{hodge-theory-ii-1.2.10.1}\skipsinglepar\[       f(A_i^{p,q})       \subset  B_i^{p,q}       \qquad \text {for }i=0,1.     \tag{1.2.10.1}     \]\end{eqn}

  Claim (iii) then follows from the following lemma:\begin{lemma}{1.2.11}\setnumber{1.2.11}\label{hodge-theory-ii-1.2.11}\skipsinglepar\par{}The \(A_i^{p,q}\) give a bigrading of \(A\).
    We have

    \begin{eqn}{1.2.11.1}\setnumber{1.2.11.1}\label{hodge-theory-ii-1.2.11.1}\skipsinglepar\[         W^n(A)         = \sum _{n+p+q\leqslant  0} A_i^{p,q}         \qquad \text {for }i=0,1       \tag{1.2.11.1}       \]\end{eqn}

    \begin{eqn}{1.2.11.2}\setnumber{1.2.11.2}\label{hodge-theory-ii-1.2.11.2}\skipsinglepar\[         F^p(A)         = \sum _{p'\geqslant  p} A_0^{p',q'}       \tag{1.2.11.2}       \]\end{eqn}

    \begin{eqn}{1.2.11.3}\setnumber{1.2.11.3}\label{hodge-theory-ii-1.2.11.3}\skipsinglepar\[         \overline {F}^q(A)         = \sum _{q'\geqslant  q} A_1^{p',q'}.       \tag{1.2.11.3}       \]\end{eqn}\begin{proof}\par{}By symmetry, it suffices to prove the claims concerning \(i=0\).
      Set \(A_0=\bigoplus  A_0^{p,q}\) and define filtrations \(W\) and \(F\) on \(A_0\) by the equations of \hyperref[hodge-theory-ii-1.2.11]{Lemma~\ref*{hodge-theory-ii-1.2.11}}.
      The canonical map \(i\) from \(A_0\) to \(A\) is compatible with the filtrations \(W\) and \(F\).
      Furthermore, by \hyperref[hodge-theory-ii-1.2.8]{Lemma~\ref*{hodge-theory-ii-1.2.8}}, \(\operatorname {Gr}_W(i)\) is an isomorphism, and induces isomorphisms of graded objects

      \begin{eqn}{1.2.11.4}\setnumber{1.2.11.4}\label{hodge-theory-ii-1.2.11.4}\skipsinglepar\[           \sum _{p+q=n} A_0^{p,q}           \xrightarrow {\sim } \operatorname {Gr}_W^{-n}(A)           = \sum _{p+q=n} A^{p,q}.         \tag{1.2.11.4}         \]\end{eqn}

      The morphism \(i\) is thus an isomorphism, and the \(A_0^{p,q}\) give a bigrading of \(A\).\par{}\hyperref[hodge-theory-ii-1.2.11.1]{Equation~\ref*{hodge-theory-ii-1.2.11.1}} then says that \(\operatorname {Gr}_W(i)\) is an isomorphism.
      By \hyperref[hodge-theory-ii-1.2.11.4]{Equation~\ref*{hodge-theory-ii-1.2.11.4}}, \(\operatorname {Gr}_F\operatorname {Gr}_W(i)\) is an isomorphism, and thus so too are \(\operatorname {Gr}_W\operatorname {Gr}_F(i)\) and \(\operatorname {Gr}_F(i)\).
      Equation \hyperref[hodge-theory-ii-1.2.11.2]{Equation~\ref*{hodge-theory-ii-1.2.11.2}} then follows.\end{proof}\end{lemma}\label{hodge-theory-ii-1.2.12}\par{}We now prove \hyperref[hodge-theory-ii-1.2.10]{Theorem~\ref*{hodge-theory-ii-1.2.10}}.
    Let \(f\colon  A\to  B\) in \(\mathscr {A}'\) and endow \(K=\operatorname {Ker}(f)\) with the filtrations induced by those of \(A\).
    By \hyperref[hodge-theory-ii-1.2.11]{Lemma~\ref*{hodge-theory-ii-1.2.11}}, \(\operatorname {Gr}_W(K)\hookrightarrow \operatorname {Gr}_W(A)\);
    furthermore, the filtration \(F\) (resp. \(\overline {F}\)) on \(K\) induces on \(\operatorname {Gr}_W(K)\) the inverse image filtration of the filtration \(F\) on \(\operatorname {Gr}_W(A)\).
    The sub-object \(\operatorname {Gr}_W(K)\) of \(\operatorname {Gr}_W(A)\) is then compatible with the bigrading of \(\operatorname {Gr}_W(A)\):
    \[       \operatorname {Gr}_W(K)       = \bigoplus _{p,q}(\operatorname {Gr}_W(K)\cap  A^{p,q}).     \]\par{}We thus deduce that
    \[       \operatorname {Gr}_F^p\operatorname {Gr}_{\overline {F}}^q\operatorname {Gr}_W^n(K)       \hookrightarrow  \operatorname {Gr}_F^p\operatorname {Gr}_{\overline {F}}^q\operatorname {Gr}_W^n(A);     \]
    
    the filtrations of \(W\), \(F\), and \(\overline {F}\) on \(K\) are thus opposite, and \(K\) is a kernel of \(f\) in \(\mathscr {A}'\).
    This, combined with the dual result, proves (ii).\par{}If \(f\) is an arrow of \(\mathscr {A}'\), then the canonical morphism from \(\operatorname {Coim}(f)\) to \(\operatorname {Im}(f)\) is an isomorphism in \(\mathscr {A}\);
    by (iii), it is also an isomorphism in \(\mathscr {A}'\), which is thus abelian.\par{}The "forget the filtrations" functor is exact by (ii).
    The exactness of the other functors in (iv) follows immediately from (ii), (iii), and (i) or (ii) of \hyperref[hodge-theory-ii-1.1.11]{Proposition~\ref*{hodge-theory-ii-1.1.11}}.\label{hodge-theory-ii-1.2.13}\par{}Let \(A\) be an object of \(\mathscr {A}\) endowed with a finite \emph{increasing} filtration \(W_\bullet \), and two finite decreasing filtrations \(F\) and \(\overline {F}\).
    The construction \cref{hodge-theory-ii-1.1.3} associates to \(W_\bullet \) a decreasing filtration \(W^\bullet \).
    We say that the filtrations \(W_\bullet \), \(F\), and \(\overline {F}\) are \emph{opposite} if the filtrations \(W^\bullet \), \(F\), and \(\overline {F}\) are, i.e. if, for all \(n\), the filtrations induced by \(F\) and \(\overline {F}\) on
    \[       \operatorname {Gr}_n^W(A)       = W_n(A)/W_{n-1}(A)     \]
    are \(n\)-opposite.\par{}\hyperref[hodge-theory-ii-1.2.10]{Theorem~\ref*{hodge-theory-ii-1.2.10}} translates trivially to this variation.\section{The two filtrations lemma}\label{hodge-theory-ii-1.3}\label{hodge-theory-ii-1.3.1}\par{}Let \(K\) be a differential complex of objects of \(\mathscr {A}\), endowed with a filtration \(F\).
    The filtration is said to be \emph{biregular} if it induces a finite filtration on each component of \(K\).\par{}We recall the definition of the terms \(E_r^{pq}(K,F)\), or simply \(E_r^{pq}\), of the spectral sequence defined by \(F\).
    We set
    \[       Z_r^{pq}       = \operatorname {Ker}(d\colon  F^p(K^{p+q})\to  F^{p+q+1}/F^{p+r}(K^{p+q+1}))     \]
    and dually we define \(B_r^{pq}\) by the formula
    \[       K^{p+q}/B_r^{pq}       = \operatorname {Coker}(d\colon  F^{p-r+1}(K^{p+q+1})\to  K^{p+q}/F^{p+1}(K^{p+q})).     \]
    These formulas still make sense for \(r=\infty \).\par{}We note that the use here of the notation \(B_r^{pq}\) is different to that of Godement \hyperref[G1958]{[G1958]}.\par{}We have, by definition:

    \begin{eqn}{1.3.1.1}\setnumber{1.3.1.1}\label{hodge-theory-ii-1.3.1.1}\skipsinglepar\[         E_r^{pq}         = \operatorname {Im}(Z_r^{pq}\to  K^{p+q}/B_r^{pq})       \tag{1.3.1.1}       \]\end{eqn}

    \begin{eqn}{1.3.1.2}\setnumber{1.3.1.2}\label{hodge-theory-ii-1.3.1.2}\skipsinglepar\[         = Z_r^{pq}/(B_r^{pq}\cap  Z_r^{pq})       \tag{1.3.1.2}       \]\end{eqn}

    \begin{eqn}{1.3.1.3}\setnumber{1.3.1.3}\label{hodge-theory-ii-1.3.1.3}\skipsinglepar\[         = \operatorname {Ker}(K^{p+q}/B_r^{pq}\to  K^{p+q}/(Z_r^{pq}+B_r^{pq})).       \tag{1.3.1.3}       \]\end{eqn}

    We can thus write

    \begin{eqn}{1.3.1.4}\setnumber{1.3.1.4}\label{hodge-theory-ii-1.3.1.4}\skipsinglepar\[         \begin {aligned}           B_r^{p\bullet }\cap  Z_r^{p\bullet }           &\coloneqq  (dF^{p-r+1}+F^{p+1}) \cap  (d^{-1}F^{p+r}\cap  F^p)         \\&= (dF^{p-r+1}\cap  F^p) + (F^{p+1}\cap  d^{-1}F^{p+r})         \end {aligned}       \tag{1.3.1.4}       \]\end{eqn}

    since \(dF^{p-r+1}\subset  d^{-1}F^{p+r}\) and \(F^{p+1}\subset  F^p\).\par{}For \(r<\infty \), the \(E_r\) form a complex graded by the degree \(p-r(p+q)\), and \(E_{r+1}\) can be expressed as the cohomology of this complex:

    \begin{eqn}{1.3.1.5}\setnumber{1.3.1.5}\label{hodge-theory-ii-1.3.1.5}\skipsinglepar\[         E_{r+1}^{pq}         = \operatorname {H}(E_r^{p-r,q+r-1} \xrightarrow {d_r} E_r^{pq} \xrightarrow {d_r} E_r^{p+r,q-r+1}).       \tag{1.3.1.5}       \]\end{eqn}\par{}For \(r=0\), we have

    \begin{eqn}{1.3.1.6}\setnumber{1.3.1.6}\label{hodge-theory-ii-1.3.1.6}\skipsinglepar\[         E_0^{\bullet \bullet }         = \operatorname {Gr}_F^\bullet (K^\bullet ).       \tag{1.3.1.6}       \]\end{eqn}\begin{proposition}{1.3.2}\setnumber{1.3.2}\label{hodge-theory-ii-1.3.2}\skipsinglepar\par{}Let \(K\) be a complex endowed with a biregular filtration \(F\).
    The following conditions are equivalent:

    
 \begin{enumerate}
      
 \item[(i)]
        The spectral sequence defined by \(F\) degenerates (\(E_1=E_\infty \)).
      

      
 \item[(ii)]
        The morphisms \(d\colon  K^i\to  K^{i+1}\) are strictly compatible with the filtrations.
      

    \end{enumerate}\begin{proof}\par{}We will prove this in the case where \(\mathscr {A}\) is a category of modules.
      For fixed \(p\) and \(q\), the hypothesis that the arrows \(d_r\) with domains \(E_r^{pq}\) be zero for \(r\geqslant 1\) implies that, if \(x\in  F^p(K^{p+q})\) satisfies \(dx\in  F^{p+1}(K^{p+q+1})\), then there exists \(y\in  K^{p+q}\) such that \(dy=0\) and such that \(x\) and \(y\) have the same image in \(E_1^{pq}\).
      Modifying \(y\) by a boundary, and setting \(z=x-y\), we then have
      \[         \forall  x\in  F^p(K^{p+q})         \left [           dx\in  F^{p+1}(K^{p+q+1})           \implies  \exists  z\text { s.t. } z\in  F^{p+1}(K^{p+q}) \text { and } dz=dx         \right ]       \]
      or, in other words,
      \[         F^{p+1}(K^{p+q+1}) \cap  dF^p(K^{p+q})         = dF^{p+1}(K^{p+q}).       \tag{1}       \]\par{}If this condition is satisfied for arbitrary \(p\) and \(q\), then by induction on \(r\) we have
      \[         F^{p+r} \cap  dF^p         = dF^{p+r}       \]
      which, for large \(p+r\), can be written as
      \[         F^p \cap  dK         = dF^p.       \tag{2}       \]\par{}Claim (2) trivially implies (1), and is equivalent to (ii), which proves the proposition.\end{proof}\end{proposition}\label{hodge-theory-ii-1.3.3}\par{}If \((K,F)\) is a filtered complex, we denote by \(\operatorname {Dec}(K)\) the complex \(K\) endowed with the \emph{shifted filtration}
    \[       \operatorname {Dec}(F)^p K^n       = Z_1^{p+n,-p}.     \]\par{}This filtration is compatible with the differentials:
    \[       \begin {aligned}         dZ_1^{p+n,-p}         &\subset  F^{p+n+1}(K^{n+1}) \cap  \operatorname {Ker}(d)       \\&\subset  Z_\infty ^{p+n+1,-p}       \\&\subset  Z_1^{p+n+1,-p}.       \end {aligned}     \]

    Since

    \begin{eqn}{1.3.3.1}\setnumber{1.3.3.1}\label{hodge-theory-ii-1.3.3.1}\skipsinglepar\[         \begin {aligned}           Z_1^{p+1+n,-p-1}           &\subset  F^{p+1+n}(K^n)         \\&\subset  B_1^{p+n,-p}         \\&\subset  Z_1^{p+n,-p}         \end {aligned}       \tag{1.3.3.1}       \]\end{eqn}

    the evident arrow from \(Z_1^{p+n,-p}/Z_1^{p+1+n,-p-1}\) to \(Z_1^{p+n,-p}/B_1^{p+n,-p}\) is a morphism

    \begin{eqn}{1.3.3.2}\setnumber{1.3.3.2}\label{hodge-theory-ii-1.3.3.2}\skipsinglepar\[         u\colon  E_0^{p,n-p}(\operatorname {Dec}(K)) \to  E_1^{p+n,-p}(K).       \tag{1.3.3.2}       \]\end{eqn}\begin{proposition}{1.3.4}\setnumber{1.3.4}\label{hodge-theory-ii-1.3.4}\skipsinglepar\begin{enumerate}
    
 \item[(i)]
      The morphisms in \hyperref[hodge-theory-ii-1.3.3.2]{Equation~\ref*{hodge-theory-ii-1.3.3.2}} form a morphism of graded complexes from \(E_0(\operatorname {Dec}(K))\) to \(E_1(K)\).
    

    
 \item[(ii)]
      This morphism induces an isomorphism on cohomology.
    

    
 \item[(iii)]
      This morphism induces step-by-step (via \hyperref[hodge-theory-ii-1.3.1.5]{Equation~\ref*{hodge-theory-ii-1.3.1.5}}) isomorphisms of graded complexes \(E_r(\operatorname {Dec}(K))\xrightarrow {\sim } E_{r+1}(K)\) (for \(r\geqslant 1\)).
    

  \end{enumerate}\begin{proof}\par{}Let \(F'\) be the filtration on \(K\) defined by
      \[         {F'}^p(K^n)         = \operatorname {Dec}(F)^{p-n}(K^n)         = Z_1^{p,n-p}.       \]
      We trivially have isomorphisms

      \begin{eqn}{1.3.4.1}\setnumber{1.3.4.1}\label{hodge-theory-ii-1.3.4.1}\skipsinglepar\[           E_r^{p,n-p}(\operatorname {Dec}(K))           = E_{r+1}^{p+n,-p}(K,F')         \tag{1.3.4.1}         \]\end{eqn}

      that are compatible with the \(d_r\) and with \hyperref[hodge-theory-ii-1.3.1.5]{Equation~\ref*{hodge-theory-ii-1.3.1.5}}.
      The map \(u\) comes from \hyperref[hodge-theory-ii-1.3.4.1]{Equation~\ref*{hodge-theory-ii-1.3.4.1}} and from the identity map
      \[         (K,F') \to  (K,F).       \]
      This proves (i), and it remains to show that, for \(r\geqslant 2\),
      \[         E_r^{pq}(K,F')         \xrightarrow {\sim } E_r^{pq}(K,F).       \]
      We have
      \[         \begin {aligned}           Z_r^{pq}(K,F')           &= Z_r^{pq}(K,F)           \qquad \text {for }r\geqslant 1         \\Z_r^{pq}(K,F') \cap  B_r^{pq}(K,F')           &= Z_r^{pq}(K,F) \cap  B_r^{pq}(K,F)           \qquad \text {for }r\geqslant 2         \end {aligned}       \]
      and we can then apply \hyperref[hodge-theory-ii-1.3.1.2]{Equation~\ref*{hodge-theory-ii-1.3.1.2}}.\end{proof}\end{proposition}\label{hodge-theory-ii-1.3.5}\par{}The construction in \cref{hodge-theory-ii-1.3.3} is not self-dual.
    The dual construction consists of defining
    \[       \operatorname {Dec}^\bullet (F)^pK^n       = B_1^{p+n-1,-p+1}.     \]
    We then have morphisms
    \[       E_0^{p,n-p}(\operatorname {Dec}(K))       \to  E_1^{p+n,p}(K)       \to  E_0^{p,n-p}(\operatorname {Dec}^\bullet (K))     \]
    and, for \(r\geqslant 1\), isomorphisms
    \[       E_r^{p,n-p}(\operatorname {Dec}(K))       \xrightarrow {\sim } E_{r+1}^{p+n,p}(K)       \xrightarrow {\sim } E_{r}^{p,n-p}(\operatorname {Dec}^\bullet (K)).     \]\par{}Recall that a morphism of complexes is said to be a \emph{quasi-isomorphism} if it induces an isomorphism on cohomology.\begin{definition}{1.3.6}\setnumber{1.3.6}\label{hodge-theory-ii-1.3.6}\skipsinglepar\begin{enumerate}
    
 \item[(i)]
      A morphism \(f\colon (K,F)\to (K',F')\) of filtered complexes with biregular filtrations is a \emph{filtered quasi-isomorphism} if \(\operatorname {Gr}_F(f)\) is a quasi-isomorphism, i.e. if the \(E_1^{pq}(f)\) are isomorphisms.
    

    
 \item[(ii)]
      A morphism \(f\colon (K,F,W)\to (K,F',W')\) of biregular bifiltered complexes is a \emph{bifiltered quasi-isomorphism} if \(\operatorname {Gr}_F\operatorname {Gr}_W(f)\) is a quasi-isomorphism.
    

  \end{enumerate}\end{definition}\label{hodge-theory-ii-1.3.7}\par{}Let \(K\) be a differential complex of objects of \(\mathscr {A}\), endowed with two filtrations \(F\) and \(W\).
    Let \(E_r^{pq}\) be the spectral sequence defined by \(W\).
    The filtration \(F\) induces on \(E_r^{pq}\) various filtrations, which we will compare.\label{hodge-theory-ii-1.3.8}\par{}\hyperref[hodge-theory-ii-1.3.1.2]{Equation~\ref*{hodge-theory-ii-1.3.1.2}} identifies \(E_r^{pq}\) with a quotient of a sub-object of \(K^{p+q}\).
    The \(E_r^{pq}\) term is thusly given by endowing with a filtration \(F_d\) induced by \(F\), called the \emph{first direct filtration}.\label{hodge-theory-ii-1.3.9}\par{}Dually, \hyperref[hodge-theory-ii-1.3.1.3]{Equation~\ref*{hodge-theory-ii-1.3.1.3}} identifies \(E_r^{pq}\) with a sub-object of a quotient of \(K^{p+q}\), whence a new filtration \(F_{d^*}\) induced by \(F\), called the \emph{second direct filtration}.\begin{lemma}{1.3.10}\setnumber{1.3.10}\label{hodge-theory-ii-1.3.10}\skipsinglepar\par{}On \(E_0\) and \(E_1\), we have \(F_d=F_{d^*}\).\begin{proof}\par{}For \(r=0,1\), we have \(B_r^{pq}\subset  Z_r^{pq}\), and we apply \hyperref[hodge-theory-ii-1.1.9]{Lemma~\ref*{hodge-theory-ii-1.1.9}}.\end{proof}\end{lemma}\label{hodge-theory-ii-1.3.11}\par{}\hyperref[hodge-theory-ii-1.3.1.5]{Equation~\ref*{hodge-theory-ii-1.3.1.5}} identifies \(E_{r+1}^{pq}\) with a quotient of a sub-object of \(E_r^{pq}\).
    We define the \emph{recurrent filtration} \(F_r\) on the \(E_r^{pq}\) by the conditions

    
 \begin{enumerate}
      
 \item[(i)]
        On \(E_0^{pq}\), \(F_r=F_d=F_{d^*}\).
      

      
 \item[(ii)]
        On \(E_{r+1}^{pq}\), the recurrent filtration is that induced by the recurrent filtration of \(E_r^{pq}\).
      

    \end{enumerate}\label{hodge-theory-ii-1.3.12}\par{}Definitions \cref{hodge-theory-ii-1.3.8} and \cref{hodge-theory-ii-1.3.9} still make sense for \(r=\infty \).
    If the filtration on \(K\) is biregular, then the direct filtrations on \(E_\infty ^{pq}\) coincide with those on \(E_r^{pq}=E_\infty ^{pq}\) for large enough \(r\), and we define the recurrent filtration on \(E_\infty ^{pq}\) as agreeing with that on \(E_r^{pq}\) for large enough \(r\).\par{}The filtrations \(F\) and \(W\) each induce a filtration on \(H^\bullet (K)\), and \(E_\infty ^{\bullet \bullet }=\operatorname {Gr}_W^\bullet (\operatorname {H}^\bullet (K))\).
    The filtration \(F\) on \(\operatorname {H}^\bullet (K)\) then induces on \(E_\infty ^{pq}\) a new filtration.\begin{proposition}{1.3.13}\setnumber{1.3.13}\label{hodge-theory-ii-1.3.13}\skipsinglepar\begin{enumerate}
    
 \item[(i)]
      For the first direct filtration, the morphisms \(d_r\) are compatible with the filtrations.
      If \(E_{r+1}^{pq}\) is considered as a quotient of a sub-object of \(E_r^{pq}\), then the first direct filtration on \(E_{r+1}^{pq}\) is finer than the filtration \(F'\) induced by the first direct filtration on \(E_r^{pq}\)Y we have \(F_d(E_{r+1}^{pq})\subset  F'(E_{r+1}^{pq})\).
    


    
 \item[(ii)]
      Dually, the morphisms \(d_r\) are compatible with the second direct filtration, and the second direct filtration on \(E_{r+1}^{pq}\) is less fine than the filtration induced by that of \(E_r^{pq}\).
    


    
 \item[(iii)]
      \(F_d(E_r^{pq})\subset  F_r(E_r^{pq})\subset  F_{d^*}(E_r^{pq})\).
    


    
 \item[(iv)]
      On \(E_\infty ^{pq}\), the filtration induced by the filtration \(F\) of \(\operatorname {H}^\bullet (K)\) \cref{hodge-theory-ii-1.3.12} is finer than the first direct filtration and less fine than the second.
    

  \end{enumerate}\begin{proof}\par{}Claim (i) is evident, (ii) is its dual, and (iii) follows by induction.
      The first claim of (iv) is easy to verify, and the second is its dual.\end{proof}\end{proposition}\label{hodge-theory-ii-1.3.14}\par{}We denote by \(\operatorname {Dec}(K)\) (resp. \(\operatorname {Dec}^\bullet (K)\)) the complex \(K\) endowed with the filtrations \(\operatorname {Dec}(W)\) and \(F\) (resp. \(\operatorname {Dec}^\bullet (W)\) and \(F\)).\par{}It is clear by \hyperref[hodge-theory-ii-1.3.4.1]{Equation~\ref*{hodge-theory-ii-1.3.4.1}} that the isomorphism in \hyperref[hodge-theory-ii-1.3.4]{Proposition~\ref*{hodge-theory-ii-1.3.4}} sends the first direct filtration on \(E_r(\operatorname {Dec}(K))\) to the second direct filtration on \(E_{r+1}(K)\) (for \(r\geqslant 1\)).
    The dual isomorphism \cref{hodge-theory-ii-1.3.5} sends the second direct filtration on \(E_r(\operatorname {Dec}^\bullet (K))\) to the second direct filtration on \(E_{r+1}(K)\).\begin{lemma}{1.3.15}\setnumber{1.3.15}\label{hodge-theory-ii-1.3.15}\skipsinglepar\par{}If the filtration \(F\) is biregular, and if, on the \(\operatorname {Gr}_W^p(K)\), the morphisms \(d\) are strictly compatible with the filtration induced by \(F\), then

    
 \begin{enumerate}
      
 \item[(i)]
        The morphism \hyperref[hodge-theory-ii-1.3.3.2]{Equation~\ref*{hodge-theory-ii-1.3.3.2}} of graded complexes filtered by \(F\)
        \[           u\colon  \operatorname {Gr}_{\operatorname {Dec}(W)}(K)           \to  E_1(K,W)         \]
        is a filtered quasi-isomorphism.
      


      
 \item[(ii)]
        Dually, the morphism in \cref{hodge-theory-ii-1.3.5}
        \[           u\colon  E_1(K,W)           \to  \operatorname {Gr}_{\operatorname {Dec}^\bullet (W)}(K)         \]
        is a filtered quasi-isomorphism.
      

    \end{enumerate}\begin{proof}\par{}It suffices, by duality, to prove (i).\par{}By \cref{hodge-theory-ii-1.3.3} and \hyperref[hodge-theory-ii-1.3.4]{Proposition~\ref*{hodge-theory-ii-1.3.4}}, the complex \(E_1(K,W)\) filtered by \(F\) is a quotient of the filtered complex \(\operatorname {Gr}_{\operatorname {Dec}(W)}(K)\).
      Let \(U\) be the filtered complex given by the kernel, which is acyclic by (ii) of \hyperref[hodge-theory-ii-1.3.4]{Proposition~\ref*{hodge-theory-ii-1.3.4}}.
      The long exact sequence in cohomology associated to the exact sequence of complexes
      \[         0         \to  \operatorname {Gr}_F(U)         \to  \operatorname {Gr}_F(\operatorname {Gr}_{\operatorname {Dec}(W)}(K))         \to  \operatorname {Gr}_F(E_1(K,W))         \to  0       \]
      shows that \(u\) is a filtered quasi-isomorphism if and only if \(\operatorname {Gr}_F(U)\) is an acyclic complex.
      By \hyperref[hodge-theory-ii-1.3.2]{Proposition~\ref*{hodge-theory-ii-1.3.2}}, and since \(U\) is acyclic, this reduces to asking that the differentials of \(U\) be strictly compatible with the filtration \(F\).
      From \hyperref[hodge-theory-ii-1.3.3.1]{Equation~\ref*{hodge-theory-ii-1.3.3.1}} we obtain that \(U\) is the sum over \(p\) of the complexes
      \[         (U^p)^n         = B_1^{p+n,-p}/Z_1^{p+1+n,-p-1}       \]
      endowed with the filtration induced by \(F\).\par{}Each differential \(d\) of each of the complexes \(U^p\) fits into a commutative diagram of filtered objects of the following type, where, for simplicity, we have omitted the total or complementary degree:

      
  \begin {tikzcd}
        B^p/Z^{p+1}
          \ar [rrr,"d"]
          \ar [dr]
        &&& B^{p+1}/Z^{p+1}
          \ar [dd]
      \\& \operatorname {Coim}(d)
          \ar [r,"(5)"]
        & \operatorname {Im}(d)
          \ar [ur] \ar [dr]
      \\W^{p+1}/Z^{p+1}
          \ar [uu]
          \ar [ur]
          \ar [rrr,near start,swap,"(3)"]
        &&& B^{p+1}/W^{p+2}
          \ar [d]
      \\W^{p+1}/W^{p+2}
          \ar [u]
          \ar [rrr,near start,swap,"(1)"]
          \ar [urrr,phantom,"(2)"description]
          \ar [uuurrr,phantom,"(4)"description]
        &&& W^{p+1}/W^{p+2}
      \end {tikzcd}


      By hypothesis, the morphism \((1)\) is strict.
      Since the square \((2)\) is exactly the canonical decomposition of \((1)\), the arrow \((3)\) is a filtered isomorphism.
      The arrows of the trapezium \((4)\) are isomorphisms;
      they are thus filtered isomorphisms, since \((3)\) is a filtered isomorphism.
      The fact that \((5)\) is a filtered isomorphism implies that \(d\) is strict.
      This proves the lemma.\end{proof}\end{lemma}\begin{theorem}{1.3.16}\setnumber{1.3.16}\label{hodge-theory-ii-1.3.16}\skipsinglepar\par{}Let \(K\) be a complex endowed with two filtrations, \(W\) and \(F\), with the filtration \(F\) biregular.
    Let \(r_0\geqslant 0\) be an integer, and suppose that, for \(0\leqslant  r<r_0\), the differentials of the graded complex \(E_r(K,W)\) are strictly compatible with the filtration \(F\).
    Then, for \(r\leqslant  r_0+1\), \(F_d=F_r=F_{d^*}\) on \(E_r^{pq}\).\begin{proof}\par{}We will prove the theorem by induction on \(r_0\).
      For \(r_0=0\), the hypothesis is empty, and we apply \hyperref[hodge-theory-ii-1.3.10]{Lemma~\ref*{hodge-theory-ii-1.3.10}} and (iii) of \hyperref[hodge-theory-ii-1.3.13]{Proposition~\ref*{hodge-theory-ii-1.3.13}}.
      For \(r_0\geqslant 1\), by the inductive hypothesis, we have \(F_d=F_r=F_{d^*}\) on \(E_r^{pq}\) for \(r\leqslant  r_0\).\par{}By \hyperref[hodge-theory-ii-1.3.15]{Lemma~\ref*{hodge-theory-ii-1.3.15}}, the morphism \(u\colon  E_0(\operatorname {Dec}(K))\to  E_1(K)\) is a filtered quasi-isomorphism.
      It thus induces a filtered isomorphism from \(\operatorname {H}^\bullet (\operatorname {Dec}(K))\) to \(\operatorname {H}^\bullet (E_1(K))\):
      \[         u\colon  (E_1(\operatorname {Dec}(K)),F_r)         \xrightarrow {\sim } (E_2(K),F_r).       \]
      Step-by-step, we thus deduce that the canonical isomorphism from \(E_s(\operatorname {Dec}(K))\) to \(E_{s+1}\) (for \(s\geqslant 1\)) is a filtered isomorphism for the recurrent filtration.\par{}On \(E_1(\operatorname {Dec}(K))\), \(F_r=F_d\) (\hyperref[hodge-theory-ii-1.3.10]{Lemma~\ref*{hodge-theory-ii-1.3.10}}), and we already know (\cref{hodge-theory-ii-1.3.14}) that \(u'\) is a filtered isomorphism
      \[         u'\colon  (E_1(\operatorname {Dec}(K)),F_d)         \xrightarrow {\sim } (E_2(K),F_d).       \]\par{}On \(E_2(K)\), we thus have \(F_d=F_r\).
      This, combined with the dual result, proves the theorem for \(r_0=1\).\par{}Suppose that \(r_0\geqslant 2\).
      Then the arrows \(d_1\) of \(E_1(K)\) are strictly compatible with the filtrations, and thus so too are the arrows \(d_0\) of \(E_0(\operatorname {Dec}(K))\) (indeed, \(u\) induces an isomorphism of spectral sequences, and we apply the criterion from \hyperref[hodge-theory-ii-1.3.2]{Proposition~\ref*{hodge-theory-ii-1.3.2}}).\par{}For \(0<s<r_0-1\), the isomorphism \((E_s(\operatorname {Dec}(K)),F_r)\cong (E_{s+1}(K),F_r)\) shows that the \(d_s\) are strictly compatible with the recurrent filtrations.\par{}By the induction hypothesis, we thus have \(F_d=F_r\) on \(E_s(\operatorname {Dec}(K))\) for \(s\leqslant  s_0\).
      The isomorphism \((E_s(\operatorname {Dec}(K)),F_d)\cong (E_{s+1}(K),F_d)\) (\hyperref[hodge-theory-ii-1.3.13]{Proposition~\ref*{hodge-theory-ii-1.3.13}}) then shows that \(F_d=F_r\) on \(E_r(K)\) for \(r\leqslant  r_0+1\).
      This, combined with the dual result, proves the theorem.\end{proof}\end{theorem}\begin{corollary}{1.3.17}\setnumber{1.3.17}\label{hodge-theory-ii-1.3.17}\skipsinglepar\par{}Under the general hypotheses of \hyperref[hodge-theory-ii-1.3.16]{Theorem~\ref*{hodge-theory-ii-1.3.16}}, suppose that, for all \(r\), the differentials \(d_r\) are strictly compatible with the recurrent filtrations on the \(E_r\).
    Then, on \(E_\infty \), the filtrations \(F_d\), \(F_r\), and \(F_{d^*}\) agree, and coincide with the filtration induced by the filtration \(F\) of \(\operatorname {H}^\bullet (K)\).\begin{proof}\par{}This follows immediately from \hyperref[hodge-theory-ii-1.3.16]{Theorem~\ref*{hodge-theory-ii-1.3.16}} and (iv) of \hyperref[hodge-theory-ii-1.3.13]{Proposition~\ref*{hodge-theory-ii-1.3.13}}.\end{proof}\end{corollary}\section{Hypercohomology of filtered complexes}\label{hodge-theory-ii-1.4}\par{}In this section, we recall some standard constructions in hypercohomology.
  We do not use the language of derived categories, which would be more natural here.\par{}\emph{Throughout this entire section, by "complex" we mean "bounded-below complex".}\label{hodge-theory-ii-1.4.1}\par{}Let \(T\) be a left-exact functor from an abelian category \(\mathscr {A}\) to an abelian category \(\mathscr {B}\).
    Suppose that every object of \(\mathscr {A}\) injects into an injective object;
    the derived functors \(\mathrm {R}^iT\colon \mathscr {A}\to \mathscr {B}\) are then defined.
    An object \(A\) of \(\mathscr {A}\) is said to be \emph{acyclic} for \(T\) if \(\mathrm {R}^iT(A)=0\) for \(i>0\).\label{hodge-theory-ii-1.4.2}\par{}Let \((A,F)\) be a filtered object with finite filtration, and \(TF\) the filtration of \(TA\) by its sub-objects \(TF^p(A)\) (these are sub-objects since \(T\) is left exact).
    If \(\operatorname {Gr}_F(A)\) is \(T\)-acyclic, then the \(F^p(A)\) are \(T\)-acyclic as successive extensions of \(T\)-acyclic objects.
    The image under \(T\) of the sequence
    \[       0       \to  F^{p+1}(A)       \to  F^p(A)       \to  \operatorname {Gr}^p(A)       \to  0     \]
    
    is thus exact, and

    \begin{eqn}{1.4.2.1}\setnumber{1.4.2.1}\label{hodge-theory-ii-1.4.2.1}\skipsinglepar\[         \operatorname {Gr}_{FT}TA \xrightarrow {\sim } T\operatorname {Gr}_FA.       \tag{1.4.2.1}       \]\end{eqn}\label{hodge-theory-ii-1.4.3}\par{}Let \(A\) be an object endowed with finite filtrations \(F\) and \(W\) such that \(\operatorname {Gr}_F\operatorname {Gr}_W A\) are \(T\)-acyclic.
    The objects \(\operatorname {Gr}_FA\) and \(\operatorname {Gr}_WA\) are then \(T\)-acyclic, as well as the \(F^q(A)\cap  W^p(A)\).
    The sequences
    \[       0       \to  T(F^q\cap  W^{p+1})       \to  T(F^q\cap  W^p)       \to  T((F^q\cap  W^p)/(F^q\cap  W^{p+1}))       \to  0     \]
    are thus exact, and \(T(F^q(\operatorname {Gr}_W^p(A)))\) is the image in \(T(\operatorname {Gr}_W^p(A))\) of \(T(F^p\cap  W^q)\).
    The diagram
    \[       \begin {CD}         T(F^q\cap  W^p) @>>> T(F^q\operatorname {Gr}_W^pA) @>>> T\operatorname {Gr}_W^pA       \\@V{\cong }VV @. @VV{\cong }V       \\TF^q\cap  TW^p @= TF^q\cap  TW^p @>>> \operatorname {Gr}_{TW}^pTA       \end {CD}     \]
    then shows that the isomorphism in \hyperref[hodge-theory-ii-1.4.2.1]{Equation~\ref*{hodge-theory-ii-1.4.2.1}} relative to \(W\) sends the filtration \(\operatorname {Gr}_{TW}(TF)\) to the filtration \(T(\operatorname {Gr}_W(F))\).\label{hodge-theory-ii-1.4.4}\par{}Let \(K\) be a complex of objects of \(\mathscr {A}\).
    The \emph{hypercohomology objects} \(\mathrm {R}^iT(K)\) are calculated as follows:

    
 \begin{enumerate}
      
 \item[(a)]
        We choose a quasi-isomorphism \(i\colon  K\to  K\) such that the components of \(K'\) are acyclic for \(T\).
        For example, we can take \(K'\) to be the simple complex associated to an injective Cartan–Eilenberg resolution of \(K\).
      

      
 \item[(b)]
        We set
        \[           \mathrm {R}^iT(K)           = \operatorname {H}^i(T(K')).         \]
      

    \end{enumerate}


    We can show that \(\mathrm {R}^iT(K)\) does not depend on the choice of \(K'\), but depends functorially on \(K\), and that a quasi-isomorphism \(f\colon  K_1\to  K_2\) induces \emph{isomorphisms}
    \[       \mathrm {R}^iT(f)\colon        \mathrm {R}^iT(K_1)       \to  \mathrm {R}^iT(K_2).     \]\label{hodge-theory-ii-1.4.5}\par{}Let \(F\) be a biregular filtration of \(K\).
    A \emph{\(T\)-acyclic filtered resolution} of \(K\) is a filtered quasi-isomorphism \(i\colon  K\to  K'\) from \(K\) to a filtered biregular complex such that the \(\operatorname {Gr}^p({K'}^n)\) are acyclic for \(T\).
    If \(K'\) is such a resolution, then the \({K'}^n\) are acyclic for \(T\), and the filtered complex (cf. \cref{hodge-theory-ii-1.4.2}) \(T(K')\) defines a spectral sequence
    \[       E_1^{pq}       = \mathrm {R}^{p+q}T(\operatorname {Gr}^p(K))       \Rightarrow  \mathrm {R}^{p+q}T(K).     \]
    This is independent of the choice of \(K'\).
    We call this the \emph{hypercohomology spectral sequence of the filtered complex \(K\)}.
    It depends functorially on \(K\), and a filtered quasi-isomorphism induces an isomorphism of spectral sequences.\par{}The differentials \(d_1\) of this spectral sequence are the connection morphisms defined by the short exact sequences
    \[       0       \to  \operatorname {Gr}^{p+1}K       \to  F^pK/F^{p+2}K       \to  \operatorname {Gr}^pK       \to  0.     \]\label{hodge-theory-ii-1.4.6}\par{}Let \(K\) be a complex.
    We denote by \(\tau _{\leqslant  p}(K)\) the following subcomplex:
    \[       \tau _{\leqslant  p}(K)^n       = \begin {cases}         K^n         &\text {for }n<p       \\\operatorname {Ker}(d)         &\text {for }n=p       \\0         &\text {for }n>p.       \end {cases}     \]\par{}The \emph{filtration}, said to be \emph{canonical}, of \(K\) by the \(\tau _{\leqslant  p}(K)\) is induced by shifting the trivial filtration \(G\) for which \(G^0(K)=K\) and \(G^1(K)=0\).
    We have, for the canonical filtration,
    \[       \begin {aligned}         E_1^{pq} = 0         &\qquad \text {if }p+q\neq -p       \\\\H^{-p}         &\qquad \text {if }p+q=-p.       \end {aligned}     \]
    A quasi-isomorphism \(f\colon  K\to  K'\) is automatically a filtered quasi-isomorphism for the canonical filtrations.\label{hodge-theory-ii-1.4.7}\par{}The subcomplexes \(\sigma _{\geqslant  p}(K)\) of \(K\)
    \[       \sigma _{\geqslant  p}(K)^n       = \begin {cases}         0         &\text {if }n<p       \\K^n         &\text {if }n\geqslant  p       \end {cases}     \]
    define a biregular filtration, called the \emph{stupid filtration} of \(K\).\par{}The hypercohomology spectral sequences attached to the stupid or canonical filtrations of \(K\) are the two \emph{hypercohomology spectral sequences} of \(K\).\begin{example}{1.4.8}\setnumber{1.4.8}\label{hodge-theory-ii-1.4.8}\skipsinglepar\par{}Let \(f\colon  X\to  Y\) be a continuous map between topological spaces, and let \(\mathscr {F}\) be an abelian sheaf on \(X\).
    Let \(\mathscr {F}^\bullet \) be a resolution of \(\mathscr {F}\) by \(f_*\)-acyclic sheaves.
    We have \(\mathrm {R}^i f_*\mathscr {F}\simeq \operatorname {\mathscr {H}}^i(f_*\mathscr {F}^\bullet )\).
    We take the functor \(T\) to be the functor \(\Gamma (Y,-)\).
    The hypercohomology spectral sequence of the complex \(f_*\mathscr {F}^\bullet \) endowed with its canonical filtration \cref{hodge-theory-ii-1.4.6}
    \[       E_1^{pq}       = \operatorname {H}^{2p+q}(Y,\mathrm {R}^{-p}f_*\mathscr {F})       \Rightarrow  \operatorname {H}^{p+q}(X,\mathscr {F})     \]
    is exactly, up to the renumbering \(E_r^{pq}\mapsto  E_{r+1}^{2p+q,-p}\), the \emph{Leray spectral sequence} for \(f\) and \(\mathscr {F}\).\end{example}\label{hodge-theory-ii-1.4.9}\par{}Let \((K,W,F)\) be a biregular bifiltered complex.
    To this complex, we associate:

    
 \begin{enumerate}
      
 \item[(a)]
        A spectral sequence
        \[           {}_W E_1^{p,n-p}           = \operatorname {H}^n(\operatorname {Gr}_W^p(K)) \Rightarrow  \operatorname {H}^n(K)         \]
        with differentials \({}_W d_1\) being the connecting morphisms induced by the short exact sequences
        \[           0           \to  \operatorname {Gr}_W^{p+1}(K)           \to  W^p(K)/W^{p+2}(K)           \to  \operatorname {Gr}_W^p(K)           \to  0;         \]
      


      
 \item[(b)]
        
        An analogous spectral sequence for the filtration \(F\);
      


      
 \item[(c)]
        Exact squares

        
  \begin {tikzcd}
          & 0 \dar 
          & 0 \dar 
          & 0 \dar 
        \\0 \rar 
          & \operatorname {Gr}_F^{p+1}\operatorname {Gr}_W^{q+1}K \rar  \dar 
          & F^p/F^{p+2}(\operatorname {Gr}_W^{q+1}K) \rar  \dar 
          & \operatorname {Gr}_F^p\operatorname {Gr}_W^{q+1}K \rar  \dar 
          & 0
        \\0 \rar 
          & \operatorname {Gr}_F^{p+1}(W^q/W^{q+2}(K)) \rar  \dar 
          & F^p/F^{p+2}(W^q/W^{q+2}(K)) \rar  \dar 
          & \operatorname {Gr}_F^p(W^q/W^{q+2}(K)) \rar  \dar 
          & 0
        \\0 \rar 
          & \operatorname {Gr}_F^{p+1}\operatorname {Gr}_W^q K \rar  \dar 
          & F^p/F^{p+2}\operatorname {Gr}_W^q K \rar  \dar 
          & \operatorname {Gr}_F^p\operatorname {Gr}_W^q K \rar  \dar 
          & 0
        \\& 0
          & 0
          & 0
        \end {tikzcd}

      

    \end{enumerate}\par{}The exterior rows and columns of this square define connection morphisms
    \[       \begin {aligned}         {}_{F,W}d_1\colon  \operatorname {H}^n\operatorname {Gr}_F^p\operatorname {Gr}_W^q K         &\to  \operatorname {H}^{n+1}\operatorname {Gr}_F^{p+1}\operatorname {Gr}_W^q K       \\{}_{W,F}d_1\colon  \operatorname {H}^n\operatorname {Gr}_F^p\operatorname {Gr}_W^q K         &\to  \operatorname {H}^{n+1}\operatorname {Gr}_F^p\operatorname {Gr}_W^{q+1} K.       \end {aligned}     \]
    These morphisms satisfy
    \[       \begin {gathered}         {}_{FW}d_1 \circ  {}_{WF}d_1 + {}_{WF}d_1 \circ  {}_{FW}d_1         = 0       \\ {}_{FW}d_1^2         = 0       \\ {}_{WF}d_1^2         = 0       \end {gathered}     \]\par{}The morphisms \({}_{FW}d_1\) are the morphisms \(d_1\) of the spectral sequences \(E^{(q)}\), with the \(E_1^{(q)p,n-p}\) term equal to

    \begin{eqn}{1.4.9.1}\setnumber{1.4.9.1}\label{hodge-theory-ii-1.4.9.1}\skipsinglepar\[         E_1^{p,q,n-p-q}         \coloneqq  \operatorname {H}^n(\operatorname {Gr}_F^p\operatorname {Gr}_W^q K)         \Rightarrow  \operatorname {H}^n(\operatorname {Gr}_W^q K)         = {}_WE_1^{q,n-q}       \tag{1.4.9.1}       \]\end{eqn}

    defined by the filtered complex \(\operatorname {Gr}_W^q(K)\).
    This spectral sequence abuts to the filtration induced by \(F\) on \(\operatorname {H}^\bullet \operatorname {Gr}_W^q K\).
    Similarly, the \({}_{WF}d_1\) are the \(d_1\) of spectral sequences with the same initial terms

    \begin{eqn}{1.4.9.2}\setnumber{1.4.9.2}\label{hodge-theory-ii-1.4.9.2}\skipsinglepar\[         E_1^{p,q,n-p-q}         \coloneqq  \operatorname {H}^n(\operatorname {Gr}_F^p\operatorname {Gr}_W^q K)         \Rightarrow  \operatorname {H}^n(\operatorname {Gr}_F^p K).       \tag{1.4.9.1}       \]\end{eqn}

    
  \begin {tikzcd}
      & \operatorname {H}^n\operatorname {Gr}_W^q K
        \ar [Rightarrow,dr,"{}_Wd_1"]
    \\\operatorname {H}^n\operatorname {Gr}_F^p\operatorname {Gr}_W^q K
        \ar [Rightarrow,ur,"{}_{F,W}d_1"]
        \ar [Rightarrow,dr,swap,"{}_{W,F}d_1"]
      && \operatorname {H}^n K
    \\& \operatorname {H}^n\operatorname {Gr}_F^p K
        \ar [Rightarrow,ur,swap,"{}_Fd_1"]
    \end {tikzcd}


    
    These constructions are symmetric in \(F\) and \(W\) via the isomorphism
    \[       \operatorname {Gr}_F^p\operatorname {Gr}_W^q       \sim  \operatorname {Gr}_W^q\operatorname {Gr}_F^p.     \]\label{hodge-theory-ii-1.4.10}\par{}We can also interpret the \({}_{W,F}d_1\) are the initial morphisms of a morphism of spectral sequences \hyperref[hodge-theory-ii-1.4.9.2]{Equation~\ref*{hodge-theory-ii-1.4.9.2}}, abutting to \({}_Wd_1\).\par{}Indeed, let \(C^q\) be the cone of the morphism
    \[       W^q(K)/W^{q+2}(K) \to  \operatorname {Gr}_W^q(K).     \]
    In the diagram
    \[       \Sigma \colon  \operatorname {Gr}_W^{q+1}(K)[1]       \xrightarrow {u} C^q       \xleftarrow {i} \operatorname {Gr}_W^q(K)     \]
    the morphism \(u\) is a quasi-isomorphism, and we have
    \[       {}_Wd_1       = \operatorname {H}(u)^{-1}\circ \operatorname {H}(i).     \]\par{}In fact, \(u\) is even a filtered (for \(F\)) quasi-isomorphism, and the above construction defines a morphism from the spectral sequence defined by \((\operatorname {Gr}_W^q(K),F)\) to that defined by \((\operatorname {Gr}_W^{q+1}(K)[1],F)\), and it abuts to \({}_Wd_1\).
    The initial term of this morphism, induced by \(\operatorname {Gr}_F(\Sigma )\) is exactly \({}_{W,F}d_1\).\label{hodge-theory-ii-1.4.11}\par{}These constructions pass as they are to hypercohomology.
    Let \(K\) be a complex endowed with two biregular filtrations \(F\) and \(W\).\par{}A \emph{bifiltered \(T\)-acyclic resolution} of \(K\) is a bifiltered quasi-isomorphism \(i\colon  K\to  K'\) such that the \(\operatorname {Gr}_F^p\operatorname {Gr}_W^n({K'}^m)\) are \(T\)-acyclic.
    Such a morphism always exists.
    In the particular case where \(\mathscr {A}\) is the category of sheaves of \(A\)-modules on a topological space \(X\), and where \(T\) is the functor \(\Gamma \) from \(\mathscr {A}\) to the category of \(A\)-modules, then an example of a bifiltered \(T\)-acyclic resolution of \(K\) is the simple complex associated to the double complex given by the Godement resolution \(\mathscr {C}^\bullet (K)\) of \(K\), filtered by the \(\mathscr {C}^\bullet (F^p(K))\) and the \(\mathscr {C}^\bullet (W^n(K))\).
    Since \(\mathscr {C}^\bullet \) is exact, we have
    \[       \operatorname {Gr}_F\operatorname {Gr}_W(\mathscr {C}^\bullet (K))       \simeq  \mathscr {C}^\bullet (\operatorname {Gr}_F\operatorname {Gr}_W(K)).     \]
    We will not have need here of any other case.\par{}If \(K'\) is a bifiltered \(T\)-acyclic resolution of \(K\), then the complex \(TK'\) is filtered by the \(TF^pK'\) and the \(TW^qK'\) (\cref{hodge-theory-ii-1.4.3}).
    Furthermore, \(\operatorname {Gr}_W^n(K')\) is a \(T\)-acyclic filtered (for \(F\)) resolution of \(\operatorname {Gr}_F^n(K)\), and \(\operatorname {Gr}_F^n\operatorname {Gr}_W^m(K')\) is a \(T\)-acyclic resolution of \(\operatorname {Gr}_F^n\operatorname {Gr}_W^m(K)\), and
    \[       \begin {gathered}         T\operatorname {Gr}_F K'         \approx  \operatorname {Gr}_F TK'         \qquad \text {as }W\text {-filtered complexes}       \\T\operatorname {Gr}_W K'         \approx  \operatorname {Gr}_W TK'         \qquad \text {as }F\text {-filtered complexes}       \\T\operatorname {Gr}_F\operatorname {Gr}_W K'         \approx  \operatorname {Gr}_F\operatorname {Gr}_W TK'.       \end {gathered}     \]\begin{lemma}{1.4.12}\setnumber{1.4.12}\label{hodge-theory-ii-1.4.12}\skipsinglepar\par{}Under the hypotheses of \cref{hodge-theory-ii-1.4.11}:

    
 \begin{enumerate}
      
 \item[(i)]
        The initial terms of the hypercohomology spectral sequences
        \[           {}_WE_1^{q,n-q}           = \mathrm {R}^nT(\operatorname {Gr}_W^qK)           \Rightarrow  \mathrm {R}^nT(K)         \tag{1}         \]
        \[           {}_FE_1^{p,n-p}           = \mathrm {R}^nT(\operatorname {Gr}_F^pK)           \Rightarrow  \mathrm {R}^nT(K)         \tag{2}         \]
        
        are abutments of the hypercohomology spectral sequences of the filtered complexes \(\operatorname {Gr}_W^qK\) and \(\operatorname {Gr}_F^pK\), with \(E_1\) pages given by
        \[           E_1^{p,q,n-p-q}           \coloneqq  \mathrm {R}^nT(\operatorname {Gr}_F^p\operatorname {Gr}_W^qK)           \Rightarrow  {}_WE_1^{q,n-q}           \qquad \text {for fixed }q         \tag{3}         \]
        \[           E_1^{p,q,n-p-q}           \coloneqq  \mathrm {R}^nT(\operatorname {Gr}_F^p\operatorname {Gr}_W^qK)           \Rightarrow  {}_FE_1^{p,n-p}           \qquad \text {for fixed }p         \tag{4}         \]
      


      
 \item[(ii)]
        The filtration of \({}_WE_1^{p,n-p}\), abutting to the spectral sequence \((3)\), is the filtration of \({}_WE_1^{p,n-p}(TK')\) induced by the filtration \(F\) of \(TK'\).
      


      
 \item[(iii)]
        For the differentials of the complexes \(\operatorname {Gr}_W^n(T(K'))\) to be strictly compatible with the filtration \(F\), it is necessary and sufficient that the hypercohomology spectral sequences \((3)\) degenerate on the \(E_1\) page.
      


      
 \item[(iv)]
        The morphisms \(d_1\) of the spectral sequence \((4)\) are the initial terms of the degree-1 morphisms of the spectral sequences \((3)\) that abut to the morphisms \(d_1\) of the spectral sequence \((1)\).
      

    \end{enumerate}\begin{proof}\par{}Claims (i) and (iv) follow from \cref{hodge-theory-ii-1.4.9} and \cref{hodge-theory-ii-1.4.10} applied to \(TK'\) via the isomorphisms \cref{hodge-theory-ii-1.4.11}.
      Claim (ii) is then trivial by the definition of the recurrent filtration \(F\) (identical to the discrete filtrations by \hyperref[hodge-theory-ii-1.3.10]{Lemma~\ref*{hodge-theory-ii-1.3.10}} and (iii) of \hyperref[hodge-theory-ii-1.3.13]{Proposition~\ref*{hodge-theory-ii-1.3.13}}), and claim (iii) follows from \hyperref[hodge-theory-ii-1.3.2]{Proposition~\ref*{hodge-theory-ii-1.3.2}}.\end{proof}\end{lemma}\chapter{Hodge structures}\thispagestyle{fancy}\label{hodge-theory-ii-2}\section{Pure structures}\label{hodge-theory-ii-2.1}\label{hodge-theory-ii-2.1.1}\par{}In all the following, we denote by \(\mathbb {C}\) an algebraic closure of \(\mathbb {R}\), and we do not suppose to have chosen a root \(i\) of the equation \(x^2+1=0\).
    The theory will be invariant under complex conjugation (cf. \cref{hodge-theory-ii-2.1.14}).\label{hodge-theory-ii-2.1.2}\par{}We denote by \(S\) the real algebraic group \(\mathbb {C}^*\), given by restriction of scalars à la Weil from \(\mathbb {C}\) to \(\mathbb {R}\) of the group \(\mathbb {G}_\mathrm {m}\):
    \[       \begin {aligned}         S         &= \prod _{\mathbb {C}{/}\mathbb {R}}\mathbb {G}_\mathrm {m}       \\S(\mathbb {R})         &= \mathbb {C}^*.       \end {aligned}     \]\par{}The group \(S\) is a torus, i.e. it is connected and of multiplicative type.
    It is thus described by the free abelian group of finite type
    \[       X(S)       = \operatorname {Hom}(S_\mathbb {C},\mathbb {G}_\mathrm {m})       = \underline {\operatorname {Hom}}(S,\mathbb {G}_\mathrm {m})(\mathbb {C})     \]
    of its complex characters, endowed with the action of \(\operatorname {Gal}(\mathbb {C}{/}\mathbb {R})=\mathbb {Z}{/}(2)\).\par{}The group \(X(S)\) has generators \(z\) and \(\overline {z}\), which induce (respectively) the identity and complex conjugation:
    \[       \mathbb {C}^*       = S(\mathbb {R})       \to  S(\mathbb {C})       \to  \mathbb {G}_\mathrm {m}(\mathbb {C})       = \mathbb {C}^*.     \]
    Complex conjugation exchanges \(z\) and \(\overline {z}\).\label{hodge-theory-ii-2.1.3}\par{}We have a canonical map

    \begin{eqn}{2.1.3.1}\setnumber{2.1.3.1}\label{hodge-theory-ii-2.1.3.1}\skipsinglepar\[         w\colon  \mathbb {G}_\mathrm {m}\to  S       \tag{2.1.3.1}       \]\end{eqn}

    
    that, on the real points, induces the inclusion of \(\mathbb {R}^*\) into \(\mathbb {C}^*\).
    We have

    \begin{eqn}{2.1.3.2}\setnumber{2.1.3.2}\label{hodge-theory-ii-2.1.3.2}\skipsinglepar\[         zw         = \overline {z}w         = \mathrm {Id}.       \tag{2.1.3.2}       \]\end{eqn}\par{}We also have a map

    \begin{eqn}{2.1.3.3}\setnumber{2.1.3.3}\label{hodge-theory-ii-2.1.3.3}\skipsinglepar\[         N\colon  S\to \mathbb {G}_\mathrm {m}       \tag{2.1.3.3}       \]\end{eqn}

    that on the real points can be identified with the norm \(N_{\mathbb {C}{/}\mathbb {R}}\colon \mathbb {C}^*\to \mathbb {R}^*\).
    We have

    \begin{eqn}{2.1.3.4}\setnumber{2.1.3.4}\label{hodge-theory-ii-2.1.3.4}\skipsinglepar\[         N = z\overline {z}       \tag{2.1.3.4}       \]\end{eqn}

    \begin{eqn}{2.1.3.5}\setnumber{2.1.3.5}\label{hodge-theory-ii-2.1.3.5}\skipsinglepar\[         N\circ  w         = (x\mapsto  x^2).       \tag{2.1.3.5}       \]\end{eqn}\begin{definition}{2.1.4}\setnumber{2.1.4}\label{hodge-theory-ii-2.1.4}\skipsinglepar\par{}A \emph{real Hodge structure} is a real vector space \(V\) of finite dimension endowed with an action of the real algebraic group \(S\).\end{definition}\label{hodge-theory-ii-2.1.5}\par{}By the general theory of groups of multiplicative type, giving a real Hodge structure on \(V\) is equivalent to giving a bigrading \(V^{p,q}\) of \(V_\mathbb {C}=\mathbb {C}\otimes _\mathbb {R} V\) that satisfies \(\overline {V^{pq}}=V^{qp}\).
    The action of \(S\) and the bigrading are mutually determined via the condition:

    \label{hodge-theory-ii-2.1.5.1}\par{}On \(V^{pq}\), \(S\) acts by multiplication by \(z^p\overline {z}^q\).\label{hodge-theory-ii-2.1.6}\par{}Let \((\mathbb {C},\times )\) be the multiplicative monoid, and let
    \[       \overline {S}       = \prod _{\mathbb {C}{/}\mathbb {R}} (\mathbb {C},\times ).     \]
    If \(V\) is a real vector space, then we can show that it is equivalent to either give an action of \(\overline {S}\) on \(V\) or to give a bigrading \(V^{pq}\) on \(V\) such that \(\overline {V^{pq}}=V^{qp}\) and \(V^{pq}=0\) for \(p<0\) or \(q<0\).\label{hodge-theory-ii-2.1.7}\par{}Let \(V\) be a real Hodge structure, defined by a representation \(\sigma \) of \(S\) and a bigrading \(V^{pq}\).
    The grading of \(V_\mathbb {C}\) by the \(V_\mathbb {C}^n=\sum _{p+q=n}V^{pq}\) is then defined over \(\mathbb {R}\).
    We call this the \emph{weight grading}.
    On \(V^n=V\cap  V_\mathbb {C}^n\), the representation \(\sigma  w\) of \(\mathbb {G}_\mathrm {m}\) is multiplication by \(x^n\).\par{}We say that \(V\) is \emph{of weight \(n\)} if \(V^{pq}=0\) for \(p+q\neq  n\), i.e. if \(\sigma  w\) is multiplication by \(x^n\).\label{hodge-theory-ii-2.1.8}\par{}Let \(V\) be a real Hodge structure.
    The \emph{Hodge filtration} on \(V_\mathbb {C}\) is defined by
    \[       F^p(V_\mathbb {C})       = \sum _{p'\geqslant  p} V^{p'q'}.     \]\par{}By \cref{hodge-theory-ii-1.2.6}, we have:\begin{proposition}{2.1.9}\setnumber{2.1.9}\label{hodge-theory-ii-2.1.9}\skipsinglepar\par{}Let \(n\) be an integer.
    The construction \cref{hodge-theory-ii-2.1.8} establishes an equivalence of categories between:

    
 \begin{enumerate}
      
 \item[(a)]
        the category of real Hodge structures of weight \(n\);
      

      
 \item[(b)]
        the category of pairs consisting of a real vector space \(V\) of finite dimension and of a filtration \(F\) on \(V_\mathbb {C}=\mathbb {C}\otimes _\mathbb {R} V\) that is \(n\)-opposite to its complex conjugate \(\overline {F}\).
      

    \end{enumerate}\end{proposition}\begin{definition}{2.1.10}\setnumber{2.1.10}\label{hodge-theory-ii-2.1.10}\skipsinglepar\par{}A \emph{Hodge structure} \(H\), of weight \(n\), consists of

    
 \begin{enumerate}
      
 \item[(a)]
        a \(\mathbb {Z}\)-module \(H_\mathbb {Z}\) of finite type (the "\emph{integral lattice}");
      

      
 \item[(b)]
        a real Hodge structure of weight \(n\) on \(H_\mathbb {R}=\mathbb {R}\otimes _\mathbb {Z} H_\mathbb {Z}\).
      

    \end{enumerate}\end{definition}\label{hodge-theory-ii-2.1.11}\par{}A morphism \(f\colon  H\to  H'\) is a homomorphism \(f\colon  H_\mathbb {Z}\to  H'_\mathbb {Z}\) such that \(f_\mathbb {R}\colon  H_\mathbb {R}\to  H'_\mathbb {R}\) is compatible with the action of \(S\) (i.e. such that \(f_\mathbb {C}\) is compatible with the bigrading, or with the Hodge filtration).\par{}Hodge structures of weight \(n\) form an abelian category.
    If \(H\) is of weight \(n\), and \(H'\) is of weight \(n'\), then we define a Hodge structure \(H\otimes  H'\) of weight \(n+n'\) by the equations:

    
 \begin{enumerate}
      
 \item[(a)]
        \((H\otimes  H')_\mathbb {Z}=H_\mathbb {Z}\otimes  H'_\mathbb {Z}\);
      

      
 \item[(b)]
        the action of \(S\) on \((H\otimes  H')_\mathbb {R}=H_\mathbb {R}\otimes  H'_\mathbb {R}\) is the tensor product of the actions of \(S\) on \(H_\mathbb {R}\) and on \(H'_\mathbb {R}\).
      

    \end{enumerate}\par{}The bigrading (resp. Hodge filtration) of \((H\otimes  H')_\mathbb {C}=H_\mathbb {C}\otimes  H'_\mathbb {C}\) is the tensor product of the bigradings (resp. Hodge filtrations (cf. \cref{hodge-theory-ii-1.1.12})) of \(H_\mathbb {C}\) and of \(H'_\mathbb {C}\).\par{}We define in an analogous manner the Hodge structure \(\underline {\operatorname {Hom}}(H,H')\) (of weight \(n'-n\)), the Hodge structures \(\wedge ^p H\) (of weight \(pn\)), and the dual Hodge structure \(H^*\).\par{}The \emph{internal hom} \(\underline {\operatorname {Hom}}\) above and the homomorphism group, are related by:

    \begin{remark}{2.1.11.1}\setnumber{2.1.11.1}\label{hodge-theory-ii-2.1.11.1}\skipsinglepar\par{}\(\operatorname {Hom}(H,H')\) is the subgroup of
        \[           \underline {\operatorname {Hom}}(H,H')_\mathbb {Z}           = \operatorname {Hom}_\mathbb {Z}(H_\mathbb {Z},H'_\mathbb {Z})         \]
        consisting of elements of degree \((0,0)\).\end{remark}\par{}The actions of \(S\) on \(H_\mathbb {R}\), \(H'_\mathbb {R}\), and \(\underline {\operatorname {Hom}}(H,H')_\mathbb {R}=\operatorname {Hom}(H_\mathbb {R},H'_\mathbb {R})\) are related by
    \[       s(f(x))       = s(f)(s(x)).     \]
    This means that if \(f\) is of degree \((0,0)\), i.e. invariant under \(S\), then it commutes with the action of \(S\).\label{hodge-theory-ii-2.1.12}\par{}If \(A\) is a Noetherian subring of \(\mathbb {R}\), then we define a \emph{Hodge \(A\)-structure of weight \(n\)} to consist of an \(A\)-module \(H_A\) of finite type along with a real Hodge structure of weight \(n\) on \(H_\mathbb {R}=H_A\otimes _A\mathbb {R}\).
    This definition is mostly used for \(A=\mathbb {Q}\).
    A Hodge \(A\)-structure consists of an \(A\)-module \(H_A\) of finite type along with a real Hodge structure on \(H_\mathbb {R}=H_A\otimes _A\mathbb {R}\) such that the weight grading is defined over the field of fractions of \(A\).\begin{definition}{2.1.13}\setnumber{2.1.13}\label{hodge-theory-ii-2.1.13}\skipsinglepar\par{}The \emph{Tate Hodge structure \(\mathbb {Z}(1)\)} is the Hodge structure of weight \(-2\), of rank \(1\), of pure bidegree \((-1,-1)\), with integral lattice \(2\pi  i\mathbb {Z}\subset \mathbb {C}\).\par{}The action of \(S\) is thus given by multiplication with the inverse of the norm (\hyperref[hodge-theory-ii-2.1.3.3]{Equation~\ref*{hodge-theory-ii-2.1.3.3}}).\par{}For \(n\in \mathbb {Z}\), we define \(\mathbb {Z}(n)\) as the \(n\)-th tensor power of \(\mathbb {Z}(1)\), so \(\mathbb {Z}(n)\) is the Hodge structure of weight \(-2n\), of rank \(1\), of pure bidegree \((-n,-n)\), with integral lattice \((2\pi  i)^n\mathbb {Z}\subset \mathbb {C}\).
    THe action of \(S\) is multiplication by \(N(x)^{-n}\).\end{definition}\label{hodge-theory-ii-2.1.14}\par{}The choice in \(\mathbb {C}\) of a solution \(i\) of the equation \(x^2+1=0\) determines, on each complex variety \(X\) of pure dimension \(n\), an orientation \(\operatorname {or}_i(X)\).
    Replace \(i\) by \(-i\) gives
    \[       \operatorname {or}_{-i}(X)       = (-1)^n\operatorname {or}_i(X).     \]\par{}The choice of \(i\) also defines an element \(C\) of order 4 in \(S(\mathbb {R})\), given by the image of \(i\) under the isomorphism \(S(\mathbb {R})\simeq \mathbb {C}^*\).\par{}Finally, it also defines an isomorphism between \(\mathbb {Z}\) and the integral lattice of \(\mathbb {Z}(n)\), given by multiplication by \((2\pi  i)^n\).\par{}When \(i\), an orientation of \(X\), \(C\), or an identification \(\mathbb {Z}\sim \mathbb {Z}(n)_\mathbb {Z}\) appear in an equation, it is implicitly understood that they all follow from a single choice of the same \(i\), and that by replacing \(i\) with \(-i\) we obtain an equivalent equation.\begin{definition}{2.1.15}\setnumber{2.1.15}\label{hodge-theory-ii-2.1.15}\skipsinglepar\par{}A \emph{polarisation} of a Hodge structure \(H\) of weight \(n\) is a homomorphism
    \[       (x,y)\colon  H\otimes  H \to  \mathbb {Z}(-n)     \]
    such that the real bilinear form \((2\pi  i)^n(x,Cy)\) on \(H_\mathbb {R}\) is symmetric and positive definite.\end{definition}\label{hodge-theory-ii-2.1.16}\par{}The \emph{real Tate Hodge structure} is the real Hodge structure \(\mathbb {R}(1)\) underlying \(\mathbb {Z}(1)\).
    We similarly define \(\mathbb {R}(n)\) as underlying \(\mathbb {Z}(1)\).
    A polarisation of a \emph{real} Hodge structure of weight \(n\) is a homomorphism
    \[       (x,y)\colon  H\otimes  H\to  \mathbb {R}(-n)     \]
    such that the real bilinear form \((2\pi  i)^n(x,Cy)\) on \(H_\mathbb {R}\) is symmetric and positive definite.
    A polarisation is entirely defined by the positive definite quadratic form \((2\pi  i)^n(x,Cy)\) on \(H_\mathbb {R}\), imposed with only the condition of being invariant under the compact sub-torus of \(S\) given by the kernel of \(N\).\par{}We have
    \[       (x,y)       = (Cx,Cy)       = (y,C^2x)       = (-1)^n(y,x).     \]
    The form \((x,y)\) is thus symmetric or alternating, depending on the parity of \(n\).\label{hodge-theory-ii-2.1.17}\par{}The reader can generalise these definitions to Hodge \(A\)-structures of weight \(n\) (\cref{hodge-theory-ii-2.1.12}).\section{Hodge theory}\label{hodge-theory-ii-2.2}\label{hodge-theory-ii-2.2.1}\par{}Let \(X\) be a compact Kähler variety (for example, smooth and projective).
    By the holomorphic Poincaré lemma, the de Rham complex \(\Omega _X^\bullet \) is a resolution of the constant sheaf \(\mathbb {C}\).
    We thus have an isomorphism (\cref{hodge-theory-ii-1.4.2})
    \[       \operatorname {H}^\bullet (X,\mathbb {C})       \sim  \operatorname {\mathbb {H}}^\bullet (X,\Omega _X^\bullet )     \]
    and the stupid filtration on \(\Omega _X^\bullet \) (\cref{hodge-theory-ii-1.4.5}) defines the hypercohomology spectral sequence

    \begin{eqn}{2.2.1.1}\setnumber{2.2.1.1}\label{hodge-theory-ii-2.2.1.1}\skipsinglepar\[         E_1^{pq}         = \operatorname {H}^q(X,\Omega _X^q)         \Rightarrow  \operatorname {H}^{p+q}(X,\mathbb {C})       \tag{2.2.1.1}       \]\end{eqn}

    which abuts to the \emph{Hodge filtration} of \(\operatorname {H}^\bullet (X,\mathbb {C})\).\par{}By Hodge theory [\hyperref[H1952]{H1952}; \hyperref[W1958]{W1958}], we have:

    
 \begin{enumerate}
      
 \item[(A)]
        The spectral sequence \hyperref[hodge-theory-ii-2.2.1.1]{Equation~\ref*{hodge-theory-ii-2.2.1.1}} degenerates: \(E_1=E_\infty \).
      

      
 \item[(B)]
        The Hodge filtration of \(\operatorname {H}^n(X,\mathbb {C})\) is \(n\)-opposite to the complex conjugate filtration.
      

    \end{enumerate}\label{hodge-theory-ii-2.2.2}\par{}Let \(V\) be a local system of real vector spaces over \(X\), i.e. a sheaf of \(\mathbb {R}\)-vector spaces that is locally isomorphic to a constant sheaf \(\mathbb {R}^n\).
    Suppose that there exists on \(V\) a bilinear form
    \[       Q\colon  V\otimes  V\to \mathbb {R}     \]
    that is locally constant and \emph{defined}.
    If \(X\) is connected, this is the case if \(V\) is defined by a representation of a finite quotient of the fundamental group of \(X\).\par{}Statements (A) and (B) \hyperref[hodge-theory-ii-2.2.1]{above} remain true as they are, without needing to change the proofs, for cohomology with coefficients in \(V_\mathbb {C}=V\otimes _\mathbb {R}\mathbb {C}\), since the spectral sequence
    \[       E_1^{pq}       = \operatorname {H}^q(X,\Omega _X^p(V))       \Rightarrow  \operatorname {H}^{p+q}(X,V_\mathbb {C})     \]
    induced by the de Rham resolution of \(V_\mathbb {C}\) by \(\Omega _X^\bullet (V_\mathbb {C})\) degenerates, and abuts to a filtration on \(\operatorname {H}^n(X,V_\mathbb {C})\) that is \(n\)-opposite to the complex conjugate filtration.\par{}The vector space \(\operatorname {H}^n(X,V)\) is thus endowed with a canonical real Hodge structure of weight \(n\).\label{hodge-theory-ii-2.2.3}\par{}We show in [\hyperref[D1968]{D1968}] that the statements in \cref{hodge-theory-ii-2.2.1} remain true when \(X\) is a non-singular complete algebraic, not necessarily Kähler, variety.
    The proof from \emph{loc. cit.}, based on a reduction to the projective case by Chow's lemma and resolution of singularities, extends to the case of \cref{hodge-theory-ii-2.2.2}.\label{hodge-theory-ii-2.2.4}\par{}Let \(\mathcal {L}\) be an invertible sheaf.
    Here are two ways of defining the class \(c_1(\mathcal {L})\).\label{hodge-theory-ii-2.2.4.1}\par{}The sheaf \(\mathcal {L}\) defines an element \(c\) in \(\operatorname {H}^1(X,\mathcal {O}^*)\).
      Its image under \(\mathrm {d} f/f\colon \mathcal {O}^*\to \Omega ^1\) lies in \(\operatorname {H}^1(X,\Omega ^1)\).
      More precisely, \(\mathrm {d} f/f\) defines a morphism of complexes
      \[         \mathrm {d}\log \colon          \mathcal {O}^*[-1]         \to  (0\to \Omega _X^1\to \Omega _X^2\to \ldots )         = \sigma _{\geqslant 1}(\Omega _X^\bullet ).       \]
      This complex maps to \(\Omega _X^\bullet \), whence
      \[         \mathrm {d}\log \colon  \mathcal {O}^*[-1]\to \Omega _X^\bullet        \]
      and the image of \(c\) under \(\mathrm {d}\log \) is in \(\operatorname {\mathbb {H}}^2(\Omega _X^\bullet )\).
      This construction still makes sense for an algebraic variety over an arbitrary field \(k\).
      For \(k=\mathbb {C}\), we further have
      \[         \operatorname {H}^2(X,\mathbb {C})         \xrightarrow {\sim } \operatorname {\mathbb {H}}^2(X,\Omega _X^\bullet )       \]
      whence a class
      \[         c'_1(\mathcal {L}) \in  \operatorname {H}^2(X,\mathbb {C}).       \]\label{hodge-theory-ii-2.2.4.2}\par{}The exponential exact sequence
      \[         0         \to  \mathbb {Z}(1)         \to  \mathcal {O}         \to  \mathcal {O}^*         \to  0       \]
      
      defines a homomorphism
      \[         \partial \colon  \operatorname {H}^1(X,\mathcal {O}^*)\to \operatorname {H}^2(X,\mathbb {Z}(1))       \]
      whence a class
      \[         \partial  c         = c''_1(\mathcal {L}) \in  \operatorname {H}^2(X,\mathbb {Z}(1)).       \]\par{}If we have a choice of \(i\), then this class is identified with
      \[         c'''_1(\mathcal {L}) \in  \operatorname {H}^2(X,\mathbb {Z})       \]
      and for \(\mathcal {L}=\mathcal {O}(D)\), this \(c'''_1(\mathcal {L})\) is exactly the integer cohomology class defined by \(D\) (with the orientations being defined by \(i\)).\label{hodge-theory-ii-2.2.5}\par{}We will prove:

    \label{hodge-theory-ii-2.2.5.1}\par{}For \(\alpha \) the natural injection of \(\mathbb {Z}(1)\) into \(\mathbb {C}\), we have
        \[           -\alpha  c''_1(\mathcal {L})           = c'_1(\mathcal {L}).         \]

    \label{hodge-theory-ii-2.2.5.2}\par{}For \(\beta \) the natural injection of \(\mathbb {Z}\) into \(\mathbb {C}\), we have
        \[           -\beta  c'''_1(\mathcal {L})           = \frac {1}{2\pi  i}c'_1(\mathcal {L}).         \]\par{}We prove this by considering the following diagram of complexes, in which \(\approx \) denotes a quasi-isomorphism, and whose upper triangle anti-commutes up to homotopy:
    \[       \begin {CD}         \mathbb {C} @>\approx >> \Omega _X^\bullet  @<<< \sigma _{\geqslant 1}(\Omega _X^\bullet )       \\@| @. @|       \\\mathbb {C} @<<< (\mathbb {C}\to \mathcal {O}) @>\approx >\mathrm {d}> \sigma _{\geqslant 1}(\Omega _X^\bullet )       \\@AAA @AAA @AA{\mathrm {d}\log }A       \\\mathbb {Z}(1) @<<< (\mathbb {Z}(1)\to \mathcal {O}) @>\approx >\exp > \mathcal {O}^*[-1]       \end {CD}     \]\label{hodge-theory-ii-2.2.6}\par{}Let \(X\) be a non-singular projective variety of pure dimension \(n\).
    A choice of \(i\) defines an orientation of \(X\) and an isomorphism \(\mathbb {Z}(1)\simeq \mathbb {Z}\).
    The corresponding trace morphism
    \[       \operatorname {H}^{2n}(X,\mathbb {Z}(n)) \to  \mathbb {Z}     \]
    that is induced does not depend on the choice of \(i\).\par{}By Hodge, for \(i\leqslant  n\), the morphism
    \[       L^{n-i}       = -\wedge  c''_1(\mathcal {O}(1))^{n-i}\colon        \operatorname {H}^i(X,\mathbb {Z}) \to  \operatorname {H}^{2n-i}(X,\mathbb {Z}(n-i))     \]
    is an isomorphism, and, combined with Poincaré duality
    \[       \operatorname {H}^i(X,\mathbb {Z})\otimes \operatorname {H}^{2n-i}(X,\mathbb {Z}(n-i))       \to  \operatorname {H}^{2n}(X,\mathbb {Z}(n-i))       \to  \mathbb {Z}(-i)     \]
    gives a polarisation on the \emph{primitive part} \(\operatorname {Ker}(L^{n-i+1})\) of \(\operatorname {H}^i(X,\mathbb {Z})\).\par{}We thus deduce that the rational Hodge structures \(\operatorname {H}^i(X,\mathbb {Q})\) are polarisable.\section{Mixed structures}\label{hodge-theory-ii-2.3}\begin{definition}{2.3.1}\setnumber{2.3.1}\label{hodge-theory-ii-2.3.1}\skipsinglepar\par{}A \emph{mixed Hodge structure \(H\) consists of}

    
 \begin{enumerate}
      
 \item[(a)]
        A \(\mathbb {Z}\)-module \(H_\mathbb {Z}\) of finite type, called the "\emph{integral lattice}";
      

      
 \item[(b)]
        A finite increasing filtration \(W_n\) on \(H_\mathbb {Q}=\mathbb {Q}\otimes _\mathbb {Z} H_\mathbb {Z}\), called the \emph{weight filtration};
      

      
 \item[(c)]
        A finite filtration \(F^p\) on \(H_\mathbb {C}=\mathbb {C}\otimes _\mathbb {Z} H_\mathbb {Z}\), called the \emph{Hodge filtration}.
      

    \end{enumerate}


    We demand that, on \(H_\mathbb {C}\), the filtration \(W_\mathbb {C}\) induced by extension of scalars of \(W\), the filtration \(F\), and the complex conjugate \(\overline {F}\) form a system \((W_\mathbb {C},F,\overline {F})\) of three opposite filtrations (\hyperref[hodge-theory-ii-1.2.7]{Definition~\ref*{hodge-theory-ii-1.2.7}} and \cref{hodge-theory-ii-1.2.13}).\end{definition}\label{hodge-theory-ii-2.3.2}\par{}We also denote by \(W\) the filtration of \(H_\mathbb {Z}\) given by the inverse image of the filtration \(W\) on \(H_\mathbb {Q}\).
    The axiom of mixed Hodge structures implies that, for each \(n\), the filtration \(F\) induces on \(\mathbb {C}\otimes _\mathbb {Z}\operatorname {Gr}_n^W(H_\mathbb {Z})\) a filtration that is \(n\)-opposite to its complex conjugate.
    By \hyperref[hodge-theory-ii-2.1.9]{Proposition~\ref*{hodge-theory-ii-2.1.9}}, \(\operatorname {Gr}_n^W(H_\mathbb {Z})\) is endowed with a Hodge structure of weight \(n\), with Hodge filtration induced by \(F\).\begin{example}{2.3.3}\setnumber{2.3.3}\label{hodge-theory-ii-2.3.3}\skipsinglepar\par{}If \(H\) is a Hodge structure of weight \(n\), then we define a mixed Hodge structure with the same integral lattice and the same Hodge filtration by setting
    \[       W_i(H_\mathbb {Q}) =       \begin {cases}         0 &\text {for } i<n     \\H_\mathbb {Q} &\text {for } i\geqslant  n.       \end {cases}     \]\end{example}\label{hodge-theory-ii-2.3.4}\par{}A \emph{morphism} \(f\colon  H\to  H'\) of Hodge structures is a homomorphism \(f\colon  H_\mathbb {Z}\to  H'_\mathbb {Z}\) that is compatible with the filtrations \(W\) and \(F\) (and thus compatible with \(\overline {F}\)).
    We immediately deduce from \hyperref[hodge-theory-ii-1.2.10]{Theorem~\ref*{hodge-theory-ii-1.2.10}} the following theorem.\begin{theorem}{2.3.5}\setnumber{2.3.5}\label{hodge-theory-ii-2.3.5}\skipsinglepar\begin{enumerate}
    
 \item[(i)]
      The category of mixed Hodge structures is abelian.
    


    
 \item[(ii)]
      The integral lattice of the kernel (resp. cokernel) of a morphism \(f\colon  H\to  H'\) is the kernel (resp. cokernel) \(K\) of \(f\colon  H_\mathbb {Z}\to  H'_\mathbb {Z}\), with \(K\otimes \mathbb {Q}\) and \(K\otimes \mathbb {C}\) being endowed with the induced filtrations (resp. quotient filtrations) of the filtrations \(W\) and \(F\) of \(H_\mathbb {Q}\) and \(H_\mathbb {C}\) (resp. of \(H'_\mathbb {Q}\) and \(H'_\mathbb {C}\)).
    


    
 \item[(iii)]
      Every morphism \(f\colon  H\to  H'\) is strictly compatible with the filtration \(W\) of \(H_\mathbb {Q}\) and \(H'_\mathbb {Q}\) and with the filtration \(F\) of \(H_\mathbb {C}\) and \(H'_\mathbb {C}\).
      It induces morphisms of Hodge \(\mathbb {Q}\)-structures
      \[         \operatorname {Gr}_n^w(f)\colon          \operatorname {Gr}_n^W(H_\mathbb {Q})         \to  \operatorname {Gr}_n^W(H'_\mathbb {Q}).       \]
      It also induces morphisms
      \[         \operatorname {Gr}_F^p(F)\colon          \operatorname {Gr}_F^p(H_\mathbb {C})         \to  \operatorname {Gr}_F^p(H'_\mathbb {C})       \]
      that are strictly compatible with the filtration induced by \(W_\mathbb {C}\).
    


    
 \item[(iv)]
      The functor \(\operatorname {Gr}_n^W\) is an exact functor from the category of mixed Hodge structures to the category of Hodge \(\mathbb {Q}\)-structures of weight \(n\).
    


    
 \item[(v)]
      The functor \(\operatorname {Gr}_F^p\) is an exact functor.
    

  \end{enumerate}\end{theorem}\label{hodge-theory-ii-2.3.6}\par{}Let \(H\) be a mixed Hodge structure.
    The \(W_n(H_\mathbb {Z})\), endowed with the filtrations induced by \(W\) and \(F\), then form mixed Hodge substructures \(W_n(H)\) of \(H\).
    The quotient \(W_n(H)/W_{n-1}(H)\) can be identified with \(\operatorname {Gr}_n^W(H_\mathbb {Z})\) endowed with its mixed Hodge structure (\cref{hodge-theory-ii-2.3.2} and \hyperref[hodge-theory-ii-2.3.3]{Example~\ref*{hodge-theory-ii-2.3.3}}).
    We denote this Hodge structure by \(\operatorname {Gr}_n^W(H)\).\label{hodge-theory-ii-2.3.7}\par{}We set
    \[       H^{p,q}       = \operatorname {Gr}_F^p\operatorname {Gr}_{\overline {F}}^q\operatorname {Gr}_{p+q}^W(H_\mathbb {C})       = (\operatorname {Gr}_{p+q}^W(H))^{p,q}.     \]
    The \emph{Hodge numbers} of \(H\) are the integers
    \[       h^{p,q}       = \dim _\mathbb {C} H^{p,q}.     \]
    The Hodge number \(h^{pq}\) of \(H\) is thus the Hodge number \(h^{pq}\) of the Hodge structure \(\operatorname {Gr}_{p+q}^W(H)\)\label{hodge-theory-ii-2.3.8}\par{}We define a \emph{mixed Hodge \(\mathbb {Q}\)-structure} \(H\) to consist of a vector space \(H_\mathbb {Q}\) of finite dimension over \(\mathbb {Q}\), a finite increasing filtration \(W\) on \(H_\mathbb {Q}\), and a finite decreasing filtration \(F\) on \(H_\mathbb {C}\), with the filtrations \(W_\mathbb {C}\), \(F\), and \(\overline {F}\) being opposite.
    \hyperref[hodge-theory-ii-2.3.5]{Theorem~\ref*{hodge-theory-ii-2.3.5}} generalises trivially to this variation.\chapter{Hodge theory of non-singular algebraic varieties}\thispagestyle{fancy}\label{hodge-theory-ii-3}\section{Logarithmic poles and residues}\label{hodge-theory-ii-3.1}\label{hodge-theory-ii-3.1.1}\par{}We recall several classical properties of "logarithmic poles" of holomorphic differential forms.
    The reader will find proofs in [\hyperref[D1970]{D1970}, II, (3.1) to (3.7)], for example.\label{hodge-theory-ii-3.1.2}\par{}A divisor \(Y\) in a smooth complex analytic variety \(X\) is said to be of \emph{normal crossing} if the inclusion of \(Y\) into \(X\) is locally isomorphic to the inclusion of a union of coordinate hyperplanes in \(\mathbb {C}^n\);
    this does not imply that \(Y\) is the union of smooth divisors.
    Let \(Y\) be a normal crossing divisor in \(X\), and \(j\) the inclusion of \(X^*=X\setminus  Y\) into \(X\).
    We denote by \(\Omega _X^1\langle  Y\rangle \) the locally free sub-\(\mathcal {O}\)-module of \(j_*\Omega _{X^*}^1\) generated by \(\Omega _X^1\) and the \(\mathrm {d} z_i/z_i\) for \(z_i\) local coordinates of a local irreducible component of \(Y\).
    The \emph{sheaf \(\Omega _X^p\langle  Y\rangle \) of differential \(p\)-forms on \(X\) with logarithmic poles along \(Y\)} is by the definition the locally free sub-sheaf \(\wedge ^p\Omega _X^1\langle  Y\rangle \) of \(j_*\Omega _{X^*}^p\).\begin{proposition}{3.1.3}\setnumber{3.1.3}\label{hodge-theory-ii-3.1.3}\skipsinglepar\begin{enumerate}
    
 \item[(i)]
      A section \(\alpha \) of \(j_*\Omega _{X^*}^p\) belongs to \(\Omega _X^p\langle  Y\rangle \) if and only if \(\alpha \) and \(\mathrm {d}\alpha \) have at worst simple poles along the divisor \(Y\).
    

    
 \item[(ii)]
      The \(\Omega _X^p\langle  Y\rangle \) form the smallest subcomplex of \(j_*\Omega _{X^*}^\bullet \) that is stable under exterior product and that contains \(\Omega _X^\bullet \) and the logarithmic differential \(\mathrm {d} f/f\) of every local section of \(j_*\Omega _{X^*}^p\) meromorphic along \(Y\).
    

  \end{enumerate}\end{proposition}\par{}We call \(\Omega _X^\bullet \langle  Y\rangle \) the \emph{logarithmic de Rham complex} of \(X\) along \(Y\).
  By (ii) of \hyperref[hodge-theory-ii-3.1.3]{Proposition~\ref*{hodge-theory-ii-3.1.3}}, this complex is contravariant in the pair \((X,X^*)\).\label{hodge-theory-ii-3.1.4}\par{}Locally on \(X\), \(Y\) is the union of various smooth divisors \(Y_i\), and we denote by \(Y^n\) (resp. \(\widetilde {Y}^n\)) the union (resp. disjoint sum) of the \(n\)-fold intersections of the \(Y_i\).
    The \(Y^n\) glue to give a subspace \(Y^n\) of \(X\), and the \(\widetilde {Y}^n\) glue to give the normalisation variety of \(Y^n\).
    We have \(\widetilde {Y}^0=Y^0=X\), and we set \(\widetilde {Y}=\widetilde {Y}^1\).\par{}We define the two-element set of \emph{orientations} of an \(n\)-element set \(E\) as the set of generators of \(\wedge ^n\mathbb {Z}^E\).
    
    For \(n\geqslant 2\), this set is that of the conjugation classes under the alternating group of total orders on \(E\).\par{}If, to each point \(y\) of \(\widetilde {Y}^n\), we associate the set of \(n\) local components of \(Y\) that contain the image in \(X\) of a neighbourhood of \(y\) in \(\widetilde {Y}^n\), then we define on \(\widetilde {Y}^n\) a local system \(E_n\) of \(n\)-element sets.
    The local system of orientations of these sets is a \(\mathbb {Z}{/}(2)\)-torsor.
    This torsor defines, via the inclusion of \(\mathbb {Z}{/}(2)\) into \(\mathbb {C}^*\), a complex local system \(\varepsilon ^n\) of rank 1 on \(\widetilde {Y}^n\), endowed with an isomorphism \((\varepsilon ^n)^{\otimes 2}\simeq \mathbb {C}\).
    We have
    \[       \varepsilon ^n       \simeq  \bigwedge ^n \mathbb {C}^{E_n}.     \]\par{}Locally on \(\widetilde {Y}^n\), \(\varepsilon ^n\) is endowed with two opposite isomorphisms \(\pm \alpha \colon \varepsilon ^n\xrightarrow {\sim }\mathbb {C}\).
    We set
    \[       \varepsilon _\mathbb {Z}^n       = \alpha ^{-1}((2\pi  i)^{-n}\mathbb {Z}).     \]
    There is a way of seeing \(\varepsilon ^n\), endowed with \(\varepsilon _\mathbb {Z}^n\), as a twisted version of \(\mathbb {Z}(-n)\) over \(\widetilde {Y}^n\).\par{}We denote by \(\varepsilon _X^n\) (resp. by \((\varepsilon _X^n)_\mathbb {Z}\)) the direct image of \(\varepsilon ^n\) (resp. of \(\varepsilon _\mathbb {Z}^n\)) under the mat from \(\widetilde {Y}^n\) to \(X\).
    We have

    \begin{eqn}{3.1.4.1}\setnumber{3.1.4.1}\label{hodge-theory-ii-3.1.4.1}\skipsinglepar\[         \varepsilon _X^n         \simeq  \bigwedge ^n\varepsilon _X^1         \qquad \text {for }n\geqslant 0.       \tag{3.1.4.1}       \]\end{eqn}\par{}If \(Y\) is the union of distinct smooth divisors \((Y_i)_{i\in  I}\), then the choice of a total order on \(I\) trivialises the \(\varepsilon ^n\).\label{hodge-theory-ii-3.1.5}\par{}We denote by \(W_n(\Omega _X^p\langle  Y\rangle )\) the sub-module of \(\Omega _X^p\langle  Y\rangle \) consisting of linear combinations of products
    \[       \alpha        \wedge  \frac {\mathrm {d} t_{i(1)}}{t_{i(1)}}       \ldots        \wedge  \frac {\mathrm {d} t_{i(m)}}{t_{i(m)}}       \qquad \text {for }m\leqslant  n     \]
    with \(\alpha \) holomorphic and the \(t_{i(j)}\) local coordinates of the distinct local components \(Y_j\) of \(Y\).
    We define the \(weight filtration\) of \(\Omega _X^\bullet \langle  Y\rangle \) to be the \emph{increasing} filtration by the subcomplexes \(W_n(\Omega _X^\bullet \langle  Y\rangle )\).
    We have

    \begin{eqn}{3.1.5.1}\setnumber{3.1.5.1}\label{hodge-theory-ii-3.1.5.1}\skipsinglepar\[         W_n(\Omega _X^p\langle  Y\rangle )         \wedge  W_m(\Omega _X^p\langle  Y\rangle )         \subset  W_{n+m}(\Omega _X^{p+q}\langle  Y\rangle ).       \tag{3.1.5.1}       \]\end{eqn}\par{}If we denote by \(i_n\) the map from \(\widetilde {Y}^n\) to \(X\), then we can show that the mapping
    \[       \alpha        \wedge  \frac {\mathrm {d} t_{i(1)}}{t_{i(1)}}       \ldots        \wedge  \frac {\mathrm {d} t_{i(m)}}{t_{i(m)}}       \longmapsto        (\alpha |Y_{i(1)}\cap \ldots \cap  Y_{i(n)})\otimes (\text {orientation }i(1)\ldots  i(n))     \]
    defines an isomorphism of complexes

    \begin{eqn}{3.1.5.2}\setnumber{3.1.5.2}\label{hodge-theory-ii-3.1.5.2}\skipsinglepar\[         \operatorname {Res}\colon \operatorname {Gr}_n^W(\Omega _X^\bullet \langle  Y\rangle )         \approx  (i_n)_*\Omega _{\widetilde {Y}^n}^\bullet (\varepsilon ^n)[-n]       \tag{3.1.5.2}       \]\end{eqn}
    
    (the \emph{Poincaré residue}).\label{hodge-theory-ii-3.1.6}\par{}The interpretation that follows, in terms of \hyperref[hodge-theory-ii-3.1.5.2]{Equation~\ref*{hodge-theory-ii-3.1.5.2}}, of the Leray spectral sequence for the inclusion of \(X^*\) into \(X\), was pointed out to me by N. Katz.
    It will allow us to prove a point that I had initially considered as evident (the first part of (ii) of \hyperref[hodge-theory-ii-3.2.5]{Theorem~\ref*{hodge-theory-ii-3.2.5}}).\label{hodge-theory-ii-3.1.7}\par{}Every point of \(X\) admits a fundamental system of Stein open neighbourhoods whose intersections on \(X^*\) are again Stein.
    For a coherent analytic sheaf \(\mathscr {F}\) on \(X^*\), we thus have \(\mathrm {R}^i j_*\mathscr {F}=0\) for \(i>0\).
    The de Rham complex \(\Omega _{X^*}^\bullet \) is thus a resolution of the constant sheaf \(\mathbb {C}\) by sheaves acyclic for the functor \(j_*\).
    Then

    \begin{eqn}{3.1.7.1}\setnumber{3.1.7.1}\label{hodge-theory-ii-3.1.7.1}\skipsinglepar\[         \operatorname {H}^\bullet (X^*,\mathbb {C})         \xrightarrow {\sim } \operatorname {\mathbb {H}}^\bullet (X^*,\Omega _{X^*}^\bullet )         \xleftarrow {\sim } \operatorname {\mathbb {H}}^\bullet (X,j_*\Omega _{X^*}^\bullet )       \]\end{eqn}

    and the Leray spectral sequence for the morphism \(j\) can be identified with the hypercohomology spectral sequence for \(j_*\Omega _{X^*}^\bullet \) corresponding to the filtration \(\tau \) by the subcomplexes \(\tau _{\leqslant  -n}(j_*\Omega _{X^*}^\bullet )\) (\cref{hodge-theory-ii-1.4.6}).\begin{proposition}{3.1.8}\setnumber{3.1.8}\label{hodge-theory-ii-3.1.8}\skipsinglepar\par{}The morphisms of filtered complexes
    \[       (\Omega _X^\bullet \langle  Y\rangle , W)       \xleftarrow {\alpha } (\Omega _X^\bullet \langle  Y\rangle , \tau )       \xhookrightarrow {\beta } (j_*\Omega _{X^*}^\bullet , \tau )     \]
    are filtered quasi-isomorphisms.
    They define an isomorphism between the Leray spectral sequence for \(j\) in complex cohomology and the hypercohomology spectral sequence of the filtered complex \((\Omega _X^\bullet \langle  Y\rangle , W)\) on \(X\).\begin{proof}\par{}By \cref{hodge-theory-ii-1.4.5} and \cref{hodge-theory-ii-3.1.7}, it suffices to prove the first claim.
      In [\hyperref[D1970]{D1970}, II, (6.9)] or [\hyperref[AH1955]{AH1955}], one can find a proof of the fact that \(\beta \) is a quasi-isomorphism, and thus a filtered quasi-isomorphism.
      We can also directly calculate the cohomology sheaves of the two sides: those of \(\Omega _X^\bullet \langle  Y\rangle \) are determined by \hyperref[hodge-theory-ii-3.1.5.2]{Equation~\ref*{hodge-theory-ii-3.1.5.2}}, while those of \(j_*\Omega _{X^*}^\bullet \) are the \(\mathrm {R}^ij_*\mathbb {C}\), which can be calculated by topological methods.\par{}For \(n\geqslant  p\), we have
      \[         W_n(\Omega _X^p\langle  Y\rangle )         = \Omega _X^p\langle  Y\rangle        \]
      and so \(\alpha \) is a morphism from \(\Omega _X^\bullet \langle  Y\rangle \), endowed with \(\tau \), to \(\Omega _X^\bullet \langle  Y\rangle \), endowed with the decreasing filtration associated to \(W\) (\cref{hodge-theory-ii-1.1.3}).
      By \hyperref[hodge-theory-ii-3.1.5.2]{Equation~\ref*{hodge-theory-ii-3.1.5.2}}, we have

      \begin{eqn}{3.1.8.1}\setnumber{3.1.8.1}\label{hodge-theory-ii-3.1.8.1}\skipsinglepar\[           \operatorname {\mathscr {H}}^i(\operatorname {Gr}_n^W(\Omega _X^\bullet \langle  Y\rangle )) =           \begin {cases}             0 &\text {for }i\neq  n           \\\varepsilon _X^n &\text {for }i=n           \end {cases}         \tag{3.1.8.1}         \]\end{eqn}

      and we deduce from the first line of this equation that \(\alpha \) is a filtered quasi-isomorphism.
      This proves \hyperref[hodge-theory-ii-3.1.8]{Proposition~\ref*{hodge-theory-ii-3.1.8}}.\end{proof}\par{}By \cref{hodge-theory-ii-3.1.7}, the isomorphism \hyperref[hodge-theory-ii-3.1.8.1]{Equation~\ref*{hodge-theory-ii-3.1.8.1}} defines an isomorphism

    \begin{eqn}{3.1.8.2}\setnumber{3.1.8.2}\label{hodge-theory-ii-3.1.8.2}\skipsinglepar\[         \mathrm {R}^n j_*\mathbb {C}         \simeq  \operatorname {\mathscr {H}}^n(j_*\Omega _{X^*}^\bullet )         \simeq  \operatorname {\mathscr {H}}^n(\Omega _X^\bullet \langle  Y\rangle )         \simeq  \varepsilon _X^n.       \tag{3.1.8.2}       \]\end{eqn}

    The isomorphisms \hyperref[hodge-theory-ii-3.1.4.1]{Equation~\ref*{hodge-theory-ii-3.1.4.1}} correspond, via \hyperref[hodge-theory-ii-3.1.8.1]{Equation~\ref*{hodge-theory-ii-3.1.8.1}}, to the cup product.\end{proposition}\begin{proposition}{3.1.9}\setnumber{3.1.9}\label{hodge-theory-ii-3.1.9}\skipsinglepar\par{}The canonical morphism from \(\mathrm {R}^nj_*\mathbb {Z}\) to \(\mathrm {R}^nj_*\mathbb {C}\) identifies, via \hyperref[hodge-theory-ii-3.1.8.2]{Equation~\ref*{hodge-theory-ii-3.1.8.2}}, the sheaf \(\mathrm {R}^nj_*\mathbb {Z}\) with \((\varepsilon _X^n)_\mathbb {Z}\) (\cref{hodge-theory-ii-3.1.4}).\begin{proof}\par{}The question is local on \(X\).
      We can thus suppose that \(X\) is an open polycylinder \(D^m\), with
      \[         D         = \{z\in \mathbb {C} \mid  |z|<1\}       \]
      and also that \(Y=\bigcup _{k=1}^\ell  Y_k\) with \(Y_k=\operatorname {pr}_k^{-1}(0)\).
      The fibre at \(0\) of \(\mathrm {R}^nj_*\mathbb {Z}\) is then the integer cohomology of \(X^*=(D^*)^\ell \times  D^{m-\ell }\), with
      \[         D^*         = \{z\in \mathbb {C} \mid  0<|z|<1\}.       \]
      
      The space \(X^*\) has the homotopy type of a torus;
      its cohomology is thus torsion free, and the cup product defines isomorphisms
      \[         \bigwedge ^n(\mathrm {R}^1j_*\mathbb {Z})         \xrightarrow {\sim } (\mathrm {R}^nj_*\mathbb {Z})_0.       \]
      It thus suffices to prove \hyperref[hodge-theory-ii-3.1.9]{Proposition~\ref*{hodge-theory-ii-3.1.9}} for \(n=1\).\par{}The integer homology \(\operatorname {H}_1(X^*)\) is generated by the loops \(\gamma _k\) that go around the various \(Y_k\).
      We have
      \[         \oint _{\gamma _k} \frac {\mathrm {d} z_k}{z_k}         = \pm 2\pi  i       \]
      and so the integer cohomology is generated by the \(\frac {1}{2\pi  i}\frac {\mathrm {d} z_k}{z_k}\), and this proves \hyperref[hodge-theory-ii-3.1.9]{Proposition~\ref*{hodge-theory-ii-3.1.9}}.\end{proof}\end{proposition}\label{hodge-theory-ii-3.1.10}\par{}Let \(\mathscr {F}\) be a coherent analytic sheaf on \(X^*\), given as the restriction to \(X^*\) of a coherent analytic sheaf \(\mathscr {F}'\) on \(X\).
    We define the \emph{meromorphic direct image} \(j_*^\mathrm {m}\mathscr {F}\) of \(\mathscr {F}\) to be the inductive limit
    \[       j_*^\mathrm {m}\mathscr {F}       = \varinjlim \mathscr {F}'(nY).     \]\par{}Locally on \(X\), \(Y\) is the sum of a finite family \((Y_i)_{i\in  I}\) of smooth divisors, and we define the \emph{pole-order filtration} \(P\) on \(j_*^\mathrm {m}\mathcal {O}_X^*\) by the equation

    \begin{eqn}{3.1.10.1}\setnumber{3.1.10.1}\label{hodge-theory-ii-3.1.10.1}\skipsinglepar\[         P^p(j_*^\mathrm {m}\mathcal {O}_{X^*})         = \sum _{n\in  A_p} \mathcal {O}_X\left (\sum (n_i+1)Y_i\right )       \tag{3.1.10.1}       \]\end{eqn}

    where
    \[       A_p       = \left \{(n_i)_{i\in  I} \mid  \sum _i n_i\leqslant  -p \text { and } n_i\geqslant 0\text { for all } i\right \}.     \]
    This construction globalises by endowing \(j_*^\mathrm {m}\mathcal {O}_{X^*}\) with an exhaustive filtration such that \(P^p=0\) for \(p>0\).\par{}We define the \emph{pole-order filtration} of the complex \(j_*^\mathrm {m}\Omega _{X^*}^\bullet =j_*^\mathrm {m}\mathcal {O}_{X^*}\otimes \Omega _X^\bullet \) to be the filtration

    \begin{eqn}{3.1.10.2}\setnumber{3.1.10.2}\label{hodge-theory-ii-3.1.10.2}\skipsinglepar\[         P^p(j_*^\mathrm {m}\Omega _{X^*}^k)         = P^{p-k}(j_*^\mathrm {m}\mathcal {O}_X)\otimes \Omega _X^k.       \tag{3.1.10.2}       \]\end{eqn}\par{}The filtration \(P\) induces, on the subcomplex \(\Omega _X^\bullet \langle  Y\rangle \) of \(j_*^\mathrm {m}\Omega _{X^*}^\bullet \), the stupid filtration by the \(\sigma _{\geqslant  p}(\Omega _X^\bullet \langle  Y\rangle )\), which we also call the \emph{Hodge filtration} \(F\).\begin{proposition}{3.1.11}\setnumber{3.1.11}\label{hodge-theory-ii-3.1.11}\skipsinglepar\par{}The inclusion morphism
    \[       (\Omega _X^\bullet \langle  Y\rangle , F)       \to  (j_*^\mathrm {m}\Omega _{X_*}^\bullet , P)     \]
    is a filtered quasi-isomorphism.\begin{proof}\par{}This statement was suggested to me by [\hyperref[G1969]{G1969}].
      A proof can be found in [\hyperref[D1970]{D1970}, II, (3.13)].\end{proof}\end{proposition}\section{Mixed Hodge theory}\label{hodge-theory-ii-3.2}\par{}\emph{We recall that, from here on, we say "scheme" to mean a scheme of finite type over \(\mathbb {C}\), and "sheaf on \(S\)" to mean a sheaf on \(S^\mathrm {an}\).}\label{hodge-theory-ii-3.2.1}\par{}Let \(X\) be a smooth and separated scheme.
    By Nagata [\hyperref[N1962]{N1962}], \(X\) is a Zariski open of a complete scheme \(\overline {X}\).
    By Hironaka [\hyperref[H1964]{H1964}], we can take \(\overline {X}\) to be smooth and such that \(Y=\overline {X}\setminus  X\) is a normal crossing divisor.\par{}The reader who wishes to avoid the reference to Nagata can suppose \(X\) to be quasi-projective.
    
    The smooth completion \(\overline {X}\) can then be chosen to be projective and such that \(Y\) is the union of smooth divisors.
    If we limit ourselves to such compactifications, then we only need Hodge theory in its standard form (\cref{hodge-theory-ii-2.2.1}).\label{hodge-theory-ii-3.2.2}\par{}By \cref{hodge-theory-ii-3.1.7} and \hyperref[hodge-theory-ii-3.1.8]{Proposition~\ref*{hodge-theory-ii-3.1.8}}, we have
    \[       \operatorname {H}^\bullet (X,\mathbb {C})       \simeq  \operatorname {\mathbb {H}}^\bullet (\overline {X},\Omega _{\overline {X}}^\bullet \langle  Y\rangle ).     \]
    We define the \emph{Hodge filtration} \(F\) on the complex \(\Omega _{\overline {X}}^\bullet \langle  Y\rangle \) as the filtration \(F^p=\sigma _{\geqslant  p}\) given by the stupid truncations (\cref{hodge-theory-ii-1.4.5}).
    On \(\Omega _{\overline {X}}^\bullet \langle  Y\rangle \), we thus have two filtrations: \(F\) and \(W\) (\cref{hodge-theory-ii-3.1.5}).\label{hodge-theory-ii-3.2.3}\par{}We will need to make use of the fact that there exist bifiltered resolutions \(i\colon \Omega _{\overline {X}}^\bullet \langle  Y\rangle \to  K^\bullet \) such that the \(\operatorname {Gr}_F^p\operatorname {Gr}_n^W(K^j)\) are \(\Gamma \)-acyclic sheaves:
    \[       \operatorname {H}^i(\overline {X},\operatorname {Gr}_F^p\operatorname {Gr}_n^W(K^j))       = 0       \qquad \text {for }i>0.     \]\par{}Here are two methods to construct such a resolution:

    
 \begin{enumerate}
      
 \item[(a)]
        We can take \(K^\bullet \) to be the canonical Godement resolution \(\mathscr {C}^\bullet (\Omega _{\overline {X}}^\bullet \langle  Y\rangle )\), filtered by the \(\mathscr {C}^\bullet (W_n(\Omega _{\overline {X}}\langle  Y\rangle ))\) and the \(\mathscr {C}^\bullet (F^p(\Omega _{\overline {X}}^\bullet \langle  Y\rangle ))\).
        This is a bifiltered resolution since \(\mathscr {C}^\bullet \) is an exact functor.
      


      
 \item[(b)]
        We can take \(K^\bullet \) to be the \(\mathrm {d}''\)-resolution of \(\Omega _{\overline {X}}^\bullet \langle  Y\rangle \).
        Let \(\Omega _{\overline {X}}^{pq}\) be the sheaf of \(C^\infty \) forms of degree \((p,q)\);
        then \(K^\bullet \) is the simple complex associated to the double complex of the \(\Omega _{\overline {X}}^p\langle  Y\rangle \otimes _\mathcal {O}\Omega _{\overline {X}}^{0,q}\) (a subcomplex of the \(j_*\Omega _X^{\bullet \bullet }\)).
        This complex is filtered by the \(F^p(\Omega _{\overline {X}}^\bullet \langle  Y\rangle )\otimes \Omega _{\overline {X}}^{0,\bullet }\) and by the \(W_n(\Omega _{\overline {X}}^\bullet \langle  Y\rangle )\otimes \Omega _{\overline {X}}^{0,\bullet }\);
        to prove that this is a bifiltered resolution, we use the fact that the sheaf \(\mathcal {O}_\infty \) of complex \(C^\infty \) functions on \(\overline {X}\) is flat over \(\mathcal {O}\) (a corollary of the Malgrange \(C^\infty \) preparation theorem).
        The sheaves \(\operatorname {Gr}_F\operatorname {Gr}_W(K^\bullet )\) are fine, since they are sheaves of modules over the soft sheaf \(\mathcal {O}_\infty \).
      

    \end{enumerate}\label{hodge-theory-ii-3.2.4}\par{}With the notation of \cref{hodge-theory-ii-3.2.3}, the complex cohomology of \(X\) appears as the cohomology of the bifiltered complex \(\Gamma (\overline {X},K^\bullet )\).
    We thus have two spectral sequences abutting to \(\operatorname {H}^\bullet (X,\mathbb {C})\).
    They can be written, with the notation of \cref{hodge-theory-ii-3.1.4}, as:

    \begin{eqn}{3.2.4.1}\setnumber{3.2.4.1}\label{hodge-theory-ii-3.2.4.1}\skipsinglepar\[         {}_WE_1^{pq}         = \operatorname {\mathbb {H}}^{p+q}(\overline {X},\varepsilon _{\overline {X}}^{-p}[p])         = \operatorname {\mathbb {H}}^{2p+q}(\widetilde {Y}^p,\varepsilon ^{-q})         \Rightarrow  \operatorname {H}^n(X,\mathbb {C})       \tag{3.2.4.1}       \]\end{eqn}

    \begin{eqn}{3.2.4.2}\setnumber{3.2.4.2}\label{hodge-theory-ii-3.2.4.2}\skipsinglepar\[         {}_FE_1^{pq}         = \operatorname {H}^q(\overline {X},\Omega _{\overline {X}}^p\langle  Y\rangle )         \Rightarrow  \operatorname {H}^n(X,\mathbb {C}).       \tag{3.2.4.2}       \]\end{eqn}

    The first of these, up to the renumbering \({}_WE_1^{pq}\mapsto  E_2^{2p+q,-p}\), is exactly the Leray spectral sequence of the inclusion \(j\).\begin{theorem}{3.2.5}\setnumber{3.2.5}\label{hodge-theory-ii-3.2.5}\skipsinglepar\begin{enumerate}
    
 \item[(i)]
      On the pages \({}_WE_r^{pq}\) of the spectral sequence \hyperref[hodge-theory-ii-3.2.4.1]{Equation~\ref*{hodge-theory-ii-3.2.4.1}}, the first direct filtration, the second direct filtration, and the recurrent filtration defined by \(F\) all coincide.
    


    
 \item[(ii)]
      The filtration on \(\operatorname {H}^n(X,\mathbb {C})\) that is the abutment of the spectral sequence \({}_WE\) comes from a filtration \(W\) of \(\operatorname {H}^n(X,\mathbb {Q})\).
      Neither it, nor the filtration \(F\) that is the abutment of the spectral sequence \({}_FE\), depend on the choice of compactification \(\overline {X}\) of \(X\) or on the choice of \(K^\bullet \).
    


    
 \item[(iii)]
      The filtrations \(W[n]\) (\hyperref[hodge-theory-ii-1.1.2]{Definition~\ref*{hodge-theory-ii-1.1.2}}) and \(F\) define on \(\operatorname {H}^n(X,\mathbb {Z})\) a mixed Hodge structure, functorially in \(X\).
    

  \end{enumerate}\end{theorem}\par{}By \cref{hodge-theory-ii-3.1.7}, the spectral sequence \({}_WE\) is the Leray spectral sequence for \(j_*\) (up to renumbering).
    It is thus induced by tensoring a \(\mathbb {Q}\)-vectorial spectral sequence with \(\mathbb {C}\), and so the first claim of (ii) is true.
    Complex conjugation acts on the \({}_WE\);
    it can be calculated via \cref{hodge-theory-ii-3.1.10}.\begin{lemma}{3.2.6}\setnumber{3.2.6}\label{hodge-theory-ii-3.2.6}\skipsinglepar\par{}The hypercohomology spectral sequences of the filtered complexes \(\operatorname {Gr}_n^W(\Omega _{\overline {X}}^\bullet \langle  Y\rangle )\), endowed with the filtration induced by the Hodge filtration, degenerate at the \(E_1\) page.\begin{proof}\par{}Define the \(Y^n\) and the \(\widetilde {Y}^n\) as in \cref{hodge-theory-ii-3.1.4}, and let \(i_n\colon \widetilde {Y}^n\to  X\).
      By \hyperref[hodge-theory-ii-3.1.5.2]{Equation~\ref*{hodge-theory-ii-3.1.5.2}}, we have
      \[         \operatorname {Gr}_n^W(\Omega _{\overline {X}}^\bullet \langle  Y\rangle )         \sim  (i_n)_*\Omega _{\widetilde {Y}^n}^\bullet (\varepsilon ^n)[-n].       \]
      Furthermore, the Hodge filtration induces the stupid filtration (\cref{hodge-theory-ii-1.4.5}), and so the spectral sequence in \hyperref[hodge-theory-ii-3.2.6]{Lemma~\ref*{hodge-theory-ii-3.2.6}} is given by translations of the degrees of the classical spectral sequence
      \[         E_1^{pq}         = \operatorname {H}^q(\widetilde {Y}^n,\Omega _{\widetilde {Y}^n}^p(\varepsilon ^n))         \Rightarrow  \operatorname {H}^{p+q}(\widetilde {Y}^n,\varepsilon ^n).       \]\par{}If \(\overline {X}\) is projective and \(Y\) the union of smooth divisors, then \(\widetilde {Y}^n\) is projective, \(\varepsilon ^n\) is a trivial local system, and classical Hodge theory (\cref{hodge-theory-ii-2.2.1}) gives the degeneration claimed in \hyperref[hodge-theory-ii-3.2.6]{Lemma~\ref*{hodge-theory-ii-3.2.6}}.
      For the general case, refer to [\hyperref[D1968]{D1968}] (see also \cref{hodge-theory-ii-2.2.2} and \cref{hodge-theory-ii-2.2.3}).\end{proof}\end{lemma}\par{}Hodge theory also tells us that the Hodge filtration on \(\operatorname {H}^k(\widetilde {Y}^n,\varepsilon ^n)\) is \(k\)-opposite to its complex conjugate (the complex conjugate being defined in terms of \(\varepsilon _\mathbb {Z}^n\) (\cref{hodge-theory-ii-3.1.4})).
  We have, taking into account the translations of the degrees:\begin{lemma}{3.2.7}\setnumber{3.2.7}\label{hodge-theory-ii-3.2.7}\skipsinglepar\par{}The filtration on
    \[       {}_WE_1^{-n,k+n}       = \operatorname {\mathbb {H}}^k(\overline {X},\operatorname {Gr}_n^W(\Omega _{\overline {X}}^\bullet \langle  Y\rangle ))       \approx  \operatorname {H}^{k-n}(\widetilde {Y}^n,\varepsilon ^n)     \]
    which is the abutment of the spectral sequence in \hyperref[hodge-theory-ii-3.2.6]{Lemma~\ref*{hodge-theory-ii-3.2.6}} is \((k+n)\)-opposite to its complex conjugate.\end{lemma}\begin{lemma}{3.2.8}\setnumber{3.2.8}\label{hodge-theory-ii-3.2.8}\skipsinglepar\par{}The differentials \(d_1\) of the spectral sequence \({}_WE\) are strictly compatible with the filtration \(F\).\begin{proof}\par{}On the \(E_1\) pages, there is only the filtration induced by \(F\) to consider (\hyperref[hodge-theory-ii-1.3.10]{Lemma~\ref*{hodge-theory-ii-1.3.10}} and (iii) of \hyperref[hodge-theory-ii-1.3.13]{Proposition~\ref*{hodge-theory-ii-1.3.13}}), and \(d_1\) is compatible with this filtration ((i) of \hyperref[hodge-theory-ii-1.3.13]{Proposition~\ref*{hodge-theory-ii-1.3.13}}).
      This filtration is the abutment of the spectral sequence in \hyperref[hodge-theory-ii-3.2.6]{Lemma~\ref*{hodge-theory-ii-3.2.6}} ((ii) of \hyperref[hodge-theory-ii-1.4.8]{Example~\ref*{hodge-theory-ii-1.4.8}}).
      By \hyperref[hodge-theory-ii-3.2.7]{Lemma~\ref*{hodge-theory-ii-3.2.7}}, the arrow
      \[         d_1\colon          \operatorname {\mathbb {H}}^k(\overline {X},\operatorname {Gr}_n^W(\Omega _{\overline {X}}^\bullet \langle  Y\rangle ))         \to  \operatorname {\mathbb {H}}^{k+1}(\overline {X},\operatorname {Gr}_{n-1}^W(\Omega _{\overline {X}}^\bullet \langle  Y\rangle ))       \]
      or

      \begin{eqn}{3.2.8.1}\setnumber{3.2.8.1}\label{hodge-theory-ii-3.2.8.1}\skipsinglepar\[           d_1\colon            \operatorname {H}^{k-n}(\widetilde {Y}^n,\varepsilon ^n)           \to  \operatorname {H}^{k-n+1}(\widetilde {Y}^{n-1},\varepsilon ^{n-1})         \tag{3.2.8.1}         \]\end{eqn}

      is compatible with the filtrations that are \((k+n)\)-opposite to their complex conjugate.
      Since \(d_1\) commutes up to complex conjugation, \(d_1\) respects the bigrading (of weight \(k+n\)) defined by \(F\) and \(\overline {F}\), which proves \hyperref[hodge-theory-ii-3.2.8]{Lemma~\ref*{hodge-theory-ii-3.2.8}}.\end{proof}\end{lemma}\par{}Furthermore, the cohomology of the complex \(E_1\) is again bigraded:\begin{lemma}{3.2.9}\setnumber{3.2.9}\label{hodge-theory-ii-3.2.9}\skipsinglepar\par{}The recurrent filtration \(F\) on \({}_WE_2^{pq}\) is \(q\)-opposite to its complex conjugate.\end{lemma}\par{}We prove by induction on \(r\) that:\begin{lemma}{3.2.10}\setnumber{3.2.10}\label{hodge-theory-ii-3.2.10}\skipsinglepar\par{}For \(r\geqslant 0\), the differentials \(d_r\) of the spectral sequence \({}_WE\) are strictly compatible with the recurrent filtration \(F\).
    For \(r\geqslant 2\), they are zero.\begin{proof}\par{}For \(r=0\) (resp. \(r=1\)) we apply \hyperref[hodge-theory-ii-3.2.6]{Lemma~\ref*{hodge-theory-ii-3.2.6}} and (iii) of \hyperref[hodge-theory-ii-1.4.8]{Example~\ref*{hodge-theory-ii-1.4.8}} (resp. \hyperref[hodge-theory-ii-3.2.8]{Lemma~\ref*{hodge-theory-ii-3.2.8}}).
      For \(r\geqslant 2\), it suffices to prove that \(d_r=0\).
      By induction, and following \hyperref[hodge-theory-ii-1.3.16]{Theorem~\ref*{hodge-theory-ii-1.3.16}}, we can assume that, on the pages \({}_WE^s\) (for \(s\geqslant  r+1\)), we have \(F_d=F_r=F_{d^*}\), and that \({}_WE_r={}_WE_2\).
      By (i) of \hyperref[hodge-theory-ii-1.3.13]{Proposition~\ref*{hodge-theory-ii-1.3.13}}, \(d_r\) is thus compatible with the filtration \(F_r\).\par{}On \({}_WE_r^{pq}={}_WE_2^{pq}\), the filtration \(F_r\) is \(q\)-opposite to its complex conjugate.
      The morphism
      \[         d_r\colon  {}_WE_r^{pq}         \to  {}_WE_r^{p+r,q-r+1}       \]
      thus satisfies, for \(r-1>0\),
      \[         \begin {aligned}           d_r({}_WE_r^{pq})           &= d_r\left (             \sum _{a+b=q} \Big (               F^a({}_WE_r^{pq}) \cap  \overline {F}^b({}_WE_r^{pq})             \Big )           \right )         \\&\subset  \sum _{a+b=q} \Big (             F^a({}_WE_r^{p+r,q-r+1}) \cap  \overline {F}^b({}_WE_r^{p+r,q-r+1})           \Big )         \\&= 0.         \end {aligned}       \]\par{}This proves \hyperref[hodge-theory-ii-3.2.10]{Lemma~\ref*{hodge-theory-ii-3.2.10}}.\end{proof}\end{lemma}\par{}\hyperref[hodge-theory-ii-3.2.10]{Lemma~\ref*{hodge-theory-ii-3.2.10}}, using \hyperref[hodge-theory-ii-1.3.16]{Theorem~\ref*{hodge-theory-ii-1.3.16}}, proves (i) of \hyperref[hodge-theory-ii-3.2.5]{Theorem~\ref*{hodge-theory-ii-3.2.5}}.\par{}By \hyperref[hodge-theory-ii-1.3.17]{Corollary~\ref*{hodge-theory-ii-1.3.17}}, the filtration on \({}_WE_\infty ^{pq}\) induced by the filtration \(F\) of \(\operatorname {H}^{p+q}(X,\mathbb {C})\) is \(q\)-opposite to its complex conjugate.
  Since \(q=-p+(p+q)\), this proves the first part of (iii) of \hyperref[hodge-theory-ii-3.2.5]{Theorem~\ref*{hodge-theory-ii-3.2.5}}\label{hodge-theory-ii-3.2.11}\par{}We now prove (ii) and (iii) of \hyperref[hodge-theory-ii-3.2.5]{Theorem~\ref*{hodge-theory-ii-3.2.5}}, which will finish the proof.\begin{enumerate}
    
 \item[(A)]
      \par{}\emph{Independence from the choice of \(K^\bullet \).}

      \par{}The filtrations \(F\) and \(W\) of \(\operatorname {H}^\bullet (X,\mathbb {C})\) are the abutments of the hypercohomology spectral sequences of \(\Omega _{\overline {X}}^\bullet \langle  Y\rangle \) for the filtrations \(F\) and \(W\).
        These whole spectral sequences do not depend on the choice of \(K^\bullet \).
    


    
 \item[(B)]
      \par{}\emph{Functoriality.}

      \par{}Let \(f\colon  X_1\to  X_2\) be a morphism of schemes.
        Suppose that we are given a morphism of smooth compactifications

        \label{hodge-theory-ii-3.2.11.1}\[             \begin {CD}               X_1 @>f>> X_2             \\@V{j_1}VV @VV{j_2}V             \\\overline {X}_1 @>>{\overline {f}}> \overline {X}_2             \end {CD}           \tag{3.2.11.1}           \]

        with the \(Y_i=\overline {X}_i\setminus  X_i\) normal crossing divisors.
        The canonical morphism (cf. \hyperref[hodge-theory-ii-3.1.3]{Proposition~\ref*{hodge-theory-ii-3.1.3}}) from \(\overline {f}^*\Omega _{\overline {X}_2}^\bullet \langle  Y_2\rangle \) to \(\Omega _{\overline {X}_1}^\bullet \langle  Y_1\rangle \) is then a morphism of bifiltered complexes;
        on hypercohomology, it induces a morphism compatible with \(F\) and \(W\), and thus \(f^*\colon \operatorname {H}^n(X_2,\mathbb {Z})\to \operatorname {H}^n(X_1,\mathbb {Z})\) is a morphism of mixed Hodge structures, for the structures defined by the compactifications \(\overline {X}_i\).
    


    
 \item[(C)]
      \par{}\emph{Independence from the choice of compactification.}

      \par{}With the notation of (B), if \(f\) is an isomorphism, then \(f^*\) is a bijective morphism of mixed Hodge structures, and thus an isomorphism (\hyperref[hodge-theory-ii-2.3.5]{Theorem~\ref*{hodge-theory-ii-2.3.5}}).

      \par{}If \(\overline {X}_1\) and \(\overline {X}_2\) are two smooth compactifications of \(X\), with \(Y_i=\overline {X}_i\setminus  X\) smooth crossing divisors, then there exists a third smooth compactification \(\overline {X}\), with \(Y=\overline {X}\setminus  X\) a smooth crossing divisor, that fits into a commutative diagram
        
        
  \begin {tikzcd}
          & X \ar [dl,hook] \ar [d,hook] \ar [dr,hook]
        \\\overline {X}_1
          & \overline {X} \ar [l] \ar [r]
          & \overline {X}_2
        \end {tikzcd}

        Namely, we can take \(\overline {X}\) to be a resolution of singularities of the closure of the diagonal image of \(X\) in \(\overline {X}_1\times \overline {X}_2\).
        The identity map from \(\operatorname {H}^n(X,\mathbb {Z})\) endowed with the mixed Hodge structure defined by \(\overline {X}_1\) to \(\operatorname {H}^n(X,\mathbb {Z})\) endowed with the mixed Hodge structure defined by \(\overline {X}_2\) is thus the composite of two isomorphisms.
    

  \end{enumerate}\par{}To finish the proof of (ii) and (iii), we note that every morphism \(f\) fits into a diagram of the form in \cref{hodge-theory-ii-3.2.11.1}: we choose compactifications \(\overline {X}'_1\) and \(\overline {X}_2\) of \(X_1\) and \(X_2\), then we take \(\overline {X}_1\) to be a resolution of singularities of the closure of the image of \(X_1\) in \(\overline {X}'_1\times \overline {X}_2\).\begin{definition}{3.2.12}\setnumber{3.2.12}\label{hodge-theory-ii-3.2.12}\skipsinglepar\par{}The \emph{mixed Hodge structure of the cohomology of a smooth separated algebraic variety} is the mixed Hodge structure from (iii) of \hyperref[hodge-theory-ii-3.2.5]{Theorem~\ref*{hodge-theory-ii-3.2.5}}.\end{definition}\begin{corollary}{3.2.13}\setnumber{3.2.13}\label{hodge-theory-ii-3.2.13}\skipsinglepar\par{}With the above notation:

    
 \begin{enumerate}
      
 \item[(i)]
        The spectral sequence in \hyperref[hodge-theory-ii-3.2.4.1]{Equation~\ref*{hodge-theory-ii-3.2.4.1}} degenerates at \(E_2\), i.e. the Leray spectral sequence for the inclusion \(j\colon  X^*\hookrightarrow  X\) degenerates at \(E_3\) (i.e. \(E_3=E_\infty \)).
      


      
 \item[(ii)]
        The spectral sequence in \hyperref[hodge-theory-ii-3.2.4.2]{Equation~\ref*{hodge-theory-ii-3.2.4.2}}
        \[           {}_FE_1^{pq}           = \operatorname {H}^q(\overline {X},\Omega _{\overline {X}}^p\langle  Y\rangle )           \Rightarrow  \operatorname {H}^{p+q}(X,\mathbb {C})         \]
        degenerates at \(E_1\).
      


      
 \item[(iii)]
        The spectral sequence defined by the sheaf \(\Omega _X^p\langle  Y\rangle \), endowed with the filtration \(W\)
        \[           \begin {aligned}             E_1^{-n,k+n}             &= \operatorname {H}^k(\overline {X},\operatorname {Gr}_n^W(\Omega _X^p\langle  Y\rangle ))           \\&\approx  \operatorname {H}^k(\widetilde {Y}^n,\Omega _{\widetilde {Y}^n}^{p-n}(\varepsilon ^n))             \Rightarrow  \operatorname {H}^k(\overline {X},\Omega _X^p\langle  Y\rangle )           \end {aligned}         \]
        degenerates at \(E_2\).
      

    \end{enumerate}\begin{proof}\par{}Claim (i) is proven in \hyperref[hodge-theory-ii-3.2.10]{Lemma~\ref*{hodge-theory-ii-3.2.10}}.\par{}Consider the four spectral sequences in \cref{hodge-theory-ii-1.4.2}, with respect to the bifiltered complex \(\Omega _{\overline {X}}^\bullet \langle  Y\rangle \).
      By (i), we have
      \[         \sum _n \dim \operatorname {H}^n(X,\mathbb {C})         = \sum _{p,q} \dim  {}_WE_2^{p,q}.       \]
      The \({}_WE_1^{n,k-n}\) pages are, for fixed \(n\), the abutments of a spectral sequence \emph{that degenerates at \(E_1\)} (\hyperref[hodge-theory-ii-3.2.6]{Lemma~\ref*{hodge-theory-ii-3.2.6}}) with initial pages \(E_1^{p,n,k-p-n}\) (with the notation of \hyperref[hodge-theory-ii-1.4.12]{Lemma~\ref*{hodge-theory-ii-1.4.12}}) the initial pages of the spectral sequence in (iii).
      
      Since \({}_Wd_1\) is strictly compatible with the filtration \(F\) which is the abutment of this spectral sequence (\hyperref[hodge-theory-ii-3.2.7]{Lemma~\ref*{hodge-theory-ii-3.2.7}}), and is (\hyperref[hodge-theory-ii-1.4.12]{Lemma~\ref*{hodge-theory-ii-1.4.12}}) the abutment of a morphism between these spectral sequences, which starts with the differentials from (iii), we have
      \[         \operatorname {Gr}_F^p({}_WE_2^{n,k-n})         \approx  \operatorname {H}^\bullet (E_1^{p,n-1,k-n-p} \to  E_1^{p,n,k-n-p} \to  E_1^{p,n+1,k-n-p})       \]
      and \(\operatorname {Gr}_F^\bullet ({}_WE_2^{\bullet \bullet })\) is the sum of the pages \(E_2^{\bullet \bullet \bullet }\) of the spectral sequences in (iii).
      We thus have

      \begin{eqn}{3.2.13.1}\setnumber{3.2.13.1}\label{hodge-theory-ii-3.2.13.1}\skipsinglepar\[           \sum _n \dim  \operatorname {H}^n(X,\mathbb {C})           = \sum  \dim  E_2^{\bullet \bullet \bullet }.         \tag{3.2.13.1}         \]\end{eqn}

      We also have
      \[         \sum _n \dim  \operatorname {H}^n(X,\mathbb {C})         \leqslant  \dim  {}_FE_1^{\bullet \bullet }       \]
      with equality if and only if the spectral sequence in (ii) degenerates at \(E_1\), and
      \[         \sum  \dim  {}_FE_1^{\bullet \bullet }         \leqslant  \sum  \dim  E_2^{\bullet \bullet \bullet }       \]
      with equality if and only if the spectral sequences in (iii) degenerate at \(E_2\).
      Comparing with \hyperref[hodge-theory-ii-3.2.13.1]{Equation~\ref*{hodge-theory-ii-3.2.13.1}}, we obtain \hyperref[hodge-theory-ii-3.2.13]{Corollary~\ref*{hodge-theory-ii-3.2.13}}.\end{proof}\end{corollary}\begin{corollary}{3.2.14}\setnumber{3.2.14}\label{hodge-theory-ii-3.2.14}\skipsinglepar\par{}Let \(\omega \) be a meromorphic differential \(p\)-form on \(\overline {X}\) that is holomorphic on \(X\) and presents at worst logarithmic poles along \(Y\).
    Then the restriction \(\omega |X\) of \(\omega \) to \(X\) is closed, and if the cohomology class in \(\operatorname {H}^p(X,\mathbb {C})\) defined by \(\omega \) is zero then \(\omega =0\).\begin{proof}\par{}This is the particular case of (ii) of \hyperref[hodge-theory-ii-3.2.13]{Corollary~\ref*{hodge-theory-ii-3.2.13}} where \({}_FE_1^{p0}={}_FE_\infty ^{p0}\).\end{proof}\end{corollary}\begin{corollary}{3.2.15}\setnumber{3.2.15}\label{hodge-theory-ii-3.2.15}\skipsinglepar\begin{enumerate}
    
 \item[(i)]
      If \(X\) is a smooth complete algebraic variety, then the mixed Hodge structure on \(\operatorname {H}^n(X,\mathbb {Z})\) is the classical Hodge structure of weight \(n\).
    


    
 \item[(ii)]
      The Hodge numbers \(h^{pq}\) of the mixed Hodge structure of \(\operatorname {H}^n(X,\mathbb {Z})\) (with \(X\) algebraic and smooth) can only be zero for \(p\leqslant  n\), \(q\leqslant  n\), and \(p+q\geqslant  n\)
    

  \end{enumerate}\begin{proof}\par{}Claim (i) is clear;
      to prove (ii), note that, if \(Y\) is the union of smooth divisors, then the rational Hodge structure \(\operatorname {Gr}_k^W(\operatorname {H}^n(X,\mathbb {Q}))\) is the quotient of a sub-object of a Hodge structure of weight \(n+k\), namely
      \[         \operatorname {H}^{n-k}(\widetilde {Y}^n,\mathbb {Q})\otimes \mathbb {Q}(-k).       \]\end{proof}\end{corollary}\label{hodge-theory-ii-3.2.16}\par{}Let \(X\) be a smooth separated scheme.
    We know that \(X\) admits smooth compactifications \(\overline {X}\) and that the subgroup of \(\operatorname {H}^n(X,\mathbb {Z})\) given by the image of \(\operatorname {H}^n(\overline {X},\mathbb {Z})\) is independent of the choice of \(\overline {X}\) (cf. [\hyperref[G1968]{G1968}, (9.1) to (9.4)]).\begin{corollary}{3.2.17}\setnumber{3.2.17}\label{hodge-theory-ii-3.2.17}\skipsinglepar\par{}Under the hypotheses of \cref{hodge-theory-ii-3.2.16}, the image of \(\operatorname {H}^n(\overline {X},\mathbb {Q})\) in \(\operatorname {H}^n(X,\mathbb {Q})\) is \(W_n(\operatorname {H}^n(X,\mathbb {Q}))\) (where \(W\) denotes the weight filtration (\hyperref[hodge-theory-ii-3.2.12]{Definition~\ref*{hodge-theory-ii-3.2.12}})).\begin{proof}\par{}We can suppose that \(\overline {X}\setminus  X\) is a normal crossing divisor.
      The claim then follows from the fact that \(W[-n]\) is the abutment of the Leray spectral sequence for the inclusion \(j\colon  X\hookrightarrow \overline {X}\).\end{proof}\end{corollary}\begin{corollary}{3.2.18}\setnumber{3.2.18}\label{hodge-theory-ii-3.2.18}\skipsinglepar\par{}Let \(f\) be a morphism from a smooth proper scheme \(Y\) to a smooth scheme \(X\), with \(X\) admitting a smooth compactification \(j\colon  X\hookrightarrow \overline {X}\).
    Then the groups \(\operatorname {H}^n(X,\mathbb {Q})\) and \(\operatorname {H}^n(\overline {X},\mathbb {Q})\) have the same image in \(\operatorname {H}^n(Y,\mathbb {Q})\).\begin{proof}\par{}Since \(f^*\) and \((jf)^*\) are strictly compatible with the weight filtration ((iii) of \hyperref[hodge-theory-ii-3.2.5]{Theorem~\ref*{hodge-theory-ii-3.2.5}}), it suffices to prove that \(\operatorname {Gr}^W(f^*)\) and \(\operatorname {Gr}^W(f^*j^*)\) have the same image in \(\operatorname {Gr}^W(\operatorname {H}^n(Y,\mathbb {Q}))\).
      By \hyperref[hodge-theory-ii-3.2.17]{Corollary~\ref*{hodge-theory-ii-3.2.17}} and \hyperref[hodge-theory-ii-3.2.15]{Corollary~\ref*{hodge-theory-ii-3.2.15}}, \(\operatorname {Gr}_n^W(j^*)\) is an isomorphism, while \(\operatorname {Gr}_m^W(f^*)=0\) for \(m\neq  n\) since \(\operatorname {Gr}_m^W(\operatorname {H}^n(Y,\mathbb {Q}))=0\) for \(m\neq  n\).\end{proof}\end{corollary}\begin{remark}{3.2.19}\setnumber{3.2.19}\label{hodge-theory-ii-3.2.19}\skipsinglepar\par{}By \hyperref[hodge-theory-ii-3.1.11]{Proposition~\ref*{hodge-theory-ii-3.1.11}} and \cref{hodge-theory-ii-1.4.5}, under the hypotheses of \hyperref[hodge-theory-ii-3.2.5]{Theorem~\ref*{hodge-theory-ii-3.2.5}}, the Hodge filtration on \(\operatorname {H}^n(X,\mathbb {C})\) is the abutment of the hypercohomology spectral sequence of the complex \(j_*^m\Omega _{X^*}^\bullet \) endowed with the filtration given by the pole-order filtration (\cref{hodge-theory-ii-3.1.10}): this spectral sequence coincides with \hyperref[hodge-theory-ii-3.2.4.2]{Equation~\ref*{hodge-theory-ii-3.2.4.2}}.\end{remark}\chapter{Applications and supplements}\thispagestyle{fancy}\label{hodge-theory-ii-4}\section{The fixed set theorem}\label{hodge-theory-ii-4.1}\begin{theorem}{4.1.1}\setnumber{4.1.1}\label{hodge-theory-ii-4.1.1}\skipsinglepar\par{}Let \(S\) be a smooth separated scheme, and \(f\colon  X\to  S\) a smooth proper morphism.\begin{enumerate}
    
 \item[(i)]
      In rational cohomology, the Leray spectral sequence
      \[         E_2^{pq}         = \operatorname {H}^p(S,\mathrm {R}^qf_*\mathbb {Q})         \Rightarrow  \operatorname {H}^{p+q}(X,\mathbb {Q})       \]
      degenerates (\(E_2=E_\infty \)).
    


    
 \item[(ii)]
      If \(\overline {X}\) is a non-singular compactification of \(X\), then the canonical morphism
      \[         \operatorname {H}^n(\overline {X},\mathbb {Q})         \to  \operatorname {H}^0(S,\mathrm {R}^nf_*\mathbb {Q})       \]
      is surjective.
    

  \end{enumerate}\begin{proof}\par{}For \(f\) smooth and projective, claim (i) is proven in \hyperref[D1968]{D1968}.
      The proof in \hyperref[D1968]{D1968} is rewritten in a more readable manner in \hyperref[G1970]{G1970};
      it does not use the smoothness of \(S\).\par{}Let us fix \(f\), and prove that (i)\(\implies \)(ii).
      We can reduce to the case where \(S\) is non-empty and connected.
      If \(s\in  S\) is a point, then the local system \(\mathrm {R}^nf_*\mathbb {Q}\) is entirely described by its fibre \((\mathrm {R}^nf_*\mathbb {Q})_s\) at \(s\) and the action of the fundamental group \(\pi =\pi _1(S,s)\) on this fibre.
      We have
      \[         \operatorname {H}^0(S,\mathrm {R}^n f_*\mathbb {Q})         \xrightarrow {\sim } [(\mathrm {R}^n f_*\mathbb {Q})_s]^\pi .       \]
      If \(X_s=f^{-1}(s)\), then the composite arrow
      \[         \operatorname {H}^0(S,\mathrm {R}^n f_*\mathbb {Q})         \to  (\mathrm {R}^n f_*\mathbb {Q})         \approx  \operatorname {H}^n(X_s,\mathbb {Q})       \]
      is thus injective.
      Consider the arrows
      \[         \operatorname {H}^n(\overline {X},\mathbb {Q})         \xrightarrow {a} \operatorname {H}^n(X,\mathbb {Q})         \xrightarrow {b} \operatorname {H}^0(S,\mathrm {R}^n f_*\mathbb {Q})         \xhookrightarrow {c} \operatorname {H}^n(X_s,\mathbb {Q}).       \]
      The arrows \(cb\) and \(cba\) have the same image, by \hyperref[hodge-theory-ii-3.2.18]{Corollary~\ref*{hodge-theory-ii-3.2.18}}.
      The arrow \(b\) is an "edge homomorphism" of the Leray spectral sequence, and is thus surjective by hypothesis.
      Since \(c\) is injective, \(b\) and \(ba\) have the same image, and \(ba\) is surjective.\par{}This proves (ii) for \(f\) projective.
      From this we will deduce the general case.
      We can again reduce to the case where \(S\) is non-empty and connected.\par{}By Chow's lemma and the resolution of singularities, there exists a quasi-projective smooth scheme \(X'\) and a projective birational morphism \(p\colon  X'\to  X\).
      
      There thus exists (by Bertini, or Sard) a non-empty Zariski open \(S_1\) of \(S\) such that \(X'_1=(fp)^{-1}(S_1)\) is smooth over \(S_1\).
      Finally, let \(X_1=f^{-1}(S_1)\), let \(\overline {X}\) be a smooth compactification of \(X\), and let \(\overline {X}'_1\) be a smooth compactification of \(X'_1\) that dominates \(\overline {X}\).

      
  \begin {tikzcd}
        \overline {X}
        && \overline {X}'_1
          \ar [ll,swap,"\overline {p}"]
      \\X
          \ar [u,hook]
          \ar [d,swap,"f"]
        & X_1
          \ar [l,hook']
          \ar [d,swap,"f_1"]
        & X'_1
          \ar [u,hook]
          \ar [l,swap,"p"]
          \ar [d,swap,"f'_1"]
      \\S
        & S_1
          \ar [l,hook',"i"]
        & S_1
          \ar [l,equals]
      \end {tikzcd}\par{}If \(s\in  S_1\), then the morphism \(i_*\colon \pi _1(S_1,s)\to \pi _1(S,s)\) is surjective, since the topological codimension of \(S\setminus  S_1\) is \(\geqslant 2\).
      We thus have
      \[         \operatorname {H}^0(S,\mathrm {R}^n f_*\mathbb {Q})         \xrightarrow {\sim } \operatorname {H}^0(S_1,\mathrm {R}^n {f_1}_*\mathbb {Q})       \]
      and it suffices to prove that
      \[         \operatorname {H}^n(\overline {X},\mathbb {Q})         \to  \operatorname {H}^0(S_1,\mathrm {R}^n {f_1}_*\mathbb {Q})       \]
      is surjective.\par{}The vertical arrows in the diagram
      \[         \begin {CD}           \operatorname {H}^n(\overline {X},\mathbb {Q}) @>>> \operatorname {H}^n(X_1,\mathbb {Q}) @>>> \operatorname {H}^0(S_1,\mathrm {R}^n{f_1}_*\mathbb {Q})         \\@VV{\overline {p}^*}V @VV{p^*}V @VV{p^*}V         \\\operatorname {H}^n(\overline {X}'_1,\mathbb {Q}) @>>> \operatorname {H}^n(X'_1,\mathbb {Q}) @>>> \operatorname {H}^0(S_1,\mathrm {R}{f'_1}_*\mathbb {Q})         \end {CD}       \]
      admit left inverses given by the \emph{Gysin morphisms} \(\overline {p}_!\), \(p_!\), and \(p_!\) (respectively), defined by Poincaré duality as the transposes of the arrows analogous to the above but in cohomology with proper support.
      Furthermore, the diagram
      \[         \begin {CD}           \operatorname {H}^n(\overline {X},\mathbb {Q}) @>{u}>> \operatorname {H}^0(S_1,\mathrm {R}^n{f_1}_*\mathbb {Q})         \\@A{\overline {p}_!}AA @AA{p_!}A         \\\operatorname {H}^n(\overline {X}',\mathbb {Q}) @>>{v}> \operatorname {H}^0(S_1,\mathrm {R}^n{f'_1}_*\mathbb {Q})         \end {CD}       \]
      commutes.
      The arrow \(u\) is thus a direct factor of the arrow \(v\).
      Since \(v\) is surjective (because \(f'_1\) is projective), so too is \(u\).\par{}We now prove (i) in the general case.
      We can reduce to the case where \(f\) is of pure relative dimension \(n\).
      Let \(f\times  f\) be the projection of \(X\times _S X\) to \(S\).\par{}Let \(\delta \) be the image of the cohomology class of the diagonal of \(X\times _S X\) in \(\operatorname {H}^0(S,\mathrm {R}^n(f\times  f)_*,\mathbb {Q})\).
      We have, by Künneth,
      \[         \mathrm {R}^n(f\times  f)_*\mathbb {Q}         = \sum _{p+q=n} (\mathrm {R}^p f_*\mathbb {Q}\otimes \mathrm {R}^q f_*\mathbb {Q}).       \]
      We denote by \(\delta '_{pq}\) (for \(p+q=n\)) the components of \(\delta \) in this decomposition, and by \(\delta _{pq}\) the classes in \(\operatorname {H}^n(X\times _S X)\) given by the images of the \(\delta '_{pq}\).\par{}The \(\delta _{pq}\) define, in the derived category \(\mathrm {D}^+(S)\), homomorphisms
      \[         \delta _p\colon          \mathrm {R} f_*\mathbb {Q}         \to  \mathrm {R} f_*\mathbb {Q}       \]
      such that \(\operatorname {\mathscr {H}}^q(\delta _p)\) is zero for \(p\neq  q\), and is the identity for \(p=q\).
      By [\hyperref[D1968]{D1968}], we thus have

      \begin{eqn}{4.1.1.1}\setnumber{4.1.1.1}\label{hodge-theory-ii-4.1.1.1}\skipsinglepar\[           \mathrm {R} f_*\mathbb {Q}           \approx  \sum _p \mathrm {R}^p f_*\mathbb {Q}[-p]         \tag{4.1.1.1}         \]\end{eqn}

      in \(\mathrm {D}^+(S)\), and the Leray spectral sequence degenerates.\end{proof}\end{theorem}\begin{corollary}{4.1.2}\setnumber{4.1.2}\label{hodge-theory-ii-4.1.2}\skipsinglepar\par{}Let \(f\colon  X\to  S\) be a proper smooth morphism, with \(S\) reduced, connected, and separated.
    Let \((\mathrm {R}^n f_*\mathbb {Q})^0\) be the largest constant local system on \(\mathrm {R}^n f_*\mathbb {Q}\), with fibre \(\operatorname {H}^0(S,\mathrm {R}^n f_*\mathbb {Q})\).
    Then, for each \(s\in  S\), \((\mathrm {R}^n f_*\mathbb {Q})_s^0\) is a Hodge sub-structure of \((\mathrm {R}^n f_*\mathbb {Q})\simeq \operatorname {H}^n(X_s,\mathbb {Q})\), and the induced Hodge structure on \(\operatorname {H}^0(S,\mathrm {R}^n f_*\mathbb {Q})\) is independent of \(s\).\begin{proof}\par{}Let \(s\in  S\), and let \(X_s=f^{-1}(s)\).
      If \(S\) is smooth, and if \(\overline {X}\) is a smooth compactification of \(X\), then, by \hyperref[hodge-theory-ii-4.1.1]{Theorem~\ref*{hodge-theory-ii-4.1.1}}, the subspace \((\mathrm {R}^n f_*\mathbb {Q})_s^0\) of \((\mathrm {R}^n f_*\mathbb {Q})_s\) is the image of \(\operatorname {H}^n(\overline {X},\mathbb {Q})\).
      Since the restriction map
      \[         \operatorname {H}^n(\overline {X},\mathbb {Q}) \to  \operatorname {H}^n(X_s,\mathbb {Q})       \]
      is a morphism of Hodge structures, its image is a Hodge sub-structure, and the induced Hodge structure on \(\operatorname {H}^0(S,\mathrm {R}^n f_*\mathbb {Q})\), given by the quotient of that on \(\operatorname {H}^n(\overline {X},\mathbb {Q})\), is independent of \(s\).\end{proof}\end{corollary}\par{}In the general case, \hyperref[hodge-theory-ii-4.1.2]{Corollary~\ref*{hodge-theory-ii-4.1.2}} also implies that, if \(a\) is a global section of \(\mathrm {R}^n f_*\mathbb {C}\), then its components \(a^{p,q}\) of degree \((p,q)\), which are a priori only continuous sections of the complex bundle defined by \(\mathrm {R}^n f_*\mathbb {C}\), are in fact locally constant, i.e. are sections of \(\mathrm {R}^n f_*\mathbb {C}\).
  This is the case because \(a^{p,q}\) is continuous, and further locally constant on the (dense) open subset of smooth points of \(S\), by the above.\label{hodge-theory-ii-4.1.3}\par{}In the case where \(S\) is compact, a generalisation of \hyperref[hodge-theory-ii-4.1.2]{Corollary~\ref*{hodge-theory-ii-4.1.2}} is proven by analytic methods in [\hyperref[G1970]{G1970}].\par{}As a corollary to \hyperref[hodge-theory-ii-4.1.2]{Corollary~\ref*{hodge-theory-ii-4.1.2}}, we note:

    \label{hodge-theory-ii-4.1.3.1}\par{}(cf. [\hyperref[G1970]{G1970}]).
        Let \(f\colon  X\to  S\) be a proper smooth morphism, as in \hyperref[hodge-theory-ii-4.1.2]{Corollary~\ref*{hodge-theory-ii-4.1.2}}.
        If a global section \(a\) of \(\mathrm {R}^n f_*\mathbb {C}\) is of Hodge degree \((p,q)\) at a point, then \(a\) is everywhere of degree \((p,q)\).
        In particular, if \(n=2\) and if \(a\) is, at some point \(s\), the cohomology class of a divisor of \(X_s\), then \(a\) is everywhere the cohomology class of a divisor, and, if \(S\) is smooth, is even defined by a divisor \(D\) on \(X\).

    \label{hodge-theory-ii-4.1.3.2}\par{}(cf. [\hyperref[G1966]{G1966}]).
        Let \(S\) be a reduced connected scheme of finite type over \(\mathbb {C}\), and let \(f_1\colon  X_1\to  S\) and \(f_2\colon  X_2\to  S\) be abelian schemes over \(S\).
        If a morphism \(u\colon \mathrm {R}^1{f_2}_*\mathbb {Z}\to \mathrm {R}^1{f_1}_*\mathbb {Z}\) comes from a morphism \(\widetilde {u}_s\colon (X_1)_s\to (X_2)_s\) of abelian varieties at some point \(s\in  S\), then \(u\) comes from a (unique) morphism of abelian schemes \(\widetilde {u}\colon  X_1\to  X_2\).

    \label{hodge-theory-ii-4.1.3.3}\par{}(cf. [\hyperref[K1971]{K1971}]).
        Let \(S\) be a smooth connected scheme, \(s\in  S\) a point, \(f\colon  X\to  S\) a proper smooth morphism, and \(P\) a direct factor of the local system \(\mathrm {R}^i f_*\mathbb {Q}\), which is pointwise a Hodge sub-structure.
        Then the following conditions are equivalent:

        
 \begin{enumerate}
          
 \item[(a)]
            The Hodge structure of \(P\) is locally constant.
          

          
 \item[(b)]
            The representation \(P_s\) of \(\pi _1(S,s)\) factors through a finite quotient of \(\pi _1(S,s)\).
          

          
 \item[(c)]
            There exists a finite non-empty étale covering \(u\colon  S'\to  S\) such that \(u^*P\) is a constant family of Hodge structures.
          

        \end{enumerate}\section{The semi-simplicity theorem}\label{hodge-theory-ii-4.2}\label{hodge-theory-ii-4.2.1}\par{}Let \(S\) be a topological space.
    A \emph{continuous family of Hodge structures} on \(S\) consists of:

    
 \begin{enumerate}
      
 \item[(a)]
        A local system \(H_\mathbb {Z}\) on \(S\) of \(\mathbb {Z}\)-modules of finite type.
      

      
 \item[(b)]
        For every point \(s\in  S\), a Hodge structure on the fibre \((H_\mathbb {Z})\) that varies continuously in \(s\).
      

    \end{enumerate}\par{}A continuous family \(H\) of Hodge structures on \(S\) is said to be \emph{of weight \(n\)} if the fibres \(H_s\) (for \(s\in  S\)) are all of weight \(n\).\par{}We similarly define a \emph{continuous family of Hodge \(\mathbb {Q}\)-structures} as a local system of \(\mathbb {Q}\)-vector spaces, endowed at each point with a Hodge \(\mathbb {Q}\)-structure, varying continuously.\par{}A \emph{polarisation} of a continuous family \(H\) of Hodge \(\mathbb {Q}\)-structures of weight \(n\) is a morphism of local systems from \(H_\mathbb {Q}\otimes  H_\mathbb {Q}\) to the constant local system \(\mathbb {Q}(-n)_\mathbb {Q}\), which defines at each point \(s\in  S\) a polarisation of \(H_s\).\label{hodge-theory-ii-4.2.2}\par{}Suppose that \(S\) is connected, and let \(\mathcal {C}\) be a strictly full subcategory of the category of continuous families of Hodge \(\mathbb {Q}\)-structures on \(S\).
    We will have need of the following conditions:

    \begin{condition}{4.2.2.1}\setnumber{4.2.2.1}\label{hodge-theory-ii-4.2.2.1}\skipsinglepar\par{}\(\mathcal {C}\) is stable under direct factors, direct sums, and tensor products;
        the Tate constant families \(\mathbb {Q}(n)\) (for \(n\in \mathbb {Z}\)) are in \(\mathcal {C}\).\end{condition}

    \begin{condition}{4.2.2.2}\setnumber{4.2.2.2}\label{hodge-theory-ii-4.2.2.2}\skipsinglepar\par{}Every homogeneous (of some weight \(n\)) Hodge structure in \(\mathcal {C}\) is polarisable.\end{condition}

    \begin{condition}{4.2.2.3}\setnumber{4.2.2.3}\label{hodge-theory-ii-4.2.2.3}\skipsinglepar\par{}For every \(H\in \operatorname {Ob}\mathcal {C}\), there exists a local system \(H'_\mathbb {Z}\) on \(S\) of free \(\mathbb {Z}\)-modules such that \(H'_\mathbb {Z}\otimes \mathbb {Q}\approx  H_\mathbb {Q}\).\end{condition}

    \begin{condition}{4.2.2.4}\setnumber{4.2.2.4}\label{hodge-theory-ii-4.2.2.4}\skipsinglepar\par{}For every \(H\in \operatorname {Ob}\mathcal {C}\), the largest constant local system \(H^f\) of \(H\) is a constant family of Hodge sub-structures of \(H\).\end{condition}\begin{lemma}{4.2.3}\setnumber{4.2.3}\label{hodge-theory-ii-4.2.3}\skipsinglepar\par{}If \(\mathcal {C}\) satisfies \hyperref[hodge-theory-ii-4.2.2.1]{Condition~\ref*{hodge-theory-ii-4.2.2.1}} and \hyperref[hodge-theory-ii-4.2.2.2]{Condition~\ref*{hodge-theory-ii-4.2.2.2}}, then:

    
 \begin{enumerate}
      
 \item[(i)]
        \(\mathcal {C}\) is a semi-simple abelian subcategory of the abelian subcategory of continuous families of Hodge \(\mathbb {Q}\)-structures on \(S\).
      

      
 \item[(ii)]
        If \(H\in \operatorname {Ob}\mathcal {C}\), then its dual \(H^*\) and the \(\wedge ^p H\) are also in \(\mathcal {C}\);
        if \(H_1,H_2\in \operatorname {Ob}\mathcal {C}\), then \(\operatorname {Hom}(H_1,H_2)\in \operatorname {Ob}\mathcal {C}\).
      

    \end{enumerate}\begin{proof}\par{}Let \(H\in \operatorname {Ob}\mathcal {C}\), and let \(H_1\) be a sub-object of \(H\), in the category of continuous families of Hodge \(\mathbb {Q}\)-structures.
      We will show that \(H_1\) is a direct factor of \(H\) in this category.
      We can suppose \(H\) to be homogeneous.
      If \(\psi \) is a polarisation form for \(H\), then the orthogonal of \(H_1\) with respect to \(\psi \) is indeed a sub-object of \(H\) that is a supplement to \(H_1\).
      This proves (i)\par{}If \(H\in \operatorname {Ob}\mathcal {C}\) is of weight \(n\), then a polarisation of \(H\) defines an isomorphism between \(H^*\) and \(H\otimes \mathbb {Q}(n)\).
      By \hyperref[hodge-theory-ii-4.2.2.1]{Condition~\ref*{hodge-theory-ii-4.2.2.1}}, we thus have that \(H^*\in \operatorname {Ob}\mathcal {C}\).
      For arbitrary \(H\in \operatorname {Ob}\mathcal {C}\), if we decompose \(H\) into its homogeneous components, then \(H^*=\bigoplus _n(H^n)^*\), whence again \(H^*\in \operatorname {Ob}\mathcal {C}\).
      Finally, \(\wedge ^p H\) is a direct factor of \(\bigotimes ^p H\), and \(\operatorname {Hom}(H_1,H_2)\simeq  H_1^*\otimes  H_2\).\end{proof}\end{lemma}\par{}Let \(f\colon  X\to  S\) be a smooth proper morphism of reduced schemes.
  The sheaf \(\mathrm {R}^i f_*\mathbb {Q}\) is then a local system, and, for \(s\in  S\), \((\mathrm {R}^i f_*\mathbb {Q})_s\simeq \operatorname {H}^i(X_s,\mathbb {Q})\) is endowed with a Hodge \(\mathbb {Q}\)-structure, which varies continuously in \(s\).\begin{definition}{4.2.4}\setnumber{4.2.4}\label{hodge-theory-ii-4.2.4}\skipsinglepar\par{}Let \(S\) be a smooth connected scheme.
    A continuous family \(H\) of Hodge \(\mathbb {Q}\)-structures on \(S^\mathrm {an}\) is said to be \emph{algebraic} if there exists a non-empty Zariski open \(U\) of \(S\), an integer \(k\), and a smooth projective morphism \(f\colon  X\to  U\) such that \(H|U\) is a direct factor of \(\mathrm {R} f_*\mathbb {Q}\otimes \mathbb {Q}(k)\).\end{definition}\begin{proposition}{4.2.5}\setnumber{4.2.5}\label{hodge-theory-ii-4.2.5}\skipsinglepar\begin{enumerate}
    
 \item[(i)]
      The category of algebraic continuous families of Hodge \(\mathbb {Q}\)-structures on \(S\) satisfies the conditions of \cref{hodge-theory-ii-4.2.2}.
    

    
 \item[(ii)]
      If a continuous family \(H\) of Hodge structures on \(S\) is such that its restriction to a dense Zariski open \(U\) of \(S\) is algebraic, then \(H\) is algebraic.
    

    
 \item[(iii)]
      If \(f\colon  X\to  S\) is smooth and proper, then \(\mathrm {R} f_*\mathbb {Q}\) is algebraic.
    

  \end{enumerate}\begin{proof}\par{}Claim (ii) is evident.
      We will prove (i).
      Let \(\mathcal {C}_0\) be the set of continuous families of Hodge structures on \(S\) that are of the form \(\mathrm {R} f_*\mathbb {Q}\) (with \(f\) smooth and projective).
      Then:

      
 \begin{enumerate}
        
 \item[(a)]
          By the Künneth formula
          \[             \mathrm {R}(f\times  g)_*\mathbb {Q}             \simeq  \mathrm {R} f_*\mathbb {Q}\otimes \mathrm {R} g_*\mathbb {Q},           \]
          \(\mathcal {C}_0\) is stable under tensor products.
        

        
        
 \item[(b)]
          Since
          \[             \mathrm {R}(f\sqcup  g)_*\mathbb {Q}             \simeq  \mathrm {R} f_*\mathbb {Q}\oplus \mathrm {R} g_*\mathbb {Q},           \]
          \(\mathcal {C}_0\) is stable under direct sums.
        


        
 \item[(c)]
          
          The homogeneous components of each \(H\in \operatorname {Ob}\mathcal {C}_0\) are polarisable (cf. \cref{hodge-theory-ii-2.2.6}).
        


        
 \item[(d)]
          The objects \(H\in \operatorname {Ob}\mathcal {C}_0\) satisfy \hyperref[hodge-theory-ii-4.2.2.3]{Condition~\ref*{hodge-theory-ii-4.2.2.3}}, since
          \[             \mathrm {R} f_*\mathbb {Q}             \simeq  \mathrm {R} f_*\mathbb {Z}\otimes \mathbb {Q}.           \]
        


        
 \item[(e)]
          The objects \(H\in \operatorname {Ob}\mathcal {C}_0\) satisfy \hyperref[hodge-theory-ii-4.2.2.4]{Condition~\ref*{hodge-theory-ii-4.2.2.4}}, by \hyperref[hodge-theory-ii-4.1.2]{Corollary~\ref*{hodge-theory-ii-4.1.2}}.
        

      \end{enumerate}\par{}Let \(\mathcal {C}_1\) be the set of direct factors of objects of \(\mathcal {C}_0\).
      Then \(\mathcal {C}_1\) is stable under tensor products, direct sums, and direct factors, and satisfies \hyperref[hodge-theory-ii-4.2.2.2]{Condition~\ref*{hodge-theory-ii-4.2.2.2}}, \hyperref[hodge-theory-ii-4.2.2.3]{Condition~\ref*{hodge-theory-ii-4.2.2.3}}, and \hyperref[hodge-theory-ii-4.2.2.4]{Condition~\ref*{hodge-theory-ii-4.2.2.4}}.
      Furthermore, \(\mathbb {Q}(-1)\) is in \(\mathcal {C}_1\), in the form of \(\mathrm {R}^2 f_*\mathbb {Q}\) for \(f\colon \mathbb {P}_S^1\to  S\) the projective bundle.
      The category \(\mathcal {C}_2\) consisting of the \(H\otimes \mathbb {Q}(k)\) (for \(k\in \mathbb {N}\)) thus satisfy all the conditions of \cref{hodge-theory-ii-4.2.2}.\par{}To thus deduce (i), it suffices to note that, if \(H\) is a continuous family of Hodge structures, and \(U\) a dense Zariski open of \(S\), then a sub-object of, a polarisation on, or an integer lattice of \(H|U\) uniquely extend to \(H\).\par{}We now prove (iii).
      If \(f\colon  X\to  S\) is smooth and projective, then there exists a dense Zariski open subset \(U\) of \(S\) along with a commutative diagram
      
  \begin {tikzcd}
        X' \ar [rr,"p"] \ar [dr,swap,"f'"]
        && X|U \ar [dl,"f|U"]
      \\& U
      \end {tikzcd}

      with \(f'\) smooth and projective, and \(p\) birational and surjective (by applying Chow's lemma and the resolution of singularities to the generic fibre of \(f\)).
      The restriction map
      \[         p^*\colon  (\mathrm {R} f_*\mathbb {Q})|U \to  \mathrm {R} f'_*\mathbb {Q}       \]
      is thus a direct injection of continuous families of Hodge structures, with left inverse given by the Gysin morphism \(p_!\), whence the algebraicity of \(\mathrm {R} f_*\mathbb {Q}\).\end{proof}\end{proposition}\begin{theorem}{4.2.6}\setnumber{4.2.6}\label{hodge-theory-ii-4.2.6}\skipsinglepar\par{}Let \(S\) be a connected topological space that is locally connected and locally simply connected, endowed with a basepoint \(s\).
    Let \(\mathcal {C}\) be a category of continuous Hodge structures on \(S\) that satisfy the conditions of \cref{hodge-theory-ii-4.2.2}.
    Then, if \(H\in \operatorname {Ob}\mathcal {C}\), the representation of \(\pi _1(S,s)\) on the fibre \((H_\mathbb {Q})_s\) is semi-simple.\end{theorem}\begin{footnote}{1}\setnumber{1}\label{hodge-theory-ii-4.2.6-footnote-1}\skipsinglepar\par{}Let \(S\) be a smooth connected scheme.
    A \emph{family of Hodge structures} on \(S\) is a continuous family \(H\) of Hodge structures on \(S\) that satisfies the following conditions:

    
 \begin{enumerate}
      
 \item[(a)]
        The Hodge filtration on \((H_\mathbb {C})_s\) varies holomorphically with \(s\), i.e. it corresponds to a filtration \(F\) of \(H_\mathcal {O}=H_\mathbb {Z}\otimes \mathcal {O}\).
      

      
 \item[(b)]
        The covariant derivative \(\nabla \) satisfies \(\nabla  F^p\subset \Omega _S^1\otimes  F^{p-1}\).
      

    \end{enumerate}\par{}Let \(\mathcal {C}\) be the category of such continuous families of Hodge \(\mathbb {Q}\)-structures on \(S\) that underlie a family of Hodge structures and whose homogeneous direct factors are polarisable.
    It is clear that \(\mathcal {C}\) satisfies \hyperref[hodge-theory-ii-4.2.2.1]{Condition~\ref*{hodge-theory-ii-4.2.2.1}}, \hyperref[hodge-theory-ii-4.2.2.2]{Condition~\ref*{hodge-theory-ii-4.2.2.2}}, and \hyperref[hodge-theory-ii-4.2.2.3]{Condition~\ref*{hodge-theory-ii-4.2.2.3}}.
    We can deduce from results of W. Schmid (August 1970, unpublished) and P.A. Griffiths [\hyperref[G1970]{G1970}] that \(\mathcal {C}\) satisfies \hyperref[hodge-theory-ii-4.2.2.4]{Condition~\ref*{hodge-theory-ii-4.2.2.4}}.
    This theorem, of which \hyperref[hodge-theory-ii-4.1.2]{Corollary~\ref*{hodge-theory-ii-4.1.2}} is a corollary, allows us to apply \hyperref[hodge-theory-ii-4.2.6]{Theorem~\ref*{hodge-theory-ii-4.2.6}} and its corollaries to the objects of \(\mathcal {C}\).\end{footnote}\par{}Let \(H\in \operatorname {Ob}\mathcal {C}\).
  The local system \(H_\mathbb {Q}\) defines a complex local system \(H_\mathbb {C}=H_\mathbb {Q}\otimes \mathbb {C}\).
  For every complex local system \(V\) on \(S\), we denote by \(V^\mathrm {c}\) the complex vector bundle that it defines, identified with the sheaf of its continuous sections.\par{}By definition, the group \(S\) (from \cref{hodge-theory-ii-2.1.2}) acts on \(H_\mathbb {C}^\mathrm {c}\).
  A sub-bundle of \(H_\mathbb {C}^\mathrm {c}\) is said to be \emph{horizontal}, or \emph{locally constant}, if it is defined by a local subsystem of \(H_\mathbb {C}\).\begin{lemma}{4.2.7}\setnumber{4.2.7}\label{hodge-theory-ii-4.2.7}\skipsinglepar\par{}Let \(V\) be a rank-1 local subsystem of \(H_\mathbb {C}\).
    Suppose that the \(n\)-th tensor power \(V^{\otimes  n}\) (for some \(n\geqslant 1\)) of \(V\) is a trivial local system, i.e. that \(\pi _1(S,s)\) acts on \(V_s\) via a finite (necessarily cyclic) group.
    Then, for all \(t\in  S\), \(tV^\mathrm {c}\) is again locally constant.\begin{proof}\par{}In order for \(tV^\mathrm {c}\) to be locally constant, it suffices for \((tV^\mathrm {c})^{\otimes  n}={tV^{c}}^{\otimes  n}\subset \bigotimes ^n H_\mathbb {C}\) to be locally constant.
      But \(V^{\otimes  n}\) is generated by a horizontal global section \(v\), and, by hypothesis, \(tv\) is again horizontal \hyperref[hodge-theory-ii-4.2.2.4]{Condition~\ref*{hodge-theory-ii-4.2.2.4}}.\par{}We proceed by induction on \(\dim (H_\mathbb {Q})_s\).
      We can suppose \(H\) to be homogeneous and non-zero.
      Let \(d\) be the minimal dimension of the non-zero complex local subsystems of \(H_\mathbb {C}\).
      The sum \(W\) of all the local subsystems of \(H_\mathbb {C}\) of dimension \(d\) (which are automatically simple) is "defined over \(\mathbb {Q}\)", i.e. is of the form \(W_\mathbb {Q}\otimes \mathbb {C}\) for \(W_\mathbb {Q}\) a local subsystem of \(H_\mathbb {Q}\).
      By construction, \(W_s\) is a semi-simple complex \(\pi _1(S,s)\)-module, so \((W_\mathbb {Q})_s\) is a semi-simple \(\pi _1(S,s)\)-module over \(\mathbb {Q}\).\par{}Let \(H_\mathbb {Z}\) be a local system of free \(\mathbb {Z}\)-modules such that \(H_\mathbb {Z}\otimes \mathbb {Q}\simeq  H_\mathbb {Q}\), and let \(W_\mathbb {Z}=H_\mathbb {Z}\cap  W_\mathbb {Q}\).
      If \(W\) is of dimension \(e\), then the local system \((\wedge ^e W)^{\otimes 2}\) is trivial.
      Indeed:
      \[         \bigwedge ^e W         \simeq  \bigwedge ^e W_\mathbb {Z}\otimes \mathbb {C}       \]
      and \(\pi _1(S,s)\) can only act on \((\wedge ^e W_\mathbb {Z})\) by \(\pm 1\).\par{}Let \(V\) be a complex local subsystem of dimension \(d\) of \(H_\mathbb {C}\), and let \(V'\) be a complement of \(V\) inside \(W\) (by the semi-simplicity of \(W\)).
      We have
      \[         \bigwedge ^e W         \simeq  \bigwedge ^d V\otimes \bigwedge ^{e-d}V'.       \]
      Applying \hyperref[hodge-theory-ii-4.2.7]{Lemma~\ref*{hodge-theory-ii-4.2.7}} to the local subsystem \(\wedge ^d V\otimes \wedge ^{e-d}V'\) of \(\wedge ^d H_\mathbb {C}\otimes \wedge ^{e-d}H_\mathbb {C}\), we see that, for \(t\in  S\),
      \[         t\left (\bigwedge ^d V\otimes \bigwedge ^{e-d}V'\right )^\mathrm {c}         = \bigwedge ^d tV^\mathrm {c}\otimes \bigwedge ^{e-d}t{V'}^\mathrm {c}         \subset  \bigwedge ^d H_\mathbb {C}\otimes \bigwedge ^{e-d}H_\mathbb {C}       \]
      is locally constant.
      From this, \(\wedge ^d tV^\mathrm {c}\subset \wedge ^d H_\mathbb {C}\) and \(tV^\mathrm {c}\subset  H_\mathbb {C}\) are locally constant.\par{}By the definition of \(W\), we have \(tV^\mathrm {c}\subset  W^\mathrm {c}\), since \(W^\mathrm {c}\) is stable under \(S\), and \(W_\mathbb {Q}\) is a Hodge sub-structure of \(H\).
      If \(\psi \) is a polarisation of \(H\), then \(H\) is thus the direct sum of \(W\) and its orthogonal, which are both in \(\mathcal {C}\), and we conclude by induction.\end{proof}\end{lemma}\begin{corollary}{4.2.8}\setnumber{4.2.8}\label{hodge-theory-ii-4.2.8}\skipsinglepar\par{}Under the hypotheses of \hyperref[hodge-theory-ii-4.2.6]{Theorem~\ref*{hodge-theory-ii-4.2.6}}:\begin{enumerate}
    
 \item[(i)]
      The algebra \(A\) of endomorphisms of the local system \(H_\mathbb {Q}\) is semi-simple.
      It admits exactly one Hodge \(\mathbb {Q}\)-structure such that, at each point \(s\), the map \(A\otimes  H_s\to  H_s\) is a morphism of Hodge structures.
    


    
 \item[(ii)]
      
      The centre of \(A\) is of degree \((0,0)\), and the composition law \(\cdot \colon  A\otimes  A\to  A\) is a morphism of Hodge structures.
    


    
 \item[(iii)]
      Let \(W\) be a complex local subsystem of dimension \(d\) of \(H_\mathbb {C}\).
      Then:

      
 \begin{enumerate}
        
 \item[(a)]
          For \(t\in  S\), \(tW^\mathrm {c}\subset  H_\mathbb {C}\) is locally constant, and defines a local subsystem \(tW\) which is isomorphic to \(W\).
        


        
 \item[(b)]
          A power \((\wedge ^d W)^{\otimes  n}\) (for some \(n\geqslant 1\)) of the local system \(\wedge ^d W\) is a trivial local system.
        

      \end{enumerate}

    

  \end{enumerate}\begin{proof}\par{}The algebra \(A\) is the commutator of \(\pi _1(S,s)\) in \((H_\mathbb {Q})_s\).
      Its semi-simplicity thus follows from \hyperref[hodge-theory-ii-4.2.6]{Theorem~\ref*{hodge-theory-ii-4.2.6}} and from Bourbaki, \emph{Algèbre commutative}, Chapter 8, §5.2, Proposition 4.
      The algebra \(A\) is also the fibre of the largest constant local subsystem of \(\underline {\operatorname {Hom}}(H_\mathbb {Q},H_\mathbb {Q})\).
      Since \(\underline {\operatorname {Hom}}(H,H)\in \mathbb {C}\) (by (ii) of \hyperref[hodge-theory-ii-4.2.3]{Lemma~\ref*{hodge-theory-ii-4.2.3}}), the existence of the Hodge structure in (i) follows from the hypotheses (\hyperref[hodge-theory-ii-4.2.2.4]{Condition~\ref*{hodge-theory-ii-4.2.2.4}}).
      It is clear that the composition law is a morphism of Hodge structures.
      It thus follows that the centre \(Z\) of \(A\) is a Hodge substructure: \(Z_\mathbb {C}\) is bigraded, as the intersection of kernels of bihomogeneous morphisms \([x,-]\), for \(x\) bihomogeneous.
      Finally, \(Z\), and thus \(Z_\mathbb {C}\), is semi-simple, and, for \((p,q)\neq (0,0)\), \(Z^{pq}\) is zero and thus nilpotent.\par{}A complex local subsystem \(W\) of \(H_\mathbb {C}\) is defined by a projector \(e\) of \(A_\mathbb {C}\).
      For \(t\in  S\), \(tW^\mathrm {c}\) is defined by the projector \(t\cdot  e\), and is thus locally constant.
      Since the group \(S\) is connected, to show that \(tW\) is isomorphic to \(W\), it suffices to show that, for \(t\) in a small enough neighbourhood of \(1\), the \(\mathbb {C}[\pi _1(S,s)]\) module \(tW_s\) is isomorphic to \(W_s\).
      Let \(H_i\) be the isotypic components of the \(\mathbb {C}[\pi _1(S,s)]\)-module \((H_\mathbb {C})_s\).
      We have
      \[         \begin {aligned}           W_s           &= \bigoplus _i(W_s\cap  H_i)         \\tW_s           &= \bigoplus _i(tW_s\cap  H_i).         \end {aligned}       \]\par{}Furthermore, for \(t\) close enough to \(1\), \(tW_s\) is close to \(W_s\) in the Grassmannian, and thus

      \begin{eqn}{4.2.8.1}\setnumber{4.2.8.1}\label{hodge-theory-ii-4.2.8.1}\skipsinglepar\[           \dim (tW_s\cap  H_i)           \leqslant  \dim (W_s\cap  H_i).         \tag{4.2.8.1}         \]\end{eqn}

      In fact, since \(\dim (tW_s)=\dim (W_s)\), we have equality in \hyperref[hodge-theory-ii-4.2.8.1]{Equation~\ref*{hodge-theory-ii-4.2.8.1}}.
      The various isotypic components of \(W_s\) and \(tW_s\) thus have, respectively, the same length, and \(W_s\) is isomorphic to \(tW_s\).\par{}We now prove (iii.b).
      It suffices to consider the case where \(H\) is homogeneous and \(W_s\) is a simple \(\mathbb {C}[\pi _1(S,s)]\)-module.
      Let \(H_1\) be the isotypic component of \((H_\mathbb {C})_s\) that contains \(W_s\).
      By (a), \(H_1\) is stable under \(S\).
      If \(W_s\) is of dimension \(d\), and \(H_1\) of length \(k\), then
      \[         \bigwedge ^{kd}H_1         \approx  \left (\bigwedge ^d W_s\right )^{\otimes  k}       \]
      as \(\mathbb {C}[\pi _1(S,s)]\)-modules.
      If \(\chi \) is the character of \(\pi _1(S,s)\) defined by \(\wedge ^d W\), then the character \(\chi _1\) defined by \(\wedge ^{kd}H_1\) is thus \(\chi ^k\).
      Let \(H_2=H_1+\overline {H}_1\).
      Since \(H_1\) is stable under \(S\), \(H_2\) is defined by a real Hodge substructure of \((H_\mathbb {R})_s\).
      Every polarisation form on \(H\) thus induces on \(H_2\) a non-degenerate bilinear form that is \(\pi _1(S,s)\)-invariant.
      
      If \(\dim (H_2)=e\), then the representation \((\wedge ^e H_2)^{\otimes 2}\) of \(\pi _1(S,s)\) is thus trivial.
      From this:

      
 \begin{enumerate}
        
 \item[(a)]
          For real \(H_1\), \(\chi ^{2k}\) is trivial.
        

        
 \item[(b)]
          For non-real \(H_1\) (so \(H_1\neq \overline {H}_1\)), \(\chi ^{2k}\cdot \overline {\chi }^{2k}\) is trivial.
        

      \end{enumerate}\par{}In both cases, \(|\chi |=1\).\par{}The representation of \(\pi _1(S,s)\) on \((H_\mathbb {C})_s\) comes from extension of scalars of a representation defined over \(\mathbb {Q}\) (namely \((H_\mathbb {Q})_s\)).
      The conjugate representations of \(W\) thus appear in \((H_\mathbb {C})_s\);
      there are at most \(N=\dim (H)\) of them, and for every automorphism \(\sigma \) of \(\mathbb {C}\) we have
      \[         |\sigma (\chi )|         = 1.       \]\par{}For \(\gamma \in \pi _1(S,s)\), \(\chi (\gamma )\) is an algebraic integer with at most \(N\) complex conjugates, and these are all of absolute value \(1\);
      thus \(\chi (\gamma )\) is a \(k\)-th root of unity, with \(k\leqslant  N\).
      Thus \(\chi ^{N!}=1\), which proves (b).\end{proof}\end{corollary}\begin{corollary}{4.2.9}\setnumber{4.2.9}\label{hodge-theory-ii-4.2.9}\skipsinglepar\par{}Let \(S\) be a smooth, connected, separated scheme, endowed with a base point \(s\in  S\).
    Let \(f\colon  X\to  S\) be a morphism of schemes such that \(\mathrm {R}^nf_*\mathbb {Q}\) is a local system on \(S\).
    Let \(G\) be the Zariski closure of the image of \(\pi _1(S,s)\) in \(\operatorname {Aut}_\mathbb {C}((\mathrm {R}^nf_*\mathbb {C})_s)\), and \(G^0\) the identity component of \(G\).
    Then:

    
 \begin{enumerate}
      
 \item[(a)]
        If \(f\) is smooth and proper, then \(G^0\) is semi-simple.
      

      
 \item[(b)]
        In general, \(G^0\) does not admit any quotient of multiplicative type (i.e. the radical of \(G^0\) is unipotent).
      

    \end{enumerate}\begin{proof}\par{}We first prove (a).
      By replacing \(S\) by a finite étale cover, if necessary, we can reduce to the case where \(G=G^0\).
      By (iii) of \hyperref[hodge-theory-ii-4.2.5]{Proposition~\ref*{hodge-theory-ii-4.2.5}}, the local system \(\mathrm {R}^nf_*\mathbb {Q}\) underlies an algebraic family \(H\) of Hodge structures, and so we can apply \hyperref[hodge-theory-ii-4.2.8]{Corollary~\ref*{hodge-theory-ii-4.2.8}} (by (i) of \hyperref[hodge-theory-ii-4.2.5]{Proposition~\ref*{hodge-theory-ii-4.2.5}}).
      The group \(G\) is reductive, since it admits a faithful semi-simple representation (namely \((\mathrm {R}^nf_*\mathbb {C})_s\)).
      If \(G\) (which we assume to be connected) were not semi-simple, then it would admit a non-trivial complex representation \(\rho \) of rank \(1\);
      since \(H_s\) is faithful, \(\rho \) would then be a direct factor of \(((H+H^*)_\mathbb {C}^{\otimes  m})_s\) for some \(m\).
      But this would contradict (iii.b) of \hyperref[hodge-theory-ii-4.2.8]{Corollary~\ref*{hodge-theory-ii-4.2.8}}, since \((H+H^*)^{\otimes  m}\) is algebraic, and no tensor power of \(\rho \) is trivial.\end{proof}\end{corollary}\par{}Let \(\pi \) be a group, and consider the following property of a representation \(\rho \) of \(\pi \) in a complex vector space \(V\) of finite dimension:

  \begin{property}{*}\setnumber{*}\label{hodge-theory-ii-4.2.star}\skipsinglepar\par{}\emph{\((*)\) The identity component \(G^0\) of the Zariski closure \(G\) of \(\rho (\pi )\) in \(\operatorname {Aut}(V)\) does not admit any quotient of multiplicative type.}\end{property}\begin{lemma}{4.2.10}\setnumber{4.2.10}\label{hodge-theory-ii-4.2.10}\skipsinglepar\par{}Let \(0\to  V'\to  V\to  V''\to 0\) be an exact sequence of representations of the above type.
    Then \(V\) satisfies \hyperref[hodge-theory-ii-4.2.star]{Property~\ref*{hodge-theory-ii-4.2.star}} if (and only if) \(V'\) and \(V''\) satisfy \hyperref[hodge-theory-ii-4.2.star]{Property~\ref*{hodge-theory-ii-4.2.star}}.\begin{proof}\par{}Let \(G\), \(G'\), and \(G''\) be the groups defined by \(V\), \(V'\), and \(V''\), respectively.
      Then \(V'\) is stable under \(G\), just as it is under \(\pi \), whence a morphism
      \[         u         = (u',u'')\colon  G\to  G'\times  G''       \]
      with unipotent kernel, and such that \(u'\) and \(u''\) are surjective.
      Let \(N\) be the identity component of the kernel of \(u''\).
      Since \(u'\) is surjective, \(u'(N)\) is distinguished in \(G'\), and thus does not admit any quotient of multiplicative type.
      
      If \(\chi \) is a character of \(G^0\), then \(\chi \) factors through \(u(G^0)\) and vanishes on \(N\);
      then a power of \(\chi \) factors through \({G''}^0\), and so \(\chi \) is trivial.\end{proof}\end{lemma}\begin{proof}\par{}We now prove part (b) of \hyperref[hodge-theory-ii-4.2.9]{Corollary~\ref*{hodge-theory-ii-4.2.9}} by dévissage, and by induction on the relative dimension of \(f\).
    We can suppose \(X\) to be reduced.\par{}Suppose first of all that \(X\) is quasi-projective.
    Then there exists a non-empty open \(U\) of \(S\), a dense open \(X'\) of \(X_U=f^{-1}(U)\) that is smooth over \(U\), and a compactification \(\overline {X}'\) of \(X'\) that is smooth and proper over \(U\).
    We can choose \(\overline {X}'\) to be such that some open \(X''\) of \(\overline {X}'\) contains \(X'\) and is proper over \(X_U\).

    
  \begin {tikzcd}[sep=huge]
      X_U
        \ar [ddr,swap,"f_U"]
      && X''
        \ar [r,hook,"k"]
        \ar [ll,swap,"p"]
        \ar [ddl,"f''"]
      & \overline {X}'
        \ar [ddll,"\overline {f}'"]
    \\&X'
        \ar [ul,hook',swap,"i"]
        \ar [ur,hook,"j"]
        \ar [d,"f'"]
    \\&U
    \end {tikzcd}


    We have long exact sequences of sheaves
    \[       \begin {gathered}         \ldots          \to  \mathrm {R}^i{f_U}_*(i_!\mathbb {Q})         \to  \mathrm {R}^i{f_U}_*\mathbb {Q}         \to  \mathrm {R}^i{f_U}_*(\mathbb {Q}{/}i_!\mathbb {Q})         \to  \ldots        \\\ldots          \to  \mathrm {R}^i{f''_U}_*(j_!\mathbb {Q})         \to  \mathrm {R}^i{f''_U}_*\mathbb {Q}         \to  \mathrm {R}^i{f''_U}_*(\mathbb {Q}{/}j_!\mathbb {Q})         \to  \ldots        \\\ldots          \to  \mathrm {R}^i f''_!\mathbb {Q}         \to  \mathrm {R}^i\overline {f}'_*\mathbb {Q}         \to  \mathrm {R}^i\overline {f}'_*(\mathbb {Q}{/}k_!\mathbb {Q})         \to  \ldots        \end {gathered}     \]\par{}For small enough \(U\), these sheaves are local systems, and the induction hypothesis applies to \(\mathrm {R}^\bullet {f_U}_*(\mathbb {Q}{/}i_!\mathbb {Q})\), to \(\mathrm {R}^\bullet {f''}_*(\mathbb {Q}{/}j_!\mathbb {Q})\), and to \(\mathrm {R}^\bullet \overline {f}'_*(\mathbb {Q}{/}k_!\mathbb {Q})\).
    By (a) of \hyperref[hodge-theory-ii-4.2.9]{Corollary~\ref*{hodge-theory-ii-4.2.9}}, \hyperref[hodge-theory-ii-4.2.star]{Property~\ref*{hodge-theory-ii-4.2.star}} is satisfied by \(\mathrm {R}^\bullet \overline {f}'_*\mathbb {Q}\).
    By \hyperref[hodge-theory-ii-4.2.10]{Lemma~\ref*{hodge-theory-ii-4.2.10}}, \hyperref[hodge-theory-ii-4.2.star]{Property~\ref*{hodge-theory-ii-4.2.star}} is then also satisfied by \(\mathrm {R}^\bullet {f''}_!\mathbb {Q}\), and by the dual local system \(\mathrm {R}^\bullet \overline {f}'_*\mathbb {Q}\) (Poincaré duality).
    By \hyperref[hodge-theory-ii-4.2.10]{Lemma~\ref*{hodge-theory-ii-4.2.10}}, \hyperref[hodge-theory-ii-4.2.star]{Property~\ref*{hodge-theory-ii-4.2.star}} is satisfied by \(\mathrm {R}^\bullet {f''}_*(j_!\mathbb {Q})\), which is isomorphic to \(\mathrm {R}^\bullet {f_U}_*(i_!\mathbb {Q})\), since \(p\) is proper, as well as by \(\mathrm {R}^\bullet {f_U}_*\mathbb {Q}\).
    Since the morphism \(\pi _1(U)\to \pi _1(S)\) is surjective, \(\mathrm {R}^if_*\mathbb {Q}\) satisfies (b) of \hyperref[hodge-theory-ii-4.2.9]{Corollary~\ref*{hodge-theory-ii-4.2.9}}.\par{}To pass from the case where \(f\) is quasi-projective to the general case, it suffices to use the Leray spectral sequence for a cover \((U_i)\) of \(X\) by quasi-projective opens:
    \[       \mathrm {R}^\bullet  f_*\left (\bigcap _i U_i,\mathbb {Q}\right )       \Rightarrow  \mathrm {R}^\bullet  f_*\mathbb {Q}     \]
    along with \hyperref[hodge-theory-ii-4.2.10]{Lemma~\ref*{hodge-theory-ii-4.2.10}}\end{proof}\section{Supplement to [D1968]}\label{hodge-theory-ii-4.3}\par{}In the footnote to (5.4) in [\hyperref[D1968]{D1968}], I claim without proof that:\begin{proposition}{4.3.1}\setnumber{4.3.1}\label{hodge-theory-ii-4.3.1}\skipsinglepar\par{}Let \(X\) be a smooth proper scheme over \(\mathbb {C}\).
    If a \(C^\infty \) form \(\alpha \) on \(X\) satisfies \(\mathrm {d}'\alpha =\mathrm {d}''\alpha =0\) and is cohomologous to zero, then there exists \(\beta \) such that \(\alpha =\mathrm {d}'\mathrm {d}''\beta \).\begin{proof}\par{}Here is the proof.
      As in \emph{loc. cit.}, we can reduce to the case where \(\alpha \) is bihomogeneous, with bidegree \((p,q)\).
      Since the spectral sequence of the double complex \(\operatorname {H}^0(X,\Omega ^{pq})\) degenerates, the exterior differential \(d\) is strictly compatible with the filtration by the first degree \(p\) (\hyperref[hodge-theory-ii-1.3.2]{Proposition~\ref*{hodge-theory-ii-1.3.2}}).
      If \(a\sim 0\), then there thus exists a form \(\beta _1\) such that \(\mathrm {d}\beta _1=\alpha \) and such that the components \(\beta _1^{p'q'}\) of \(\beta _1\) are zero for \(p'<p\).\par{}By complex conjugation, there similarly exists \(\beta _2\) such that \(\mathrm {d}\beta _2=\alpha \) and such that the components \(\beta _2^{p'q'}\) of \(\beta _2\) are zero for \(q'<q\).
      We have \(\mathrm {d}(\beta _1-\beta _2)=0\).\par{}Let \(\gamma \) be a form satisfying \(\mathrm {d}'\gamma =\mathrm {d}''\gamma =0\) in the cohomology class of \(\beta _1-\beta _2\) (\emph{loc. cit.}).
      Let \(\gamma _1\) (resp. \(\gamma _2\)) be the sum of the components of \(\gamma \) of degree \((p',q')\) for \(p'\geqslant  p\) (resp. \(q'\geqslant  q\)).
      If we set \(\beta '_1=\beta _1-\gamma _1\) and \(\beta '_2=\beta _2+\gamma _2\), then we still have that \(\mathrm {d}\beta '_1=\mathrm {d}\beta '_2=\alpha \);
      furthermore, \(\beta '_1-\beta '_2=\beta _1-\beta _2-\gamma \) is cohomologous to zero: there exists \(\delta \) such that
      \[         \mathrm {d}\delta          = \beta '_1 - \beta '_2.       \]
      Let \(\delta _1\) (resp. \(\delta _2\)) be the sum of the components of \(\delta \) of degree \((p',q')\) for \(p'<p-1\) (resp. \(q'>q-1\));
      let \(\beta \) be the component of \(\delta \) of degree \((p-1,q-1)\).\par{}We have that
      \[         \beta '_2         =\mathrm {d}\delta _2 + \mathrm {d}''\beta        \]
      and
      \[         \begin {aligned}           \alpha            = \mathrm {d}\beta '_2           &= \mathrm {d}\mathrm {d}\delta _2 + \mathrm {d}\mathrm {d}''\beta          \\&= \mathrm {d}'\mathrm {d}''\beta .         \end {aligned}       \]\end{proof}\end{proposition}\section{Homomorphisms of abelian schemes}\label{hodge-theory-ii-4.4}\par{}We intend to show how, in some cases, we can remove from \cref{hodge-theory-ii-4.1.3.2} the assumption of algebraicity at a point.\label{hodge-theory-ii-4.4.1}\par{}Let \(f\colon  X\to  S\) be a smooth proper morphism;
    we denote by \(\mathrm {R}_i f_*\mathbb {Z}\) the local system on \(S^\mathrm {an}\) whose homology is the fibres of \(f\).
    At each point \(s\in  S\), \((\mathrm {R}_i f_*\mathbb {Z})_s\simeq \operatorname {H}_i(X_s,\mathbb {Z})\) is endowed with a Hodge structure of weight \(-i\), dual to the Hodge structure on \(\operatorname {H}^i(X_s,\mathbb {Z})\).
    This structure varies continuously with \(s\in  S\).\label{hodge-theory-ii-4.4.2}\par{}If \(f\colon  X\to  S\) is an abelian scheme, then the exponential defines an exact sequence
    \[       0       \to  \mathrm {R}_1f_*\mathbb {Z}       \to  \operatorname {Lie}(X)       \to  X       \to  0     \]
    of sheaves on \(S^\mathrm {an}\), and the Hodge filtration is the short exact sequence
    \[       0       \to  \operatorname {Ker}(\alpha )       \to  \mathrm {R}_1f_*\mathbb {Z}\otimes \mathcal {O}       \xrightarrow {\alpha } \operatorname {Lie}(X)       \to  0.     \]\par{}It is known that:\begin{reminder}{4.4.3}\setnumber{4.4.3}\label{hodge-theory-ii-4.4.3}\skipsinglepar\par{}Let \(S\) be a smooth scheme of finite type over \(\mathbb {C}\).
    The construction of \cref{hodge-theory-ii-4.4.2} establishes an equivalence of categories between

    
 \begin{enumerate}
      
 \item[(a)]
        the category of abelian schemes over \(S\); and
      

      
 \item[(b)]
        the category of polarisable continuous families of Hodge structures \(H\) of type \((-1,0)+(0,-1)\) on \(S\) such that \(H_\mathbb {Z}\) is torsion free and the Hodge filtration varies holomorphically on \(S\).
      

    \end{enumerate}\begin{proof}\par{}The essential surjectivity of the functor that goes from (a) to (b) follows from Borel (\emph{unpublished});
      the full faithfulness follows from the fact that, for two abelian schemes \(X_1\) and \(X_2\) over \(S\), the \(S\)-scheme \(\underline {\operatorname {Hom}}_S(X_1,X_2)\) is unramified over \(S\), and thus has the same algebraic and holomorphic sections.\end{proof}\end{reminder}\label{hodge-theory-ii-4.4.4}\par{}Let \(Z\) be a commutative field, \(H\) a simple central algebra over \(Z\), and \(*\) a \(Z\)-linear anti-automorphism of \(H\).
    After extension of scalars, we have \(H\simeq \operatorname {End}(V)\) for some vector space \(V\) over \(Z\), and \(*\) is then the transposition with respect to a bilinear form \(\Phi \) on \(V\), unique up to a factor.
    Suppose that \(*\) is involutive;
    the forms \(\Phi (X,Y)\) and \(\Phi (Y,X)\) are then proportional, i.e. \(\Phi (X,Y)=\lambda \Phi (Y,X)\) for some \(\lambda ^2=1\).
    \emph{We set \(\varepsilon _*=\lambda \)}.
    If \([H:Z]=d^2\), then
    \[       \operatorname {Tr}(*\colon  H\to  H)       = \varepsilon _*\cdot  d.     \]\par{}Every \(Z\)-linear automorphism \(C\) of \(H\) is interior: \(Cx=cxc^{-1}\).
    If \(C\) is involutive and commutes with \(*\), then the elements \(c^2\) and \(cc^*\) are central, so that \(c^*=\lambda  c\) with \(\lambda \) central.
    Since \(c^{**}=c\), we have that \(\lambda ^2=1\).
    \emph{We set \(\varepsilon _C=\lambda \)}.\par{}We easily see that the involution \(C\circ  *\) satisfies
    \[       \varepsilon _{C\circ  *}       = \varepsilon _C\cdot \varepsilon _*.     \]\begin{reminder}{4.4.5}\setnumber{4.4.5}\label{hodge-theory-ii-4.4.5}\skipsinglepar\par{}Let \(X\) be a simple abelian variety over a field \(K\), and let \(H\) be the field \(\operatorname {End}_K(X)\otimes \mathbb {Q}\).
    Every polarisation of \(X\) defines a positive involution \(*\) of \(H\), i.e.
    \[       \operatorname {Tr}(hh^*) > 0       \quad \text {if }h\neq 0.     \]
    If \(Z_0\) is the \(*\)-invariant subfield of the centre \(Z\) of \(H\), then \(Z_0\) is completely real and we are in one of the following cases:

    
 \begin{enumerate}
      
 \item[(a)]
        \(Z=Z_0\) and, for every real place \(\rho \) of \(Z\), \(H\otimes _\rho \mathbb {R}\) is an algebra of matrices over \(\mathbb {R}\), and \(\varepsilon _*=1\).
      


      
 \item[(b)]
        \(Z=Z_0\) and, for every real place \(\rho \) of \(Z\), \(H\otimes _\rho \mathbb {R}\) is an algebra of matrices over \(\mathbb {H}\), and \(\varepsilon _*=-1\).
      


      
 \item[(c)]
        \(Z\) is a completely imaginary quadratic extension of \(Z_0\).
      

    \end{enumerate}\end{reminder}\label{hodge-theory-ii-4.4.6}\par{}Let \(f_i\colon  X_i\to  S\) (for \(i=1,2\)) be an abelian scheme over a smooth scheme \(S\).
    The free abelian group
    \[       \operatorname {Hom}(\mathrm {R}_1{f_1}_*\mathbb {Z},\mathrm {R}_1{f_2}_*\mathbb {Z})       \simeq  \Gamma (S,\underline {\operatorname {Hom}}(\mathrm {R}_1{f_1}_*\mathbb {Z},\mathrm {R}_1{f_2}_*\mathbb {Z}))     \]
    
    is then endowed with a Hodge structure of weight \(0\) (\hyperref[hodge-theory-ii-4.2.8]{Corollary~\ref*{hodge-theory-ii-4.2.8}}).
    Further, by \hyperref[hodge-theory-ii-2.1.11.1]{Remark~\ref*{hodge-theory-ii-2.1.11.1}} and \hyperref[hodge-theory-ii-4.4.3]{Reminder~\ref*{hodge-theory-ii-4.4.3}}, the group \(\operatorname {Hom}_S(X_1,X_2)\) is the degree \((0,0)\) component of this group.\label{hodge-theory-ii-4.4.7}\par{}Let \(S\) be a non-empty, connected, smooth scheme, with generic point \(\eta \), and let \(f\colon  X\to  S\) be an abelian scheme over \(S\).
    The centre \(Z\) of \(H=\operatorname {End}(\mathrm {R}_1f_*\mathbb {Z})\) is then of degree \((0,0)\) (\hyperref[hodge-theory-ii-4.2.8]{Corollary~\ref*{hodge-theory-ii-4.2.8}}), and thus contained in \(\operatorname {End}_S(X)\simeq \operatorname {End}_{k(\eta )}(X_\eta )\).\par{}A polarisation of \(X\) defines an involution \(*\) of \(H\).
    The restriction of this to \(\operatorname {End}_S(X)\) is positive, so that \(Z\) is the product of completely real fields or quadratic imaginary extensions of completely real fields.\par{}Denote by \(C\) the actions of \(i\in  S(\mathbb {R})\) on \(H_\mathbb {R}=H\otimes \mathbb {R}\) and on the \((\mathrm {R}_1f_*\mathbb {R})_s\) for \(s\in  S\).
    On \(H_\mathbb {R}\), \(C^2=1\), and \(C\) commutes with \(*\);
    on \((\mathrm {R}_1f_*\mathbb {R})_s\), \(C^2=-1\).
    If \(\psi \) is the alternating form on \(\mathrm {R}_1f_*\mathbb {Z}\) with integer values induced by the polarisation of \(X\), then
    \[       \psi (hx,Cy)       = \psi (x,h^*Cy)       = \psi (x,C\cdot  C(h^*)\cdot  y).     \]
    We thus deduce that

    \begin{eqn}{4.4.7.1}\setnumber{4.4.7.1}\label{hodge-theory-ii-4.4.7.1}\skipsinglepar\[         \operatorname {Tr}(h\cdot  C(h^*)) > 0         \quad \text {for }h\neq 0.       \tag{4.4.7.1}       \]\end{eqn}\label{hodge-theory-ii-4.4.8}\par{}Now suppose that \(X_\eta \) is simple, so that \(Z\) is a field.
    Let \(\rho \colon  Z\to \mathbb {C}\) be a complex embedding.
    Then the algebra
    \[       H_\rho        = H\otimes _{Z,\rho }\mathbb {C}     \]
    is an algebra of matrices and a bigraded direct factor of \(H_\mathbb {C}=H\otimes _\mathbb {Q}\mathbb {C}\).
    Pick an isomorphism

    \begin{eqn}{4.4.8.1}\setnumber{4.4.8.1}\label{hodge-theory-ii-4.4.8.1}\skipsinglepar\[         H_\rho          \simeq  \operatorname {End}(V_\rho ).       \tag{4.4.8.1}       \]\end{eqn}

    From the fact that the infinitesimal generators of \(S\) act on \(H_\rho \), by derivations, which are necessarily interior, we deduce that \(V_\rho \) admits a bigrading, unique up to translation, such that \hyperref[hodge-theory-ii-4.4.8.1]{Equation~\ref*{hodge-theory-ii-4.4.8.1}} is a bigraded isomorphism.
    Let \((\mathrm {R}_1f_*\mathbb {C})_\rho \) be the bigraded local system
    \[       (\mathrm {R}_1f_*\mathbb {C})\otimes _{Z\otimes \mathbb {C},\rho }\mathbb {C}     \]
    which is a direct factor of \(\mathrm {R}_1f_*\mathbb {C}\).
    Denote by \(T_\rho \) the local system
    \[       T_\rho        = \operatorname {Hom}_{H_\rho }(V_\rho ,(\mathrm {R}_1f_*\mathbb {C})_\rho ).     \]\par{}We have the bigraded isomorphism

    \begin{eqn}{4.4.8.2}\setnumber{4.4.8.2}\label{hodge-theory-ii-4.4.8.2}\skipsinglepar\[         V_\rho \otimes _\mathbb {C} T_\rho          \xrightarrow {\sim } (\mathrm {R}_1f_*\mathbb {C})_\rho .       \tag{4.4.8.2}       \]\end{eqn}

    Since \((\mathrm {R}_1f_*\mathbb {C})_\rho \) is of degree \((-1,0)+(0,-1)\), one of the following conditions is satisfied:

    
 \begin{enumerate}
      
 \item[(a)]
        \(V_\rho \) is bihomogeneous;
      

      
 \item[(b)]
        \(T_\rho \) is bihomogeneous.
      

    \end{enumerate}\par{}It is clear that condition (a) (resp. (b)) is simultaneously satisfied by \(\rho \) and \(\overline {\rho }\).
    If condition (a) is satisfied at each place, then \(H\) is of degree \((0,0)\).
    If condition (b) is satisfied at each place, then the Hodge structure on \(\mathrm {R}_1f_*\mathbb {Q}\) is locally constant.\label{hodge-theory-ii-4.4.9}\par{}Now suppose that \(Z\) is completely real.
    For each real embedding \(\varphi \) of \(Z\), we set \(\varepsilon _{H,\varphi }=+1\) if \(H\otimes _{Z,\varphi }\mathbb {R}\) is an algebra of matrices; we set \(\varepsilon _{H,\varphi }=-1\) otherwise.
    If a complex embedding \(\rho \) induces a real embedding \(\varphi \), then we set \(\varepsilon _{H\rho }=\varepsilon _{H\varphi }\).\par{}Let \(\rho \colon  Z\to \mathbb {C}\) factor through \(\varphi \colon  Z\to \mathbb {R}\).
    Complex conjugation on \(H_\rho \) is induced by an antilinear automorphism \(\sigma _1\) of \(V_\rho \), unique up to real scalar, with scalar square.
    Since \(\sigma _1\) commutes with \(\sigma _1^2\), \(\sigma _1^2\) is real and, by normalising \(\sigma _1\), we can assume that \(\sigma _1^2=\pm 1\).
    We then have that

    \begin{eqn}{4.4.9.1}\setnumber{4.4.9.1}\label{hodge-theory-ii-4.4.9.1}\skipsinglepar\[         \sigma _1^2         = \varepsilon _{H\rho }.       \tag{4.4.9.1}       \]\end{eqn}\par{}The complex conjugation automorphism on
    \[       (\mathrm {R}_1f_*\mathbb {C})_\rho        =((\mathrm {R}_1f_*\mathbb {Q})\otimes _{Z,\varphi }\mathbb {R})\otimes _\mathbb {R}\mathbb {C}     \]
    is "compatible" with \hyperref[hodge-theory-ii-4.4.8.2]{Equation~\ref*{hodge-theory-ii-4.4.8.2}} and can be written as
    \[       \sigma        = \sigma _1\otimes \sigma _2.     \]
    Since \(\sigma ^2=1\), we have, by \hyperref[hodge-theory-ii-4.4.9.1]{Equation~\ref*{hodge-theory-ii-4.4.9.1}}, that

    \begin{eqn}{4.4.9.2}\setnumber{4.4.9.2}\label{hodge-theory-ii-4.4.9.2}\skipsinglepar\[         \sigma _1^2         = \sigma _2^2         = \varepsilon _{H\rho }.       \tag{4.4.9.2}       \]\end{eqn}\par{}Let \(\Phi \) be a polarisation form on \(\mathrm {R}_1f_*\mathbb {Z}\).
    This form is necessarily of the form
    \[       \Phi        = \operatorname {Tr}_{Z/\mathbb {Q}}(\psi )     \]
    where \(\psi \) is an alternating \(Z\)-bilinear form on \(\mathrm {R}_1f_*\mathbb {Q}\).\par{}Let \(\psi _\rho \) be the form induced on \(V_\rho \otimes _\mathbb {C} T_\rho \).
    The alternating form \(\psi _\rho \) is invariant under \(\pi _1(S,s)\).
    Since \(T_\rho \) is irreducible, we can write \(\psi _\rho \) as

    \begin{eqn}{4.4.9.3}\setnumber{4.4.9.3}\label{hodge-theory-ii-4.4.9.3}\skipsinglepar\[         \psi _\rho          = A_\rho \otimes  B_\rho        \tag{4.4.9.3}       \]\end{eqn}

    with \(A_\rho \) on \(V_\rho \), and \(B_\rho \) on \(T_\rho \) invariant under \(\pi _1(S,s)\).\par{}Since \(T_\rho \) is irreducible, \(B_\rho \) is either symmetric or alternating.
    Since \(\psi _\rho \) is alternating, \(A_\rho \) is either symmetric or alternating.
    Since the \(T_\rho \) are conjugate to one another, whether we are in one case or the other does not depend on \(\rho \).
    We set

    \begin{eqn}{4.4.9.4}\setnumber{4.4.9.4}\label{hodge-theory-ii-4.4.9.4}\skipsinglepar\[         A_\rho (X,Y)         = \varepsilon _V A_\rho (Y,X)       \tag{4.4.9.4}       \]\end{eqn}

    \begin{eqn}{4.4.9.5}\setnumber{4.4.9.5}\label{hodge-theory-ii-4.4.9.5}\skipsinglepar\[         B_\rho (X,Y)         = \varepsilon _T B_\rho (Y,X).       \tag{4.4.9.5}       \]\end{eqn}

    We thus have

    \begin{eqn}{4.4.9.6}\setnumber{4.4.9.6}\label{hodge-theory-ii-4.4.9.6}\skipsinglepar\[         \varepsilon _V\varepsilon _T         = -1       \tag{4.4.9.6}       \]\end{eqn}

    and it is clear by \cref{hodge-theory-ii-4.4.4} that if \(*\) is the involution of \(H\) induced by \(\psi \), then

    \begin{eqn}{4.4.9.7}\setnumber{4.4.9.7}\label{hodge-theory-ii-4.4.9.7}\skipsinglepar\[         \varepsilon _*         = \varepsilon _V.       \tag{4.4.9.7}       \]\end{eqn}\par{}In case (a) (resp. (b)) of \cref{hodge-theory-ii-4.4.8}, we normalise the bigradings so that \(V_\rho \) (resp. \(T_\rho \)) is of bidegree \((0,0)\).
    We then have, for non-zero \(x\in  V_\rho \) and non-zero \(y\in  T_\rho \), the two cases:

    
 \begin{enumerate}
      
 \item[(a)]
        \(A_\rho (x,\sigma _1 x)B_\rho (y,C\sigma _2 y)           >0\);
      

      
 \item[(b)]
        \(A_\rho (x,C\sigma _1 x)B_\rho (y,\sigma _2 y)           >0\).
      

    \end{enumerate}\par{}By normalising \(A_\rho \) and \(B_\rho \) (of which only their tensor product is given), we can suppose, in each respective case, that

    
 \begin{enumerate}
      
 \item[(a)]
        \(A_\rho (x,\sigma _1x)>0\) and \(B_\rho (y,C\sigma _2y)>0\);
      

      
 \item[(b)]
        \(A_\rho (x,C\sigma _1x)>0\) and \(B_\rho (y,\sigma _2y)\).
      

    \end{enumerate}\par{}Furthermore, \(A_\rho \) and \(B_\rho \), just like \(\Phi \), are invariant under the compact sub-torus \(\operatorname {Ker}(N)\) of \(S\), and in particular under \(C\).
    In case (a), we have that
    \[       0       < A_\rho (\sigma _1x,\sigma _1\sigma _1x)       = \varepsilon _V\cdot \varepsilon _{H,\rho }A_\rho (x,\sigma _1x)     \]
    whence \(\varepsilon _V\cdot \varepsilon _{H\rho }=1\).
    Similarly, in case (b), we see that \(\varepsilon _T\cdot \varepsilon _{H\rho }=1\).\begin{lemma}{4.4.10}\setnumber{4.4.10}\label{hodge-theory-ii-4.4.10}\skipsinglepar\par{}For \(\varepsilon _T=\varepsilon _{H\rho }\), we are in case (b);
    for \(\varepsilon _T=-\varepsilon _{H\rho }\), we are in case (a).\end{lemma}\begin{proposition}{4.4.11}\setnumber{4.4.11}\label{hodge-theory-ii-4.4.11}\skipsinglepar\par{}Let \(S\) be a non-empty, connected, smooth scheme, with generic point \(\eta \), and let \(f\colon  X\to  S\) be an abelian scheme over \(S\).
    Let \(H=\operatorname {End}(\mathbb {R}_1f_*\mathbb {Q})\), and let \(Z\) be the centre of \(H\).
    Suppose that

    
 \begin{enumerate}
      
 \item[(a)]
        \(X_\eta \) is simple over \(k(\eta )\).
      


      
 \item[(b)]
        \(X\) does not become constant on any finite covering of \(S\).
      


      
 \item[(c)]
        One of the following conditions is satisfied (cf. \hyperref[hodge-theory-ii-4.4.5]{Reminder~\ref*{hodge-theory-ii-4.4.5}}):

        
 \begin{enumerate}
          
 \item[(1)]
            \(Z\) is imaginary \emph{quadratic};
          

          
 \item[(2)]
            \(Z\) is completely real, and for each real embedding \(\varphi \colon  Z\to \mathbb {R}\), \(H\otimes _{Z,\varphi }\mathbb {R}\) is an algebra of matrices over \(\mathbb {R}\);
          

          
 \item[(3)]
            \(Z\) is completely real, and for each real embedding \(\varphi \colon  Z\to \mathbb {R}\), \(H\otimes _{Z,\varphi }\mathbb {R}\) is an algebra of matrices over \(\mathbb {H}\).
          

        \end{enumerate}

      

    \end{enumerate}\par{}Then
    \[       \operatorname {End}(X)       \xrightarrow {\sim } \operatorname {End}(\mathrm {R}_1f_*\mathbb {Z}).     \]\begin{proof}\par{}We will show that we are either in case (a) or case (b) of \cref{hodge-theory-ii-4.4.8}, for every place of \(Z\) simultaneously.
      This is clear (\cref{hodge-theory-ii-4.4.8}) under the hypothesis (c.1), and follows from \hyperref[hodge-theory-ii-4.4.10]{Lemma~\ref*{hodge-theory-ii-4.4.10}} under the hypotheses (c.2) or (c.3), since not only \(\varepsilon _T\), but also \(\varepsilon _{H\rho }\), is then independent of \(\rho \).\par{}Condition (b) of \cref{hodge-theory-ii-4.4.8} cannot be satisfied at every place of \(Z\), since the Hodge structure would then be locally constant, which (by \cref{hodge-theory-ii-4.1.3.3} and \hyperref[hodge-theory-ii-4.4.3]{Reminder~\ref*{hodge-theory-ii-4.4.3}}) contradicts (b).
      Condition (a) is thus satisfied at each place of \(Z\), and so \(H\) is of degree \((0,0)\) (\cref{hodge-theory-ii-4.4.8}).
      Taking \hyperref[hodge-theory-ii-4.4.3]{Reminder~\ref*{hodge-theory-ii-4.4.3}} into account, this finishes the proof.\end{proof}\end{proposition}\begin{proposition}{4.4.12}\setnumber{4.4.12}\label{hodge-theory-ii-4.4.12}\skipsinglepar\par{}Let \(f\colon  X\to  S\) be an abelian scheme over a smooth scheme \(S\).
    Then the following conditions are equivalent:

    
 \begin{enumerate}
      
 \item[(i)]
        For every abelian scheme \(Y\to  S\), we have
        \[           \operatorname {Hom}_S(X,Y)           \xrightarrow {\sim } \operatorname {Hom}_S(\mathrm {R}_1f_*\mathbb {Z},\mathrm {R}_1g_*\mathbb {Z}).         \]
      

      
 \item[(ii)]
        Condition (i) is satisfied for \(X=Y\), and the centre \(Z\) of \(\operatorname {End}_S(X)\otimes \mathbb {Q}\) does not admit any complex place \(\rho \colon  Z\to \mathbb {C}\) such that the direct factor \(\mathrm {R}_1f_*\mathbb {Q}\otimes _{Z,\rho }\mathbb {C}\) of \(\mathrm {R}_1f_*\mathbb {C}\) is of pure Hodge degree \((-1,0)\).
      

    \end{enumerate}\begin{proof}\par{}We can reduce to the case where \(S\) is non-empty and connected, with generic point \(\eta \), and where \(X_\eta \) is a simple abelian variety over \(k(\eta )\).
      If (i) or (ii) is satisfied, then \(D=\operatorname {End}(X_\eta )\otimes \mathbb {Q}\) is a field, and, for \(s\in  S\), \((\mathrm {R}_1f_*\mathbb {Q})_s\) is an irreducible representation of \(\pi _1(S,s)\).\par{}To prove that (ii) \(\implies \) (i), we can suppose \(Y_\eta \) to be simple, and \((\mathrm {R}_1g_*\mathbb {Q})_s\) to be isotypic of type \((\mathrm {R}_1f_*\mathbb {Q})_s\).
      We then have an isomorphism of continuous families of Hodge \(\mathbb {Q}\)-structures (with \(D\) of degree \((0,0)\))
      \[         \mathrm {R}_1g_*\mathbb {Q}         \approx  \mathrm {R}_1f_*\mathbb {Q}\otimes _D\operatorname {Hom}(\mathrm {R}_1f_*\mathbb {Q},\mathrm {R}_1g_*\mathbb {Q}).       \]\par{}We have that \(D\otimes \mathbb {C}\simeq \mathbb {M}_n(Z\otimes \mathbb {C})\).
      If \(e\in  D\otimes \mathbb {C}\) is sent by such an isomorphism to the idempotent \(e_{11}\), then we again have an isomorphism of bigraded vector spaces
      \[         (\mathrm {R}_1g_*\mathbb {C})_s         \approx  (\mathrm {R}_1f_*\mathbb {C})_s\otimes _{Z\otimes \mathbb {C}}e(\operatorname {Hom}(\mathrm {R}_1f_*\mathbb {C},\mathrm {R}_1g_*\mathbb {C})).       \]\par{}If condition (ii) is satisfied, and if \(e(\operatorname {Hom}(\mathrm {R}_1f_*\mathbb {C},\mathrm {R}_1g_*\mathbb {C}))\) is not of pure degree \((0,0)\), then \((\mathrm {R}_1g_*\mathbb {C})_s\) cannot be of pure degree \((-1,0)+(0,-1)\), which is a contradiction.
      We thus deduce that \(\operatorname {Hom}(\mathrm {R}_1f_*\mathbb {Z},\mathrm {R}_1g_*\mathbb {Z})\) is of pure degree \((0,0)\), and (i) then follows by \hyperref[hodge-theory-ii-4.4.3]{Reminder~\ref*{hodge-theory-ii-4.4.3}} and \hyperref[hodge-theory-ii-2.1.11.1]{Remark~\ref*{hodge-theory-ii-2.1.11.1}}.\par{}Let \(Z\) be a completely imaginary quadratic extension of a completely real number field \(Z^0\).
      If \(Z\otimes \mathbb {R}\approx \mathbb {C}\oplus \ldots \oplus \mathbb {C}\), then we define an action by homotheties of \(S\) on \(Z\otimes \mathbb {R}\) by sending \(s\in  S(\mathbb {R})\simeq \mathbb {C}^*\) to \((s^2\cdot  N(s)^{-1},1,\ldots ,1)\).
      The pair consisting of \(Z\) and this action of \(S\) on \(Z\otimes \mathbb {R}\) is a polarisable Hodge structure \(Z_\rho \) of degree \((-1,1)+(0,0)+(1,-1)\) that depends on \(Z\) and the chose place \(\rho \colon  Z\to \mathbb {C}\).\par{}Let \(X\) be an abelian scheme over \(S\) such that \(X_\eta \) is simple.
      If the centre \(Z\) of \(\operatorname {End}(\mathrm {R}_1f_*\mathbb {Q})\) is not completely real, then it is a completely imaginary quadratic extension of a completely real number field.
      If the second hypothesis of (ii) is not satisfied, then the direct factor \(\mathrm {R}_1f_*\mathbb {Q}\otimes _Z Z_\rho \) of the continuous family of polarisable Hodge structures \(\mathrm {R}_1f_*\mathbb {Q}\otimes _\mathbb {Q} Z_\rho \) is of pure degree \((-1,0)+(0,-1)\).\par{}By \hyperref[hodge-theory-ii-4.4.3]{Reminder~\ref*{hodge-theory-ii-4.4.3}}, the choice of an integer lattice in \(\mathrm {R}_1f_*\mathbb {Q}\otimes _Z Z_\rho \) defines an abelian scheme \(g\colon  Y\to  S\), and \(\operatorname {Hom}(\mathrm {R}_1f_*\mathbb {Q},\mathrm {R}_1g_*\mathbb {Q})\) is not of pure degree \((0,0)\), which contradicts (i).\end{proof}\end{proposition}\begin{corollary}{4.4.13}\setnumber{4.4.13}\label{hodge-theory-ii-4.4.13}\skipsinglepar\par{}Let \(S\) be a smooth connected scheme with generic point \(\eta \), and let \(f\colon  X\to  S\) be an abelian scheme over \(S\) of relative dimension \(g\leqslant 3\).
    Suppose that, on any finite étale covering of \(S\), \(X\) does not admit a constant abelian subscheme.
    Then, for every abelian scheme \(g\colon  Y\to  S\), we have
    \[       \begin {aligned}         \operatorname {Hom}(X,Y)         &\xrightarrow {\sim } \operatorname {Hom}(\mathrm {R}_1f_*\mathbb {Z},\mathrm {R}_1g_*\mathbb {Z})       \\\operatorname {Hom}(Y,X)         &\xrightarrow {\sim } \operatorname {Hom}(\mathrm {R}_1g_*\mathbb {Z},\mathrm {R}_1f_*\mathbb {Z}).       \end {aligned}     \]\begin{proof}\par{}By duality, it suffices to prove the first claim.
      We can also suppose \(X_\eta \) to be simple.
      We will show that \(X\) satisfies one of the conditions of (c) in \hyperref[hodge-theory-ii-4.4.11]{Proposition~\ref*{hodge-theory-ii-4.4.11}}.\par{}Set \(H=\mathbb {M}_c(D)\), with \(D\) a skew field of rank \(b^2\) over its centre \(Z\), which we know to be either completely real or completely imaginary.
      If we set \(a=[Z:\mathbb {Q}]\), then we have

      \begin{eqn}{4.4.13.1}\setnumber{4.4.13.1}\label{hodge-theory-ii-4.4.13.1}\skipsinglepar\[           ab^2c \mathbin {\vert } 2g \leqslant  6.         \tag{4.4.13.1}         \]\end{eqn}\begin{enumerate}
      
 \item[(i)]
        \par{}If \(Z\) is completely imaginary, and not imaginary quadratic, then we must have \(a=2g=4\).
          For \(s\in  S\), the centraliser of \(Z\) acting on \((\mathrm {R}_1f_*\mathbb {Q})_s\) is thus commutative, and so the action of \(\pi _1(S,s)\) is abelian, and thus factors through a finite group ((iii.b) of \hyperref[hodge-theory-ii-4.2.8]{Corollary~\ref*{hodge-theory-ii-4.2.8}} or \hyperref[hodge-theory-ii-4.2.9]{Corollary~\ref*{hodge-theory-ii-4.2.9}}).
          This is a contradiction (\cref{hodge-theory-ii-4.1.3.3} and \hyperref[hodge-theory-ii-4.4.3]{Reminder~\ref*{hodge-theory-ii-4.4.3}}).

        \par{}With the notation of \cref{hodge-theory-ii-4.4.8}, we can also note that the \(T_\rho \) are then of rank \(1\), and so condition (b) of \cref{hodge-theory-ii-4.4.8} is satisfied at each place, and the Hodge structure is locally constant.
      


      
 \item[(ii)]
        If \(Z\) is completely real, and if conditions (c.1) and (c.2) are not satisfied, then \(a\geqslant 2\) and \(b\geqslant 2\), which is a contradiction.
      

    \end{enumerate}\par{}We will show that \(X\) satisfies the conditions of \hyperref[hodge-theory-ii-4.4.12]{Proposition~\ref*{hodge-theory-ii-4.4.12}}.
      If \(Z\) is imaginary quadratic, and if there exists \(\rho \colon  Z\to \mathbb {C}\) such that \(\mathrm {R}_1f_*\mathbb {Q}\otimes _{Z,\rho }\mathbb {C}\) is of pure Hodge type \((-1,0)\), then the Hodge structure on \(\mathrm {R}_1f_*\mathbb {Q}\) is locally constant, which is a contradiction.
      This proves \hyperref[hodge-theory-ii-4.4.13]{Corollary~\ref*{hodge-theory-ii-4.4.13}}.\end{proof}\end{corollary}\label{hodge-theory-ii-4.4.14}\par{}Let \(f\colon  X\to  S\) be a polarised abelian scheme of dimension \(g\) over a smooth connected base \(S\).
    For \(s\in  S\), the fundamental group \(\pi _1(S,s)\) acts on the fibre \(\operatorname {H}^1(X_s,\mathbb {Z})\) of \(\mathrm {R}_1f_*\mathbb {Z}\) over \(s\).
    The polarisation of \(X\) defines an alternating form on \(\operatorname {H}^1(X_s,\mathbb {Z})\), with values in \(\mathbb {Z}(-1)\), that is non-degenerate over \(\mathbb {Q}\), and the action of \(\pi _1(S,s)\) respects this form.
    Consider the hypothesis:

    \begin{hypothesis}{4.4.14.1}\setnumber{4.4.14.1}\label{hodge-theory-ii-4.4.14.1}\skipsinglepar\par{}The morphism defined by \(X\) from \(S\) to one of the irreducible components of the corresponding scheme of modules of polarised abelian varieties is a dominant morphism.\end{hypothesis}

    \par{}The following result was suggested to me by J.-P. Serre.\begin{corollary}{4.4.15}\setnumber{4.4.15}\label{hodge-theory-ii-4.4.15}\skipsinglepar\par{}Let \(S\) be a smooth connected scheme, and \(f\colon  X\to  S\) a polarised abelian scheme over \(S\) that satisfies \hyperref[hodge-theory-ii-4.4.14.1]{Hypothesis~\ref*{hodge-theory-ii-4.4.14.1}}.
    Then, for every abelian scheme \(g\colon  Y\to  S\), we have
    \[       \operatorname {Hom}(X,Y)       \xrightarrow {\sim } \operatorname {Hom}(\mathrm {R}_1f_*\mathbb {Z},\mathrm {R}_1g_*\mathbb {Z}).     \]\begin{proof}\par{}Let \(s\in  S\).
      By \hyperref[hodge-theory-ii-4.4.11]{Proposition~\ref*{hodge-theory-ii-4.4.11}} and \hyperref[hodge-theory-ii-4.4.12]{Proposition~\ref*{hodge-theory-ii-4.4.12}}, it suffices to show that the representation of \(\pi _1(S,s)\) on \((\mathrm {R}_1f_*\mathbb {Q})_s\) is absolutely irreducible, so that \(H=\mathbb {Q}\).
      This follows from the following lemma.\end{proof}\end{corollary}\begin{lemma}{4.4.16}\setnumber{4.4.16}\label{hodge-theory-ii-4.4.16}\skipsinglepar\par{}If \(X\) satisfies \hyperref[hodge-theory-ii-4.4.14.1]{Hypothesis~\ref*{hodge-theory-ii-4.4.14.1}}, then the image of \(\pi _1(S,s)\) in the group of symplectic automorphisms of \(\operatorname {H}^1(X_s,\mathbb {Z})\) is of finite index.\begin{proof}\par{}Let \(n\) be an integer.
      Replacing \(S\) by a surjective étale covering, defined by a finite-index subgroup of \(\pi _1(S,s)\), if necessary, we can assume that the local system \({}_nX=\operatorname {Ker}(n\cdot \mathrm {Id}\colon  X\to  X)\) is trivial on \(S\).
      Replacing \(X\) by an isogenous abelian scheme, we can then assume that \(X\) is polarised in the principal series.
      Suppose that \(n\geqslant  3\).\par{}The abelian scheme \(X\) then defines a dominant morphism \(x\) from the \(S\) to the scheme of "Jacobi" modules \(M_n\) of polarised abelian schemes in the principal series endowed with a symplectic isomorphism between \({}_nX\) and \((\mathbb {Z}{/}n)^{2g}\);
      furthermore, \(X\) is an inverse image under \(x\) of the universal abelian scheme on \(M_n\).
      
      We know that \(\pi _1(M_n,x(s))\) is sent isomorphically to the subgroup of \(\operatorname {Sp}(\operatorname {H}^1(X_s,\mathbb {Z}))\) given by the kernel of the reduction \(\operatorname {mod}{n}\) map.
      It thus remains to show that:\end{proof}\end{lemma}\begin{lemma}{4.4.17}\setnumber{4.4.17}\label{hodge-theory-ii-4.4.17}\skipsinglepar\par{}Let \(p\colon  S\to  M\) be a dominant morphism of smooth connected schemes.
    The subgroup \(p(\pi _1(S,s))\) of \(\pi _1(M,p(s))\) is then of finite index.\begin{proof}\par{}Let \(t\) be a closed point of the generic fibre of \(p\), and let \(T'\) be its closure in \(S\).
      There exists a dense Zariski open \(U\) of \(M\) such that \(T=p^{-1}(U)\cap  T'\) is an étale covering of \(U\).\par{}We lose no generality by assuming that \(s\in  T\).
      In the diagram
      
  \begin {tikzcd}
        \pi _1(T,s)
          \ar [r]
          \ar [d,swap,"a"]
        & \pi _1(S,s)
          \ar [d,"c"]
      \\\pi _1(U,p(s))
        \ar [r,swap,"b"]
        & \pi _1(M,p(s))
      \end {tikzcd}

      the image of the arrow \(a\) is of finite index, since \(T\) is a finite étale covering, and the arrow \(b\) is surjective, since \(M\setminus  U\) is of real codimension \(\geqslant 2\).
      The image of the arrow \(c\) is thus of finite index.\end{proof}\end{lemma}